\documentclass[a4paper,11pt]{article}

\usepackage[utf8]{inputenc}
\usepackage[T1]{fontenc}
\usepackage{newtxtext}
\usepackage[left=2.5cm, right=2.5cm, top=2.5cm, bottom=2.5cm]{geometry}
\usepackage[french]{babel}
\usepackage{setspace}
\usepackage{graphicx}
\usepackage{etoolbox}
\usepackage{csquotes}
\usepackage{caption}
\usepackage{fancyhdr}
\usepackage{titlesec}
\usepackage{tocloft}
\usepackage{amsmath,amssymb,amsfonts}
\usepackage{mathrsfs}
\usepackage{float}
\usepackage{dsfont}
\usepackage{tikz}
\usepackage{url}
\usepackage{color, xcolor}
\usepackage{array}
\usepackage{booktabs}
\usepackage{colortbl}
\usepackage{longtable}
\usepackage{pdflscape}
\usepackage{enumitem}
\usepackage{mdframed}
\usepackage{hyperref}
\usetikzlibrary{arrows.meta, positioning, calc, shapes.geometric, fit, backgrounds}

% ---------------------------------------------------------------- charte ----
% Charte reprise de la note technique Solvabilite II : bleu financier pour la
% structure, vert pour l'accent et les sorties, deux gris de tableau.
\definecolor{maincolor}{RGB}{0,70,140}
\definecolor{greenline}{RGB}{0,180,80}
\definecolor{tableheader}{RGB}{200,220,250}
\definecolor{tablerow}{RGB}{245,245,245}
\definecolor{financeblue}{RGB}{0,70,140}

\newcommand{\annee}{2026}
\newcommand{\datedoc}{\today}
\newcommand{\versiondoc}{2.0}
\newcommand{\statutdoc}{Projet}
% Titre court utilise dans le pied de page (fancyhdr) : remplace la
% mention « Note technique Solvabilite II » du document source.
\newcommand{\titredoc}{Documentation de l'outil de calibrage USP}

\captionsetup[figure]{labelfont={bf, color=maincolor}, textfont={normalfont}, labelsep=colon}
\captionsetup[table]{labelfont={bf, color=maincolor}, textfont={normalfont}, labelsep=colon}

\hypersetup{colorlinks=true, linkcolor=financeblue, urlcolor=financeblue,
            pdftitle={Documentation de l'outil de calibrage USP - Solvabilite II}}

\renewcommand{\baselinestretch}{1.08}
\pagestyle{fancy}
\fancyhf{}
\fancyhead[L]{\color{maincolor}\nouppercase{\leftmark}}
\renewcommand{\headrulewidth}{0pt}
\renewcommand{\footrulewidth}{0.4pt}
\fancyfoot[R]{\color{greenline}\textbf{page \thepage}}
\fancyfoot[L]{\color{maincolor}\titredoc}
\fancypagestyle{plain}{\fancyhf{}%
  \renewcommand{\headrulewidth}{0pt}\renewcommand{\footrulewidth}{0.4pt}%
  \fancyfoot[R]{\color{greenline}\textbf{page \thepage}}%
  \fancyfoot[L]{\color{maincolor}\titredoc}}

% ------------------------------------------------------ titres de section ---
% Le bandeau bleu pleine largeur avec le titre en vert marque le niveau
% section, qui joue ici le role du chapitre de la note technique.
\newcommand{\titrebandeau}[1]{%
  \noindent{\color{financeblue}\rule{\linewidth}{2.5pt}}
  \par\vspace{5mm}
  {\color{greenline}\raggedright\fontsize{19}{23}\bfseries\selectfont #1\par}
  \vspace{6mm}}

\titleformat{\section}[display]
  {\normalfont}{}{0pt}
  {\titrebandeau}
\titlespacing*{\section}{0pt}{0pt}{0pt}
\titleformat{\subsection}
  {\color{financeblue}\normalfont\large\bfseries}{\thesubsection}{0.7em}{}
\titleformat{\subsubsection}
  {\color{financeblue}\normalfont\normalsize\bfseries}{\thesubsubsection}{0.6em}{}

\renewcommand{\cftsecfont}{\bfseries\color{financeblue}}
\renewcommand{\cftsecaftersnum}{\color{financeblue}}
\renewcommand{\cftsecpagefont}{\color{financeblue}}
\renewcommand{\cftsubsecfont}{\color{black}}
\renewcommand{\cftsubsecpagefont}{\color{black}}
\setlength{\cftsecnumwidth}{2.2em}
\setcounter{tocdepth}{2}

\newcommand{\code}[1]{\texttt{\small #1}}
\newcommand{\refl}[1]{{\small\textcolor{maincolor!70!black}{\textit{Référence.} #1}}\par\smallskip}

% Encadre de citation reglementaire : filet vertical bleu en marge.
\newcommand{\reglement}[2]{%
  \par\vspace{2mm}\noindent\begin{minipage}{\linewidth}
  \small\setlength{\parindent}{0pt}
  \textcolor{financeblue}{\vrule width 2pt}\hspace{4mm}%
  \begin{minipage}{\dimexpr\linewidth-7mm\relax}
  \textcolor{financeblue}{\textbf{#1}}
  \par\vspace{2pt}
  {\itshape #2}
  \end{minipage}
  \end{minipage}\par\vspace{3mm}}

% L'encadre generique reprend le filet bleu de la charte.
\newmdenv[linecolor=financeblue, linewidth=0pt, leftline=true,
          innerleftmargin=6mm, innertopmargin=4pt, innerbottommargin=4pt,
          backgroundcolor=tablerow, skipabove=3mm, skipbelow=3mm]{encadre}

% Fiches de test : six rubriques fixes.
\newcounter{ficheid}
\newenvironment{fiche}[1]{%
  \stepcounter{ficheid}%
  \phantomsection
  \subsection{#1}%
  \label{fiche:\theficheid}%
}{%
  \end{description}\vspace{0.5em}%
}
\newcommand{\Cadre}{%
  \begin{description}[leftmargin=0pt, style=unboxed, itemsep=3pt, topsep=3pt, parsep=2pt]%
  \item[\textcolor{financeblue}{\textbf{1. Cadre, notations et hypothèses de validité (non testées).}}]}
\newcommand{\Hzero}{\item[\textcolor{financeblue}{\textbf{2. Hypothèse nulle.}}]}
\newcommand{\Stat}{\item[\textcolor{financeblue}{\textbf{3. Statistique de test.}}]}
\newcommand{\Loi}{\item[\textcolor{financeblue}{\textbf{4. Loi sous $H_0$.}}]}
\newcommand{\Pval}{\item[\textcolor{financeblue}{\textbf{5. p-value.}}]}
\newcommand{\Usage}{\item[\textcolor{financeblue}{\textbf{6. Usage et limites.}}]}


% Titre de partie : allege par rapport a une premiere version qui doublait le
% bandeau chapitre (filet bleu, texte "Partie", grand titre vert, PUIS un
% second filet vert) -- le document source (main.tex) ne comporte aucun
% niveau au-dessus du chapitre et ne double jamais le filet. On reprend donc
% ici le meme principe qu'un titre de chapitre (un seul filet bleu, titre en
% vert dessous), avec un corps de texte plus petit pour rester subordonne au
% bandeau \titrebandeau des sections qui suivent.
\newcommand{\partie}[2]{%
  \clearpage
  \vspace*{1cm}
  \noindent{\color{financeblue}\rule{\linewidth}{2pt}}\par\vspace{3mm}
  {\color{financeblue}\raggedright\normalsize\bfseries\scshape #1\par}
  \vspace{1.5mm}
  {\color{greenline}\raggedright\fontsize{20}{24}\bfseries\selectfont #2\par}
  \vspace{5mm}
  \addcontentsline{toc}{section}{\protect\numberline{}\textsc{#1} --- #2}}

\emergencystretch=1.5em

\begin{document}

% ============================================================ page de garde ==
\thispagestyle{empty}
\begin{tikzpicture}[remember picture, overlay]
  \fill[financeblue]
    ([xshift=2.5cm,yshift=-3.2cm]current page.north west) rectangle
    ([xshift=-2.5cm,yshift=-3.28cm]current page.north east);
  \fill[greenline]
    ([xshift=2.5cm,yshift=-3.2cm]current page.north west) rectangle
    ([xshift=7.5cm,yshift=-3.28cm]current page.north west);
  \fill[financeblue]
    ([xshift=2.5cm,yshift=3.6cm]current page.south west) rectangle
    ([xshift=-2.5cm,yshift=3.53cm]current page.south east);
\end{tikzpicture}%
\vspace*{2.1cm}
\noindent
{\color{financeblue}\normalsize\bfseries\scshape Exercice \annee}

\vspace{1.7cm}
\noindent
{\color{financeblue}\fontsize{28}{34}\bfseries\selectfont
 Documentation\\ de l'outil\par}

\vspace{1.1cm}
\noindent
{\color{black}\Large
 Calibrage des paramètres propres à l'entreprise\\
 Moteur de calcul, tests statistiques et application\par}

\vspace{0.7cm}
\noindent
{\color{black}\normalsize\itshape
 Règlement délégué (UE) 2015/35, annexe XVII --- profondeur d'intérêt $T = 8$\par}

\vfill
\noindent
{\footnotesize\color{black}
 Établie le \datedoc\quad\textbar\quad Version \versiondoc\quad\textbar\quad
 Statut : \statutdoc\quad\textbar\quad Moteur \code{R/engine.R}\par}
\vspace{2.6cm}
\clearpage

% ================================================================= resume ==
\titrebandeau{Objet du document}
\markboth{Objet du document}{}

\noindent
Ce document constitue la documentation de l'\textbf{outil de calibrage des
paramètres propres à l'entreprise} (USP) au titre des méthodes « risque de
prime » et « risque de réserve~1 » de l'annexe~XVII du règlement délégué
(UE)~2015/35. Il couvre le moteur de calcul \code{R/engine.R}, les tests
statistiques qu'il met en \oe uvre, et l'application interactive qui l'exploite.

\medskip\noindent
Il est organisé en quatre parties. La \textbf{partie~I} décrit l'architecture de
l'outil, les notations et les contrôles de recevabilité des données. La
\textbf{partie~II} pose le cadre statistique commun~: nature des p-values
(exactes, asymptotiques, de Monte-Carlo) et séparation de l'erreur numérique
liée au nombre $B$ de réplications d'avec l'erreur d'approximation liée à la
taille $T$ de l'échantillon. La \textbf{partie~III} rassemble les fiches de
tests, groupées par hypothèse réglementaire H1 à H4, chacune suivant une
structure uniforme en six rubriques~: cadre et hypothèses de validité (non
testées), hypothèse nulle, statistique de test, loi sous l'hypothèse nulle,
p-value, usage et limites. La \textbf{partie~IV} rassemble la validation
empirique de l'implémentation, les références bibliographiques et la synthèse
de lecture.

\medskip\noindent
Une attention particulière est portée à la profondeur d'intérêt du projet,
$T = 8$, à laquelle plusieurs des tests présentés sont structurellement
incapables de rejeter au seuil usuel de 5~\%. Le sommaire détaille les
sous-sections afin de permettre l'accès direct à chaque test.

\vspace{8mm}
\reglement{Article 218 du règlement délégué (UE) 2015/35}{Les entreprises
d'assurance et de réassurance peuvent remplacer un sous-ensemble des paramètres
standard du module « souscription non-vie », du module « souscription santé » et
du module « souscription vie » par des paramètres propres à l'entreprise.}

\vspace{6mm}
\noindent
\begin{center}
\small
\rowcolors{2}{tablerow}{white}
\begin{tabular}{@{}p{5.2cm}p{9.6cm}@{}}
\rowcolor{tableheader}
\textcolor{financeblue}{\textbf{Élément}} & \textcolor{financeblue}{\textbf{Référence}} \\
Base réglementaire & Directive 2009/138/CE, art.~104(7) ; règlement délégué
(UE) 2015/35, art.~218 à 220 et annexe XVII \\
Procédure d'agrément & Règlement d'exécution (UE) 2015/498 \\
Transposition & Code des assurances, art.~R.~352-5 (V) et R.~356-19 (II) \\
Moteur de calcul & \code{R/engine.R} (R base et \code{stats} uniquement) \\
Interface & \code{app.R}, \code{R/display\_helpers.R} (aucun calcul) \\
Profondeur d'intérêt & $T = 8$ années de survenance \\
\end{tabular}
\end{center}

\clearpage
\titrebandeau{Sommaire}
\markboth{Sommaire}{}
% \tableofcontents composerait son propre titre (\contentsname, soit
% « Table des matières »), qui ferait doublon avec le bandeau ci-dessus.
% On n'appelle donc que \@starttoc, qui compose les entrées seules.
\makeatletter
\@starttoc{toc}
\makeatother
\newpage


\partie{Partie I}{Présentation de l'outil}

% =============================================================================
\section{Architecture du moteur de calcul}
\label{sec:archi}
% =============================================================================

\subsection{Principe de séparation}

Le projet applique une séparation stricte entre le \emph{moteur quantitatif} et
la \emph{couche d'affichage}. La règle appliquée est la suivante~: si un
résultat peut être calculé indépendamment de l'interface graphique, son calcul
figure dans \code{R/engine.R}.

\begin{center}
\small
\rowcolors{2}{tablerow}{white}
\begin{tabular}{@{}p{4.3cm}p{5.1cm}p{5.4cm}@{}}
\rowcolor{tableheader}
\textcolor{financeblue}{\textbf{Fichier}} &
\textcolor{financeblue}{\textbf{Contenu}} &
\textcolor{financeblue}{\textbf{Interdit}} \\
\code{R/engine.R} & Préparation des données, contrôles, estimation, tests,
p-values, bootstrap, jackknife, calibration, quantités des graphiques &
Toute primitive Shiny (\code{input\$}, \code{output\$}, \code{reactive()},
\code{render*()}) \\
\code{app.R} & Saisie, déclenchement, stockage réactif, agencement des onglets &
Tout calcul statistique ou actuariel \\
\code{R/display\_helpers.R} & Formatage des nombres et des tableaux, tracés à
partir de \code{res\$plots\_data} & Tout calcul de statistique, de p-value ou
de paramètre \\
\end{tabular}
\end{center}

\medskip
Cette séparation a été vérifiée automatiquement~: aucune primitive Shiny n'est
présente dans le moteur, et aucun appel de fonction de calcul
(\code{optim}, \code{lm}, \code{*.test}, \code{rnorm}, \code{quantile},
\code{sd}, \code{var}, \code{mean}) ne subsiste dans la couche d'affichage.
Le moteur s'exécute seul~:

\begin{encadre}
\small
\code{source("R/engine.R")}\\
\code{res <- run\_engine(xt = ..., yt = ..., methode = "premium", segment = 1,}\\
\code{\phantom{res <- run\_engine(}T = 8, B = 999, seed = 20260831)}\\
\code{res\$parametre\_final\$sigma\_usp \ \ ;\ \ engine\_table\_tests(res)}
\end{encadre}

\subsection{Chaîne de calcul et flux de données}

La figure~\ref{fig:archi} donne l'enchaînement des fonctions du moteur. Les blocs
sont les fonctions, les flèches portent les objets transmis. Les blocs bleus sont
les entrées, les blancs les étapes de calcul, les verts les sorties exposées à
l'application. Le point d'aiguillage est le paramètre \code{methode}~: il oriente
vers l'une des deux chaînes, dont les géométries de données et les jeux de tests
diffèrent. L'annexe~\ref{sec:annexe-flux} en donne le détail fonction par
fonction.

\begin{figure}[H]
\centering
\begin{tikzpicture}[
  font=\small,
  ent/.style   = {rectangle, rounded corners=2pt, draw=financeblue, line width=0.7pt,
                  fill=tableheader, text=financeblue, align=center, inner sep=3pt,
                  minimum height=6mm, text width=46mm, font=\small\bfseries},
  fonc/.style  = {rectangle, rounded corners=2pt, draw=financeblue, line width=0.7pt,
                  fill=white, text=black, align=center, inner sep=3pt,
                  minimum height=6mm, text width=46mm},
  sortie/.style= {rectangle, rounded corners=2pt, draw=greenline, line width=0.9pt,
                  fill=greenline!8, text=black, align=center, inner sep=3pt,
                  minimum height=6mm, text width=52mm},
  test/.style  = {diamond, aspect=2.6, draw=financeblue, line width=0.7pt,
                  fill=tableheader, text=financeblue, align=center,
                  inner sep=0pt, text width=24mm, font=\scriptsize\bfseries},
  fl/.style    = {-{Latex[length=2mm]}, draw=financeblue, line width=0.6pt},
  li/.style    = {draw=financeblue, line width=0.6pt},
  flx/.style   = {font=\scriptsize\itshape, text=greenline, inner sep=1.5pt}
]

\node[ent] (in) {Entrées utilisateur\\[1pt]
  \scriptsize \code{xt}, \code{yt} \emph{ou} \code{triangle},\\
  \scriptsize \code{methode}, \code{segment}, \code{annexe},\\
  \scriptsize \code{T}, \code{B}, \code{alpha}, \code{seed}};

\node[fonc, below=5mm of in] (orch) {\code{run\_engine()}\\[1pt]
  \scriptsize orchestrateur unique};
\node[test, below=5mm of orch] (aig) {\code{methode}\\\code{= reserve2} ?};

% ---------------- branche lognormale (gauche) ------------------------------
\node[fonc, below left=7mm and 14mm of aig] (vln) {\code{engine\_valider\_donnees()}};
\node[fonc, below=4mm of vln] (fit) {\code{usp\_ajuster()}\\[1pt]
  \scriptsize L-BFGS-B sur \code{usp\_objectif()},\\
  \scriptsize via \code{usp\_noyau()} et \code{usp\_pi()}};
\node[fonc, below=4mm of fit] (bln) {\code{usp\_bootstrap()}\\[1pt]
  \scriptsize \code{.stats\_bootstrapables()}};
\node[fonc, below=4mm of bln] (tln) {\code{usp\_tests()}\\[1pt]
  \scriptsize H1 à H4, stabilité, robustesse};
\node[fonc, below=4mm of tln] (rln) {\code{usp\_jackknife()}, \code{usp\_profil()}};
\node[fonc, below=4mm of rln] (cln) {\code{usp\_parametre()}\\[1pt]
  \scriptsize $\times$ \code{usp\_credibilite()}};
\node[fonc, below=4mm of cln] (pln) {\code{engine\_plots\_data()}};

% ---------------- branche Merz--Wüthrich (droite) --------------------------
\node[fonc, below right=7mm and 14mm of aig] (vmw) {\code{mw\_valider\_triangle()}};
\node[fonc, below=4mm of vmw] (amw) {\code{mw\_ajuster()}\\[1pt]
  \scriptsize $\hat f_j$, $\hat\sigma^2_j$, $\hat C(i,j)$};
\node[fonc, below=4mm of amw] (mmw) {\code{mw\_msep()}\\[1pt]
  \scriptsize erreur de prédiction à un an};
\node[fonc, below=4mm of mmw] (bmw) {\code{mw\_bootstrap()}\\[1pt]
  \scriptsize \code{mw\_simuler\_triangle()}, \code{.mw\_stats()}};
\node[fonc, below=4mm of bmw] (tmw) {\code{mw\_tests()}\\[1pt]
  \scriptsize M1 à M6};
\node[fonc, below=4mm of tmw] (cmw) {\code{mw\_parametre()}\\[1pt]
  \scriptsize $\times$ \code{usp\_credibilite()}};
\node[fonc, below=4mm of cmw] (pmw) {\code{mw\_plots\_data()}};

% ---------------- sorties ---------------------------------------------------
\coordinate (jonc) at ($(pln)!0.5!(pmw) + (0,-9mm)$);
\node[sortie, below=13mm of aig |- pmw, yshift=-9mm] (out)
  {objet \code{usp\_engine}\\[1pt]
   \scriptsize \code{\$donnees \ \$tests \ \$bootstrap \ \$calibration}\\
   \scriptsize \code{\$parametre\_final \ \$plots\_data \ \$metadata}};
\node[sortie, below=5mm of out] (app)
  {Couche d'affichage\\[1pt]
   \scriptsize \code{app.R}, \code{display\_helpers.R} : aucun calcul};

% ---------------- liaisons --------------------------------------------------
\draw[fl] (in) -- (orch);
\draw[fl] (orch) -- (aig);
\draw[fl] (aig) -| node[flx, above left] {non} (vln);
\draw[fl] (aig) -| node[flx, above right] {oui} (vmw);
\foreach \a/\b in {vln/fit, fit/bln, bln/tln, tln/rln, rln/cln, cln/pln,
                   vmw/amw, amw/mmw, mmw/bmw, bmw/tmw, tmw/cmw, cmw/pmw}
  \draw[fl] (\a) -- (\b);
\draw[li] (pln) |- (jonc);
\draw[li] (pmw) |- (jonc);
\draw[fl] (jonc) -- (out);
\draw[fl] (out) -- node[flx, right] {\code{res}} (app);
\end{tikzpicture}
\caption{Architecture du moteur \code{R/engine.R}. Le paramètre \code{methode}
aiguille vers l'une des deux chaînes~: à gauche les méthodes lognormales
(sections~B et~C de l'annexe~XVII), à droite la méthode Merz--Wüthrich
(section~D). Les deux convergent vers un objet de sortie unique, que la couche
d'affichage se contente de lire. \code{usp\_credibilite()} est la seule fonction
de calcul partagée par les deux branches.}
\label{fig:archi}
\end{figure}

\subsection{Inventaire des fonctions}

\begin{center}
\small
\rowcolors{2}{tablerow}{white}
\begin{tabular}{@{}p{4.6cm}p{4.0cm}p{6.2cm}@{}}
\rowcolor{tableheader}
\textcolor{financeblue}{\textbf{Fonction}} &
\textcolor{financeblue}{\textbf{Entrées}} &
\textcolor{financeblue}{\textbf{Sorties}} \\
\multicolumn{3}{@{}l@{}}{\textit{\textcolor{financeblue}{Contrôles et données}}} \\
\code{engine\_valider\_donnees()} & \code{xt}, \code{yt}, \code{T\_min} &
\code{\$ok}, \code{\$erreurs}, \code{\$avertissements} \\
\code{usp\_controle\_donnees()} & \code{xt}, \code{yt}, \code{alpha} &
liste de contrôles de recevabilité \\
\multicolumn{3}{@{}l@{}}{\textit{\textcolor{financeblue}{Noyau du modèle}}} \\
\code{usp\_pi()} & $\delta$, $\gamma$, $x$ & vecteur $\pi_t$ \\
\code{usp\_noyau()} & $\delta$, $\gamma$, $x$, $y$ &
$\ln\hat\beta$, $v_t$, $z_t$, $\sigma$, objectif \\
\code{usp\_objectif()} & \code{par}, $x$, $y$ & $-2\ln L$ (à une constante près) \\
\code{usp\_ajuster()} & $x$, $y$ & \code{fit} complet, FOC, convergence \\
\multicolumn{3}{@{}l@{}}{\textit{\textcolor{financeblue}{Lois exactes}}} \\
\code{mk\_p\_exacte()} & série $r_t$ & p-value exacte (loi mahonienne) \\
\code{runs\_p\_exacte()} & résidus $z_t$ & p-value exacte (Swed \& Eisenhart) \\
\multicolumn{3}{@{}l@{}}{\textit{\textcolor{financeblue}{Statistiques et tests}}} \\
\code{stat\_*()}, \code{test\_*()} & $x$, $y$ ou $z$ & statistique et p-value \\
\code{test\_tost\_intercept()} & $x$, $y$, $\theta$ ou $\Delta$ &
p-value d'équivalence (TOST) \\
\code{.stats\_bootstrapables()} & $x$, $y$, $z$ & vecteur de 34 statistiques \\
\code{usp\_bootstrap()} & \code{fit}, $B$, \code{seed} &
\code{p\_mc}, \code{err\_mc}, \code{sigma\_boot} \\
\code{usp\_tests()} & \code{fit}, \code{boot}, \code{alpha} &
liste des tests (H0, H1, statistique, 3 p-values, nature) \\
\multicolumn{3}{@{}l@{}}{\textit{\textcolor{financeblue}{Robustesse et calibration}}} \\
\code{usp\_jackknife()} & \code{fit}, $\sigma_{\text{std}}$ &
table des $\sigma_{\mathrm{USP}}^{(-t)}$ \\
\code{usp\_profil()} & \code{fit} & profils en $\delta$ et $\gamma$, tests LR \\
\code{usp\_segment\_infos()} & segment, annexe & libellé, $\sigma$ standard, barème \\
\code{usp\_credibilite()} & $T$, barème & facteur $c$ \\
\code{usp\_parametre()} & \code{fit}, $\sigma_{\text{std}}$, barème &
$\sigma_{\mathrm{USP}}$ et ses composantes \\
\multicolumn{3}{@{}l@{}}{\textit{\textcolor{financeblue}{Sorties}}} \\
\code{engine\_plots\_data()} & \code{fit}, \code{boot}, \code{profil} &
toutes les quantités tracées \\
\code{engine\_table\_tests()} & \code{res} & \code{data.frame} auditable \\
\code{run\_engine()} & entrées utilisateur & objet \code{usp\_engine} complet \\
\end{tabular}
\end{center}

\subsection{Périmètres couverts et paramètres standard}
\label{sec:segments}

L'outil couvre les deux segmentations pour lesquelles l'annexe~XVII autorise des
paramètres propres au titre du risque de primes et de réserve~: l'annexe~II
(engagements non-vie) et l'annexe~XIV (engagements de santé non similaires à la
vie, dits « non-SLT »). Les écarts-types standard ci-dessous sont ceux du
règlement délégué (UE)~2015/35, \emph{Journal officiel de l'Union européenne},
L~12 du 17 janvier 2015, p.~230--231 et p.~269. Ils sont codés dans les tables
\code{ANNEXE\_II} et \code{ANNEXE\_XIV} de \code{R/engine.R}.

\begin{center}
\small
\rowcolors{2}{tablerow}{white}
\begin{tabular}{@{}cp{8.0cm}ccc@{}}
\rowcolor{tableheader}
\textcolor{financeblue}{\textbf{Seg.}} &
\textcolor{financeblue}{\textbf{Annexe XIV --- santé non-SLT}} &
\textcolor{financeblue}{\textbf{LoB}} &
\textcolor{financeblue}{\textbf{$\sigma$ prime}} &
\textcolor{financeblue}{\textbf{$\sigma$ réserve}} \\
1 & Assurance frais médicaux et réassurance proportionnelle y afférente & 1 et 13 & 5 \% & 5 \% \\
2 & Assurance protection du revenu et réassurance proportionnelle y afférente & 2 et 14 & 8,5 \% & 14 \% \\
3 & Assurance indemnisation des travailleurs et réassurance proportionnelle y afférente & 3 et 15 & 8 \% & 11 \% \\
4 & Réassurance santé non proportionnelle & 25 & 17 \% & 20 \% \\
\end{tabular}
\end{center}

\noindent
Le facteur de crédibilité applicable est déterminé par la section~G de
l'annexe~XVII, dont le paragraphe~2 range explicitement \emph{tous} les segments
de l'annexe~XIV sous le barème court, aux côtés des segments 2 à 4 et 7 à 12 de
l'annexe~II. Le barème long du paragraphe~1 est réservé aux seuls segments 1, 5
et 6 de l'annexe~II. L'annexe d'appartenance fait donc partie de la clé~: à
$T = 8$, le segment~1 de l'annexe~XIV (frais médicaux) reçoit $c = 81~\%$, alors
que le segment~1 de l'annexe~II (responsabilité civile automobile) reçoit
$c = 59~\%$. La fonction \code{usp\_bareme\_segment(segment, annexe)} implémente
cette distinction, et \code{usp\_segment\_infos()} renvoie l'ensemble des
caractéristiques réglementaires d'un segment.

\begin{encadre}
\textbf{Portée.} L'outil ne couvre pas les deux autres paramètres remplaçables
de l'article~218~: le facteur d'ajustement pour la réassurance non
proportionnelle (segments 1, 4 et 5 de l'annexe~II) et l'écart-type du risque de
révision. Il ne met pas davantage en \oe uvre la méthode « risque de réserve
n\textsuperscript{o}~2 » (chaîne d'échelle et erreur de prédiction à un an), qui
relève d'un dispositif distinct.
\end{encadre}

\subsection{Reproductibilité et gestion des graines}

Une seule graine pilote l'ensemble des tirages~: elle est passée à
\code{run\_engine(seed = )} puis à \code{usp\_bootstrap()}, qui appelle
\code{set.seed()} une fois avant la boucle de réplications. Aucune autre
fonction du moteur ne tire de nombres aléatoires, à l'exception de
\code{engine\_plots\_data()}, qui fixe sa propre graine pour l'enveloppe de
simulation du QQ-plot afin que le graphique soit stable d'un appel à l'autre.

Le script \code{tests/test\_reproductibilite.R} vérifie, pour chacune des trois
méthodes, deux propriétés qu'il ne faut pas confondre.

\emph{La reproductibilité.} Deux appels de \code{run\_engine()} avec les mêmes
données, les mêmes paramètres et la même graine, \textbf{sur une même machine},
produisent des objets identiques au bit près. C'est ce que contrôle la
comparaison de deux appels successifs par \code{identical()}, et cette propriété
tient (vérification exécutée).

\emph{La non-régression.} Le résultat est en outre confronté à des valeurs de
référence enregistrées (\code{tests/reference/}), produites sur une plateforme
unique désignée et versionnées avec le code. Cette seconde comparaison n'est
\textbf{pas} au bit près~: elle emploie \code{all.equal()} à une tolérance
relative de $10^{-8}$, qui porte sur une différence relative \emph{moyenne} et
absorbe donc une dérive diffuse, si nombreuses qu'en soient les valeurs
touchées. Une telle dérive existe entre plateformes et a été mesurée~: sur le cas
de contrôle de la méthode prime, 5\,515 des 12\,409 valeurs élémentaires de
l'objet diffèrent au sens strict, pour une différence relative moyenne de
$6{,}2\cdot 10^{-10}$ et un écart relatif maximal de $3{,}5\cdot 10^{-7}$ sur les
tirages bootstrap~; l'écart sur $\sigma_{\mathrm{USP}}$ lui-même est de l'ordre
de $10^{-14}$ en relatif, très en deçà de la précision à laquelle le paramètre
est lu. La branche Merz--Wüthrich est, elle, indemne (aucun écart sur le triangle
de contrôle). L'origine vraisemblable de cette dérive est l'optimiseur employé
par la branche lognormale sur une vraisemblance quasi plate en $\delta$, dont le
chemin dépend de l'ordre des opérations en virgule flottante~; cette explication
est cohérente avec la répartition observée des écarts, mais elle n'a pas été
démontrée par une expérience isolant la bibliothèque d'algèbre linéaire ou la
version de R. Ce qui précède ne remet en cause aucun résultat publié~: c'est la
portée de la garantie qui doit être énoncée exactement.

\begin{encadre}
\textbf{Point de revue.} L'objet retourné contient un champ \code{\$metadata}
consignant la méthode, le segment, $\sigma_{\text{standard}}$, le barème de
crédibilité, $T$, $B$, la granularité $1/(B+1)$, le seuil $\alpha$, la graine,
l'horodatage, la durée et la version de R. Ce champ suffit à rejouer un calcul
à l'identique.
\end{encadre}

\subsection{Index des fonctions de statistique et de p-value}
\label{sec:index-fonctions}

L'inventaire de la section~1.3 regroupe les fonctions par catégorie
(\code{stat\_*()}, \code{test\_*()}). Le tableau ci-dessous descend d'un niveau~:
il relie \textbf{chaque test documenté} --- méthode lognormale en partie~III,
méthode Merz--Wüthrich en sections~\ref{sec:mwm1} à~\ref{sec:mwm6} --- à la ou aux fonctions R
qui calculent sa statistique, et à celle(s) qui calculent chacune des p-values
disponibles. Les entrées sont groupées par hypothèse~: H1 à H4 pour la méthode
lognormale, M1 à M6 pour la méthode Merz--Wüthrich. La colonne « Réf. » renvoie à la fiche correspondante. La méthode
générale de p-value de Monte-Carlo est décrite à part
(fiche~\ref{fiche:1}, page~\pageref{fiche:1})~: la colonne « Clé MC » indique,
pour chaque test, le nom de la statistique dans \code{.stats\_bootstrapables()}
qui lui correspond, lorsqu'une p-value de Monte-Carlo est calculée.

Certains tests s'appuient sur un appel direct à une fonction de \code{stats}
plutôt que sur une fonction dédiée du moteur~: cela est indiqué explicitement,
la fonction \code{stats::}\emph{...} restant néanmoins un élément identifiable et
auditable du code de \code{usp\_tests()}.

\begin{landscape}
\begin{center}
\scriptsize
\setlength{\tabcolsep}{4.5pt}
\renewcommand{\arraystretch}{1.22}
\begin{longtable}{@{}p{4.4cm}p{0.9cm}p{5.7cm}p{7.8cm}p{2.1cm}@{}}
\toprule
\textbf{Test (fiche)} & \textbf{Réf.} & \textbf{Fonction(s) --- statistique} &
\textbf{Fonction(s) / méthode --- p-value non simulée} & \textbf{Clé MC} \\
\midrule
\endfirsthead
\toprule
\textbf{Test (fiche)} & \textbf{Réf.} & \textbf{Fonction(s) --- statistique} &
\textbf{Fonction(s) / méthode --- p-value non simulée} & \textbf{Clé MC} \\
\midrule
\endhead
\multicolumn{5}{@{}l@{}}{\textit{Suite page suivante}} \\
\endfoot
\bottomrule
\endlastfoot

\multicolumn{5}{@{}l@{}}{\textbf{\textcolor{financeblue}{Hypothèse H1 --- linéarité et proportionnalité}}} \\
\midrule
Student, constante & \ref{fiche:2} & \code{test\_intercept()} &
\code{test\_intercept()} (p exacte, loi $t$) & \code{Intercept} \\
Équivalence (TOST) & \ref{fiche:3} & \code{test\_tost\_intercept()} &
\code{test\_tost\_intercept()} (p exacte si $\Delta$ a priori) & --- \\
Student, pente & \ref{fiche:4} & \code{test\_lm\_complet()} &
\code{test\_lm\_complet()} (p exacte, loi $t$) & --- \\
Fisher global & \ref{fiche:5} & \code{test\_lm\_complet()} &
\code{test\_lm\_complet()} (p exacte, loi $F$) & --- \\
Coefficient $R^2$ & \ref{fiche:6} & \code{test\_lm\_complet()} &
indicateur, aucune p-value & --- \\
RESET & \ref{fiche:7} & \code{test\_reset()} &
\code{test\_reset()} (p asymptotique, loi $F$ approchée) & \code{RESET} \\
Spearman ($\times 2$) & \ref{fiche:8} & \code{stats::cor.test()} (appel direct) &
\code{stats::cor.test(..., exact = TRUE)} (p exacte, permutation) & \code{SpearVol} /
\code{SpearTps} \\
Mann--Kendall & \ref{fiche:9} & \code{test\_mann\_kendall()} &
\code{mk\_p\_exacte()} (p exacte, loi mahonienne, via \code{.mk\_loi\_exacte()},
\code{.poly\_mult()}) & \code{MK} \\
Cox--Stuart & \ref{fiche:10} & \code{test\_cox\_stuart()} &
\code{test\_cox\_stuart()} (p exacte, loi binomiale) & \code{CoxStuart} \\

\multicolumn{5}{@{}l@{}}{\textbf{\textcolor{financeblue}{Hypothèse H2 --- structure de variance quadratique}}} \\
\midrule
Breusch--Pagan (Koenker) & \ref{fiche:11} & \code{test\_breusch\_pagan()} &
\code{test\_breusch\_pagan()} (p asymptotique, $\chi^2$) & \code{BP} \\
Breusch--Pagan (1979) & \ref{fiche:12} & \code{test\_breusch\_pagan\_original()} &
idem (p asymptotique, $\chi^2$ sous normalité) &
\code{BP79} \\
White & \ref{fiche:13} & \code{test\_white()} &
\code{test\_white()} (p asymptotique, $\chi^2$) & \code{White} \\
Goldfeld--Quandt & \ref{fiche:14} & \code{test\_goldfeld\_quandt()} &
\code{test\_goldfeld\_quandt()} (p asymptotique, loi $F$ approchée) & \code{GQ} \\
Brown--Forsythe & \ref{fiche:15} & \code{test\_brown\_forsythe()} &
\code{test\_brown\_forsythe()} (p asymptotique, loi $F$ approchée) & \code{BF} \\
Smirnov (2 éch.) & \ref{fiche:16} & \code{stats::ks.test()} (appel direct) &
\code{stats::ks.test()} (p exacte combinatoire, sans ex \ae quo) & \code{Smirnov} \\
Position de $\hat\delta$ & \ref{fiche:17} & \code{usp\_ajuster()} (champ
\code{\$delta}) & diagnostic, aucune p-value & --- \\

\multicolumn{5}{@{}l@{}}{\textbf{\textcolor{financeblue}{Hypothèse H3 --- lognormalité des pertes agrégées}}} \\
\midrule
Shapiro--Wilk & \ref{fiche:18} & \code{.shapiro\_sur()} (encapsule
\code{stats::shapiro.test()}) & \code{.shapiro\_sur()} (p approx. Royston 1992) &
\code{SW} \\
Shapiro--Francia & \ref{fiche:19} & \code{test\_shapiro\_francia()} &
\code{test\_shapiro\_francia()} (p approx. Royston 1993) & \code{SF} \\
Anderson--Darling & \ref{fiche:20} & \code{stat\_ad()} &
\code{ad\_p\_stephens()} (p approx. D'Agostino \& Stephens) & \code{AD} \\
Cramér--von Mises & \ref{fiche:21} & \code{stat\_cvm()} &
\code{cvm\_p\_stephens()} (p approx. D'Agostino \& Stephens) & \code{CvM} \\
KS contre $\mathcal{N}(0,1)$ & \ref{fiche:22} & \code{stat\_ks()} &
formule de Kolmogorov, calculée en ligne dans \code{usp\_tests()} & \code{KS} \\
Lilliefors & \ref{fiche:23} & \code{stat\_lilliefors()} &
\code{lillie\_p()} (p approx. Dallal \& Wilkinson) & \code{Lillie} \\
Shapiro--Wilk sans Royston & \ref{fiche:24} & \code{.shapiro\_sur()} (statistique
$W$ réutilisée) & \code{sw\_p\_loi\_nulle()} (p exacte simulée, via
\code{sw\_loi\_nulle()}, mise en cache) & --- \\
Jarque--Bera & \ref{fiche:25} & \code{test\_jarque\_bera()} &
\code{test\_jarque\_bera()} (p asymptotique, $\chi^2_2$) & \code{JB} \\
D'Agostino (asymétrie) & \ref{fiche:26} & \code{test\_dagostino\_skew()} &
\code{test\_dagostino\_skew()} (p asymptotique, transf. $S_U$) & \code{DAgo} \\
Anscombe--Glynn (aplat.) & \ref{fiche:27} & \code{test\_anscombe\_kurt()} &
\code{test\_anscombe\_kurt()} (p asymptotique, Wilson--Hilferty) & --- \\

\multicolumn{5}{@{}l@{}}{\textbf{\textcolor{financeblue}{Hypothèse H4 --- indépendance et validité du maximum de vraisemblance}}} \\
\midrule
Durbin--Watson & \ref{fiche:28} & \code{stat\_dw()} &
\code{dw\_p\_exacte()} (p exacte, méthode d'Imhof via \code{.imhof\_p\_sup0()}) &
\code{DW} \\
Ljung--Box / Box--Pierce & \ref{fiche:29} & \code{stats::Box.test()} (appel
direct, 2 retards) & \code{stats::Box.test()} (p asymptotique, $\chi^2$) &
\code{LB1}, \code{LB2}, \code{BP2} \\
Suites (Wald--Wolfowitz) & \ref{fiche:30} & \code{test\_runs()} &
\code{runs\_p\_exacte()} (p exacte combinatoire, via \code{.runs\_dens()}) &
\code{Runs} \\
FOC (diagnostic) & \ref{fiche:31} & \code{usp\_ajuster()} (champ \code{\$foc}) &
diagnostic, aucune p-value & --- \\
Centrage des résidus (diagnostic) & \ref{fiche:32} & \code{mean()} sur $z$
(appel direct dans \code{usp\_tests()}) & diagnostic, aucune p-value & --- \\
Variance unitaire (diagnostic) & \ref{fiche:33} & \code{stats::var()} sur $z$
(appel direct dans \code{usp\_tests()}) & diagnostic, aucune p-value & --- \\

\multicolumn{5}{@{}l@{}}{\textbf{\textcolor{financeblue}{Stabilité, ruptures et points aberrants}}} \\
\midrule
Rupture sup-F & \ref{fiche:34} & \code{stat\_supF()} &
aucune (théorique seulement, Andrews 1993) & \code{supF} \\
OLS-CUSUM & \ref{fiche:35} & \code{stat\_cusum()} &
formule de Kolmogorov, calculée en ligne dans \code{usp\_tests()} & \code{CUSUM} \\
Grubbs & \ref{fiche:36} & \code{test\_grubbs()} &
\code{test\_grubbs()} (p asymptotique, borne de Bonferroni) & \code{Grubbs} \\
Rosner (ESD) & \ref{fiche:37} & \code{test\_rosner()} &
procédure de décision, aucune p-value & --- \\
Distance de Cook & \ref{fiche:38} & \code{stats::cooks.distance()} (appel
direct) & diagnostic, aucune p-value & --- \\
Leviers & \ref{fiche:39} & \code{stats::hatvalues()} (appel direct) &
diagnostic, aucune p-value & --- \\

\multicolumn{5}{@{}l@{}}{\textbf{\textcolor{financeblue}{Robustesse de l'estimation (méthode lognormale)}}} \\
\midrule
Rapport de vraisemblance sur $\delta$ & \ref{fiche:40} & fonction interne
\code{lr()} dans \code{usp\_profil()} & calcul direct \code{stats::pchisq()}
(p asymptotique, correction de Chernoff non appliquée) & --- \\
Jackknife & \ref{fiche:41} & \code{usp\_jackknife()} &
diagnostic, aucune p-value & --- \\
IC bootstrap & \ref{fiche:42} & \code{usp\_bootstrap()} (champ
\code{\$sigma\_boot}) & diagnostic, aucune p-value & --- \\
Convergence multi-démarrages & \ref{fiche:43} & \code{usp\_ajuster()} (champ
\code{\$part\_starts\_convergents}) & diagnostic, aucune p-value & --- \\

\multicolumn{5}{@{}l@{}}{\textbf{\textcolor{financeblue}{Méthode Merz--Wüthrich --- M1 : proportionnalité des cumulés (annexe XVII, D(2)(h)(iii))}}} \\
\midrule
Tendance avec le cumulé, à colonne donnée & \ref{mw:m1} &
\code{.mw\_lm\_intra()} (via \code{.mw\_pente\_intra()}) &
idem (p de Student, approchée) & \code{Intercept} \\
Ordonnée à l'origine par colonne & \ref{mw:m1a} &
\code{mw\_test\_ordonnee\_origine()} &
idem (p de Fisher, via \code{.fisher\_combine()}) & \code{Origine} \\
Homogénéité de $f_j$ & \ref{mw:m1b} & \code{mw\_test\_homogeneite\_f()} &
idem (p de Fisher) & \code{HomogF} \\
Absence de courbure & \ref{mw:m1c} & \code{mw\_test\_courbure()} &
idem (p de Fisher) & \code{Courbure} \\
Famille $\alpha$ & \ref{mw:m1d} & \code{mw\_famille\_alpha()} &
\textbf{aucune} (pas de loi analytique) & \code{Alpha} \\

\multicolumn{5}{@{}l@{}}{\textbf{\textcolor{financeblue}{Méthode Merz--Wüthrich --- M2 : variance proportionnelle au cumulé (annexe XVII, D(2)(h)(iv))}}} \\
\midrule
Breusch--Pagan sur résidus de Mack & \ref{mw:m2bp} &
\code{stats::lm(r\textasciicircum2 \textasciitilde\ C)} dans \code{mw\_tests()} &
\code{stats::pchisq()} (p asymptotique) & \code{BP} \\
Exposant de variance par colonne & \ref{mw:m2b} &
\code{mw\_test\_exposant\_variance()} & idem (p de Fisher) & \code{ExpVar} \\
Variance unitaire des résidus & \ref{mw:m2var} & \code{stats::var()} sur
\code{mw\_residus()} & diagnostic, aucune p-value & --- \\

\multicolumn{5}{@{}l@{}}{\textbf{\textcolor{financeblue}{Méthode Merz--Wüthrich --- M3 : indépendance (annexe XVII, D(2)(h)(i) et (ii))}}} \\
\midrule
Homogénéité entre années de survenance & \ref{mw:m3acc} &
\code{mw\_test\_homogeneite\_accident()} &
\code{stats::kruskal.test()} (p asymptotique) & \code{KruskalAcc} \\
Effets d'année calendaire & \ref{mw:m3cal} &
\code{mw\_test\_annees\_calendaires()} (via \code{.mack\_moments\_Z()}) &
idem (p asymptotique ; moments \emph{exacts}) & \code{Calendrier} \\
Corrélation entre années de développement & \ref{mw:m3cor} &
\code{mw\_stat\_correlation\_dev()} & \textbf{aucune} (pas de loi analytique) &
\code{CorrDev} \\
Autocorrélation (Durbin--Watson) & \ref{mw:m3dw} & \code{stat\_dw()} &
\textbf{aucune} (loi d'Imhof inapplicable aux résidus de Mack) & \code{DW} \\
Test des suites & \ref{mw:m3runs} & \code{test\_runs()} &
\code{test\_runs()} (approximation normale) & \code{Runs} \\

\multicolumn{5}{@{}l@{}}{\textbf{\textcolor{financeblue}{Méthode Merz--Wüthrich --- M4 : points aberrants}}} \\
\midrule
Cellule aberrante (Grubbs) & \ref{mw:m4gr} & \code{test\_grubbs()} &
\code{test\_grubbs()} (borne de Bonferroni) & \code{Grubbs} \\
Cellules aberrantes multiples (Rosner) & \ref{mw:m4ros} & \code{test\_rosner()} &
procédure de décision, aucune p-value & --- \\

\multicolumn{5}{@{}l@{}}{\textbf{\textcolor{financeblue}{Méthode Merz--Wüthrich --- M5 : normalité (diagnostic, non exigée par le modèle)}}} \\
\midrule
Shapiro--Wilk sur résidus de Mack & \ref{mw:m5sw} & \code{.shapiro\_sur()} &
\code{sw\_p\_loi\_nulle()} (p exacte simulée, via \code{sw\_loi\_nulle()}) & --- \\
Lilliefors sur résidus de Mack & \ref{mw:m5li} & \code{stat\_lilliefors()} &
\code{lillie\_p()} (approx. Dallal \& Wilkinson) & --- \\

\multicolumn{5}{@{}l@{}}{\textbf{\textcolor{financeblue}{Méthode Merz--Wüthrich --- M6 : robustesse de l'estimation}}} \\
\midrule
Extrapolation de $\hat\sigma^2_{J-1}$ & \ref{mw:m6sig} &
\code{mw\_ajuster()} (champ \code{\$sigma2}) & diagnostic, aucune p-value & --- \\
Concentration de la réserve & \ref{mw:m6conc} &
\code{mw\_ajuster()} (champ \code{\$reserve\_par\_annee}) &
diagnostic, aucune p-value & --- \\

\multicolumn{5}{@{}l@{}}{\textbf{\textcolor{financeblue}{Méthode Merz--Wüthrich --- chaîne de calcul (hors tests)}}} \\
\midrule
Facteurs et variances de développement & \S\ref{sec:mw} &
\code{mw\_ajuster()} & sans objet & --- \\
Erreur de prédiction à un an & \S\ref{sec:mw} & \code{mw\_msep()} &
sans objet & --- \\
Paramètre propre $\sigma_{(\text{res},s,\mathrm{USP})}$ & \S\ref{sec:mw} &
\code{mw\_parametre()} $\times$ \code{usp\_credibilite()} & sans objet & --- \\
Résidus standardisés de Mack & \S\ref{sec:mw} & \code{mw\_residus()} &
sans objet & --- \\
Bootstrap semi-paramétrique & \S\ref{sec:mw} & \code{mw\_bootstrap()}
(via \code{mw\_simuler\_triangle()} et \code{.mw\_stats()}) &
sans objet & toutes les clés ci-dessus \\
Leviers et DFBETA du triangle & \ref{fiche:39} & \code{mw\_influence()} &
diagnostic, aucune p-value & --- \\
Contributions par année de survenance & \ref{fiche:39} &
\code{mw\_contributions()} & diagnostic, aucune p-value & --- \\

\end{longtable}
\end{center}
\end{landscape}

\begin{encadre}
\textbf{Fonctions internes de support.} Trois fonctions préfixées d'un point ne
sont jamais appelées directement par \code{usp\_tests()} mais seulement par les
fonctions de p-value ci-dessus~: \code{.imhof\_p\_sup0()} (intégrale d'Imhof,
utilisée par \code{dw\_p\_exacte()}), \code{.mk\_loi\_exacte()} et
\code{.poly\_mult()} (distribution exacte du $S$ de Kendall par produit de
polynômes, utilisées par \code{mk\_p\_exacte()}), et \code{.runs\_dens()}
(densité exacte du nombre de suites, utilisée par \code{runs\_p\_exacte()}).
\code{.shapiro\_sur()} et \code{sw\_loi\_nulle()} sont deux autres fonctions
internes~: la première encapsule \code{stats::shapiro.test()} pour éviter
qu'un échantillon dégénéré n'interrompe la boucle de bootstrap, la seconde
met en cache la loi nulle simulée de $W$ par valeur de $T$.
\end{encadre}

\begin{encadre}
\textbf{Lecture de la colonne « Clé MC ».} Un tiret signifie qu'aucune p-value
de Monte-Carlo n'est calculée pour ce test~: soit parce que l'hypothèse nulle
n'appartient pas au modèle simulé (Student sur la pente, Fisher)~; soit parce
que le test n'est pas de nature à produire de p-value (indicateurs, procédures
de décision, diagnostics numériques)~; soit, troisième motif, parce que la
grandeur est \emph{rivée par l'estimation} et qu'il n'y a donc rien à simuler
--- la loi simulée porterait un atome en la valeur observée (centrage et
variance unitaire des résidus standardisés, fiches~\ref{fiche:32}
et~\ref{fiche:33}). Dans tous les autres cas, la clé
correspond au nom de colonne retourné par \code{.stats\_bootstrapables()} et
consommé par \code{usp\_bootstrap()}.
\end{encadre}


\newpage
\clearpage
% =============================================================================
\section{Notations générales}
% =============================================================================

L'outil couvre deux familles de méthodes, aux géométries de données distinctes.
Les notations sont donc présentées séparément~: la sous-section suivante concerne
les méthodes lognormales (sections~B et~C de l'annexe~XVII), la
section~\ref{sec:notations-mw} la méthode Merz--Wüthrich (section~D). Les
conventions communes --- lois, fonctions de répartition, convention de
restitution --- valent pour les deux.

\subsection*{Méthodes lognormales : données et modèle}

Pour un segment donné et une profondeur $T$~:
\begin{itemize}[itemsep=1pt]
  \item $x_1,\dots,x_T$ : primes acquises par année de survenance (méthode risque de prime) ou
        meilleure estimation de la provision en début d'exercice (méthode risque de réserve~1) ;
  \item $y_1,\dots,y_T$ : pertes agrégées (méthode prime) ou paiements de l'exercice augmentés de
        la meilleure estimation de clôture sur les mêmes sinistres (méthode réserve~1) ;
  \item $r_t = y_t / x_t$ : ratio observé (ratio S/P, ou ratio de liquidation) ;
  \item $\bar x = T^{-1}\sum_t x_t$.
\end{itemize}

Le modèle de l'annexe~XVII postule $Y_t \mid x_t \sim \mathcal{LN}(\mu_t,\psi_t)$ avec
\[
  \mathbb{E}[Y_t] = \beta x_t,
  \qquad
  \operatorname{Var}(Y_t) = \sigma^2\big[(1-\delta)\,\bar x\, x_t + \delta\, x_t^2\big],
  \qquad \delta\in[0,1],\ \ \sigma = \beta e^{\gamma}.
\]
On en déduit la précision logarithmique et l'estimateur du niveau~:
\[
  \pi_t(\delta,\gamma) = \Big[\ln\!\big(1 + e^{2\gamma}(\delta + (1-\delta)\bar x / x_t)\big)\Big]^{-1},
  \qquad
  \ln\hat\beta = \frac{\tfrac{T}{2} + \sum_t \pi_t \ln r_t}{\sum_t \pi_t}.
\]

\subsection*{Résidus}

Les tests des sections~5 à~8 portent sur les \emph{résidus normalisés}
\[
  z_t \;=\; \sqrt{\hat\pi_t}\;\Big(\ln r_t + \tfrac{1}{2\hat\pi_t} - \ln\hat\beta\Big),
  \qquad t = 1,\dots,T ,
\]
qui, si le modèle est correctement spécifié, sont \emph{approximativement} indépendants et de loi
$\mathcal{N}(0,1)$. On note $z_{(1)} \le \dots \le z_{(T)}$ la statistique d'ordre associée,
$\bar z$ la moyenne empirique, $s_z^2 = (T-1)^{-1}\sum_t (z_t - \bar z)^2$ la variance empirique
corrigée, et $p_{(i)} = \Phi(z_{(i)})$.

\subsection*{Lois et fonctions de répartition}

$\Phi$ et $\varphi$ désignent la fonction de répartition et la densité de $\mathcal{N}(0,1)$ ;
$F_{t,\nu}$, $F_{\chi^2_\nu}$, $F_{\mathcal{F},\nu_1,\nu_2}$ celles des lois de Student, du
khi-deux et de Fisher. Les moments empiriques centrés d'ordre $k$ sont notés
$m_k = T^{-1}\sum_t (z_t - \bar z)^k$, l'asymétrie $\sqrt{b_1} = m_3/m_2^{3/2}$ et
l'aplatissement $b_2 = m_4/m_2^{2}$.

\subsection*{Régression auxiliaire}

Plusieurs tests de la section~4 s'appuient sur la régression auxiliaire non contrainte
$y_t = a + b\,x_t + \varepsilon_t$, dont on note
$S_{xx} = \sum_t (x_t-\bar x)^2$, $S_{xy} = \sum_t (x_t-\bar x)(y_t-\bar y)$,
$\hat b = S_{xy}/S_{xx}$, $\hat a = \bar y - \hat b \bar x$, les résidus
$e_t = y_t - \hat a - \hat b x_t$, ainsi que
$\mathrm{SCR} = \sum_t e_t^2$, $\mathrm{SCT} = \sum_t (y_t - \bar y)^2$ et
$\hat s^2 = \mathrm{SCR}/(T-2)$.

\medskip
\begin{encadre}
\textbf{Attention.} Cette régression auxiliaire n'est \emph{pas} le modèle réglementaire, qui
impose $\mathbb{E}[Y_t] = \beta x_t$ sans constante. Elle sert uniquement de support de test aux
diagnostics de la section~4.
\end{encadre}

\subsection*{Méthode Merz--Wüthrich : données et modèle}
\label{sec:notations-mw}

La donnée d'entrée n'est plus un couple de vecteurs mais un \emph{triangle} de
paiements cumulés. Les notations sont celles de l'annexe~XVII, section~D,
paragraphe~3~:

\begin{itemize}[itemsep=1pt]
  \item $i = 0,\dots,I$ : années de survenance, la plus ancienne portant
        l'indice~0 ;
  \item $j = 0,\dots,J$ : années de développement ;
  \item $C(i,j)$ : sinistres cumulés de l'année de survenance $i$ à l'année de
        développement $j$, observés pour $i + j \le I$ ;
  \item $F(i,j) = C(i,j+1)/C(i,j)$ : facteur de développement individuel ;
  \item $n_j = I - j$ : nombre d'observations disponibles dans la colonne $j$.
\end{itemize}

\noindent
Le modèle de chaîne d'échelle postule, pour tout $i$ et tout $j$~:
\[
  \mathbb{E}\big[C(i,j+1)\mid C(i,j)\big] = f_j\,C(i,j),
  \qquad
  \operatorname{Var}\big[C(i,j+1)\mid C(i,j)\big] = \sigma_j^2\,C(i,j),
\]
les années de survenance étant mutuellement indépendantes. \textbf{Aucune loi
n'est spécifiée}~: seuls les deux premiers moments conditionnels le sont, ce qui
constitue la différence de nature la plus importante avec les méthodes
lognormales et commande l'ensemble des choix de p-values.

Les estimateurs sont
\[
  \hat f_j = \frac{\sum_{i=0}^{I-j-1} C(i,j+1)}{\sum_{i=0}^{I-j-1} C(i,j)},
  \qquad
  \hat\sigma_j^2 = \frac{1}{I-j-1}\sum_{i=0}^{I-j-1} C(i,j)
  \big(F(i,j) - \hat f_j\big)^2 ,
\]
la variance de la dernière année de développement étant \emph{extrapolée} et non
estimée (section~\ref{sec:mw}).

\subsection*{Résidus des deux méthodes}

Chaque méthode a son résidu de référence, support de ses tests~:

\begin{center}
\small
\rowcolors{2}{tablerow}{white}
\begin{tabular}{@{}p{3.0cm}p{5.6cm}p{6.3cm}@{}}
\rowcolor{tableheader}
\textcolor{financeblue}{\textbf{Méthode}} &
\textcolor{financeblue}{\textbf{Résidu}} &
\textcolor{financeblue}{\textbf{Loi sous le modèle}} \\
Lognormale & $z_t = \sqrt{\hat\pi_t}\big(\ln r_t + \tfrac{1}{2\hat\pi_t} - \ln\hat\beta\big)$ &
i.i.d. $\mathcal{N}(0,1)$~: la loi est entièrement spécifiée \\
Merz--Wüthrich & $r(i,j) = \dfrac{\sqrt{C(i,j)}}{\hat\sigma_j}\big(F(i,j) - \hat f_j\big)$ &
centré, d'échelle approximativement unitaire, non corrélé~; \textbf{aucune loi
n'est spécifiée} \\
\end{tabular}
\end{center}

\noindent
Cette asymétrie est déterminante~: pour la méthode lognormale, la normalité des
résidus est une \emph{hypothèse à tester} (H3)~; pour la méthode Merz--Wüthrich,
elle n'est qu'un \emph{diagnostic facultatif} (M5).

\subsection*{Convention de restitution : statistique, estimation, p-value}

Les sorties du script distinguent explicitement trois grandeurs, qui ne doivent jamais être
confondues et figurent dans des champs séparés~:

\begin{center}
\small
\begin{tabular}{@{}p{2.3cm}p{12.9cm}@{}}
\toprule
\code{stat} & \textbf{Statistique de test}~: la quantité dont la loi sous $H_0$ est connue (ou
simulée). C'est elle, et elle seule, qui est comparée à une loi de référence pour produire une
p-value. Le champ \code{loi} rappelle cette loi. \\
\code{estim} & \textbf{Estimation ou grandeur descriptive}~: pente, constante, $\hat\rho_s$,
$\hat\delta$, $R^2$, nombre de valeurs aberrantes, distance de Cook maximale\dots Ces quantités
n'ont \emph{aucune} loi de référence dans ce dispositif et ne doivent pas être lues comme des
statistiques de test. \\
\code{p}, \code{p\_mc} & \textbf{p-values}~: respectivement asymptotique (ou exacte selon le
test) et par bootstrap paramétrique. \\
\bottomrule
\end{tabular}
\end{center}

Le champ \code{type} classe chaque ligne en \emph{test} (hypothèse nulle et loi de référence),
\emph{indicateur} (descriptif, comme $R^2$), \emph{diagnostic} (numérique ou d'influence),
\emph{procédure de décision} (Rosner, qui ne produit pas de p-value) ou \emph{non applicable}.
Le champ \code{sens} indique enfin si l'on souhaite rejeter ou ne pas rejeter $H_0$~: pour les
tests de Student sur la pente et de Fisher, un $p$ faible est un \emph{bon} résultat, et le
verdict du script est inversé en conséquence.

\clearpage
% =============================================================================
\section{Contrôles de qualité des données}
\label{sec:controles}
% =============================================================================

Les contrôles de recevabilité ne sont \textbf{pas} des tests statistiques~: ils ne
comportent ni hypothèse nulle, ni statistique, ni loi de référence. Ce sont des
vérifications déterministes au titre des articles~19 et~219 du règlement délégué
et des exigences de données propres à chaque méthode. Chaque méthode a sa
fonction de contrôle dans le moteur, exécutable indépendamment de l'interface~:
\code{engine\_valider\_donnees()} et \code{usp\_controle\_donnees()} pour les
méthodes lognormales, \code{mw\_valider\_triangle()} pour la méthode
Merz--Wüthrich.

\subsection*{Méthodes lognormales (annexe XVII, section B/C, paragraphe 2)}

\begin{center}
\small
\rowcolors{2}{tablerow}{white}
\begin{tabular}{@{}p{4.4cm}p{5.0cm}p{5.4cm}@{}}
\rowcolor{tableheader}
\textcolor{financeblue}{\textbf{Contrôle}} &
\textcolor{financeblue}{\textbf{Critère}} &
\textcolor{financeblue}{\textbf{Justification}} \\
Profondeur minimale & $T \ge 5$ & Annexe XVII, B/C(2)(b). Bloquant. \\
Valeurs manquantes & aucun \code{NA} & Art.~19 : exhaustivité. \\
Positivité de $x_t$ & $\min_t x_t > 0$ & $r_t = y_t/x_t$ doit être défini. \\
Positivité de $y_t$ & $\min_t y_t > 0$ & Support de la loi lognormale. \\
Doublons & aucun couple $(x_t,y_t)$ répété & Signale une duplication de ligne. \\
Plausibilité du ratio & $r_t \in\, ]0\,;5[$ & Erreurs d'unité ou de saisie. \\
Amplitude du volume & $\max_t x_t/\min_t x_t < 10$ & Rupture de périmètre probable. \\
Crédibilité pleine & $T \ge 10$ (barème court) & Informatif, section~G annexe~XVII. \\
\end{tabular}
\end{center}

\subsection*{Méthode Merz--Wüthrich (annexe XVII, section D, paragraphe 2)}

\begin{center}
\small
\rowcolors{2}{tablerow}{white}
\begin{tabular}{@{}p{4.4cm}p{5.0cm}p{5.4cm}@{}}
\rowcolor{tableheader}
\textcolor{financeblue}{\textbf{Contrôle}} &
\textcolor{financeblue}{\textbf{Critère}} &
\textcolor{financeblue}{\textbf{Justification}} \\
Années de survenance & $I + 1 \ge 5$ & D(2)(b). Bloquant. \\
Années de développement & $J + 1 \ge 5$ & D(2)(c). Bloquant. \\
Forme du triangle & $I + 1 \ge J + 1$ & D(2)(e). Bloquant. \\
Géométrie carrée & $I = J$ & Restriction d'implémentation~: la formule du
paragraphe~4 fait intervenir $C(i,I-i)$, indéfinie si $I - i > J$. Bloquant. \\
Cellules observées & aucune valeur manquante pour $i + j \le I$ & Le
chain-ladder exige la partie supérieure gauche complète. \\
Positivité des cumulés & $C(i,j) > 0$ & Les facteurs $F(i,j)$ et la pondération
en $C(i,j)$ doivent être définis. \\
Monotonie des cumulés & $C(i,j+1) \ge C(i,j)$ & Avertissement~: un recul traduit
un recouvrement ou un boni de liquidation, licite mais à signaler. \\
Crédibilité pleine & $I + 1 \ge 10$ (barème court) & Informatif~; à défaut,
$\hat\sigma_j^2$ repose sur deux ou trois observations en fin de triangle. \\
\end{tabular}
\end{center}

\medskip
\begin{encadre}
\textbf{Différence de nature.} Les contrôles lognormaux portent sur la
\emph{plausibilité économique} des séries (ratio dans une plage admissible,
amplitude des volumes). Ceux de la méthode Merz--Wüthrich portent d'abord sur la
\emph{géométrie} du triangle, dont dépend la simple définition des estimateurs~:
un triangle non carré ou incomplet ne permet pas d'appliquer la formule du
paragraphe~4. C'est pourquoi la plupart sont bloquants et non de simples
avertissements.
\end{encadre}

\newpage

\partie{Partie II}{Cadre statistique général}

% =============================================================================
\section{Bootstrap paramétrique et p-values de Monte-Carlo}
\label{sec:mc}
% =============================================================================

La profondeur exigée par l'annexe~XVII (cinq années au minimum) rend inutilisables, en toute
rigueur, les lois asymptotiques de la majorité des tests documentés ci-après. Le script
recalcule donc, pour les statistiques les plus sensibles, une p-value de Monte-Carlo par
bootstrap paramétrique.

\begin{fiche}{Procédure générale de p-value de Monte-Carlo}
\refl{B.~Efron, \emph{Bootstrap methods: another look at the jackknife}, \emph{Annals of
Statistics}, 7(1), 1979, p.~1--26 ; G.\,A. Barnard, discussion, \emph{JRSS~B}, 25, 1963, p.~294 ;
J.\,G. MacKinnon, \emph{Bootstrap hypothesis testing}, in \emph{Handbook of Computational
Econometrics}, Wiley, 2009.}
\Cadre
On dispose d'un modèle paramétrique entièrement estimé
$\widehat{M} = (\hat\delta,\hat\gamma,\hat\beta)$ et d'une statistique $S$ dont la loi exacte sous
$H_0$ est inconnue ou mal approchée. L'hypothèse de validité, non testée, est que $\widehat{M}$
constitue une approximation suffisante du vrai modèle sous $H_0$ (bootstrap \emph{paramétrique},
par opposition au rééchantillonnage non paramétrique).
\Hzero
Celle du test sous-jacent, inchangée~: la statistique $S$ a été calculée sur des données conformes
au modèle.
\Stat
$S_{\text{obs}}$, calculée sur l'échantillon réel. Pour $b = 1,\dots,B$, on simule
$y_t^{(b)} \sim \mathcal{LN}\!\big(\ln(\hat\beta x_t) - \tfrac{1}{2\hat\pi_t},\, \hat\pi_t^{-1}\big)$,
on \textbf{réestime intégralement} $(\delta,\gamma)$ sur $(x_t, y_t^{(b)})$, et l'on recalcule
$S^{(b)}$.
\Loi
La loi empirique de $\{S^{(1)},\dots,S^{(B)}\}$ est, par construction, un échantillon de la loi
exacte de $S$ sous $H_0$ conditionnellement au modèle ajusté -- \emph{y compris} l'effet de
l'estimation des paramètres, source principale du biais des lois asymptotiques à petit $T$.
\Pval
Pour un test unilatéral à droite,
\[
  \hat p_{\text{MC}} = \frac{1 + \#\{b : S^{(b)} \ge S_{\text{obs}}\}}{B+1},
\]
avec l'adaptation évidente à gauche, et
$\hat p = 2\min(\hat p_{\text{gauche}}, \hat p_{\text{droite}})$ au bilatéral.

\medskip
La portée de l'ajout de~1 au numérateur et au dénominateur mérite d'être énoncée avec précision,
car elle est souvent surestimée. Deux situations doivent être distinguées.

\emph{Lorsque la loi de $S$ sous $H_0$ est exactement simulable} --- cas des tests de
permutation et de randomisation --- les $B$ valeurs simulées et la valeur observée sont
échangeables sous $H_0$. Le rang de $S_{\text{obs}}$ est alors uniforme sur $\{1,\dots,B+1\}$ et
la règle ci-dessus fournit un test \textbf{exactement} de niveau $\alpha$ dès que $\alpha(B+1)$
est entier (Dwass, 1957~; Hope, 1968). Vérification numérique sur 200\,000
réplications avec $B = 99$~: la fréquence de rejet vaut $0{,}0506$ pour $\alpha = 0{,}05$ et
$0{,}1003$ pour $\alpha = 0{,}10$.

\emph{Dans le cas présent}, la loi de référence n'est pas connue~: elle est celle du
\emph{modèle estimé} $\widehat{M}$. Les $S^{(b)}$ sont tirés sous $\widehat{M}$ et non sous le vrai
modèle, de sorte que l'échangeabilité avec $S_{\text{obs}}$ n'est pas acquise. Le résultat
d'exactitude ci-dessus \textbf{ne se transporte donc pas} au bootstrap paramétrique. L'ajout
de~1 corrige la seule discrétisation liée à un $B$ fini~; il ne corrige pas l'écart entre la loi
sous $\widehat{M}$ et la loi sous le vrai modèle, dont la disparition est un résultat
\emph{asymptotique en $T$} (Bickel \& Freedman, 1981~; Singh, 1981~; Hall, 1992). La validité du
dispositif reste par conséquent asymptotique, et l'étude de simulation de la
section~\ref{sec:verif} en constitue la seule mesure disponible à $T = 8$.
\Usage
C'est la seule p-value réellement calibrée aux tailles d'échantillon du régime USP. Sa
granularité est bornée par $1/(B+1)$~: avec $B = 999$, la plus petite p-value atteignable vaut
$0{,}001$. Sa limite est qu'elle teste la conformité au \emph{modèle ajusté}~: si le modèle est
mal spécifié dans une direction que $S$ ne détecte pas, la p-value sera élevée à tort. Dans les
sorties du script, elle figure en colonne \code{p\_mc} et prime sur la p-value asymptotique
lorsque les deux sont disponibles.
\end{fiche}

\subsection{Bootstrap semi-paramétrique de la méthode Merz--Wüthrich}
\label{sec:boot-mw}

Le modèle de Mack ne spécifiant aucune loi, aucun bootstrap \emph{paramétrique}
n'est possible. Le moteur emploie un rééchantillonnage des résidus standardisés
(England \& Verrall, 2002, \emph{British Actuarial Journal} 8(3), 443--518), qui
ne suppose que leur échangeabilité~: les p-values correspondantes sont donc
\emph{semi-paramétriques} et non paramétriques.

\begin{encadre}
\textbf{Limite structurelle du rééchantillonnage de résidus.} Cette procédure
tire dans la loi \emph{empirique} des résidus observés. Son hypothèse nulle est
« les résidus suivent leur propre loi empirique », et non « les résidus sont
normaux ». Une p-value de Monte-Carlo pour un test de normalité y serait donc
vide de sens~; la vérification par simulation le confirme, la fréquence de rejet
tombant à zéro au lieu de 5~\%. Les tests de normalité de la famille M5 sont en
conséquence exclus du bootstrap et rapportés avec leur propre loi nulle~: loi
exacte de $W$ évaluée par simulation directe pour Shapiro--Wilk, approximation de
Dallal \& Wilkinson pour Lilliefors.
\end{encadre}


\clearpage
% =============================================================================
\section{Nature des p-values et validité à $T = 8$}
\label{sec:nature}
% =============================================================================

\subsection{Quatre statuts épistémiques à ne jamais confondre}

Toute conclusion du dispositif relève de l'un des quatre statuts suivants. La
distinction est reprise systématiquement dans les tableaux~\ref{tab:nature} et
\ref{tab:conv}.

\begin{center}
\small
\begin{tabular}{@{}p{3.0cm}p{12.3cm}@{}}
\toprule
\textbf{Statut} & \textbf{Signification et portée} \\
\midrule
\textbf{Exact} & La loi de la statistique sous $H_0$ est connue pour la taille
d'échantillon considérée, sans approximation (loi combinatoire, de permutation,
binomiale, ou loi de Student/Fisher sous normalité exacte des erreurs). La
p-value est alors valide à $T = 8$ comme à $T = 1000$. \\
\textbf{Asymptotique} & La loi n'est connue qu'à la limite $T \to \infty$.
L'existence du résultat limite ne dit \emph{rien} de la qualité de
l'approximation à $T$ fini~: il faut disposer, séparément, d'un résultat sur la
\emph{vitesse} de convergence ou d'une étude de simulation. \\
\textbf{Approximation numérique} & La loi exacte existe mais n'est pas
calculable analytiquement, et l'on emploie une formule d'ajustement (Royston
pour Shapiro--Wilk, Stephens pour Anderson--Darling). La qualité de
l'approximation dépend du domaine de validité annoncé par ses auteurs, qui doit
être vérifié pour $T = 8$. \\
\textbf{Comportement empirique} & Résultat d'une étude de simulation. Il
documente le niveau effectif observé, mais ne transforme jamais une
approximation en résultat exact et ne vaut que pour le modèle sous lequel la
simulation a été conduite. \\
\bottomrule
\end{tabular}
\end{center}

\medskip
\begin{encadre}
\textbf{Principe de sélection retenu dans le moteur.} Pour chaque test, la
p-value effectivement utilisée pour le verdict suit la hiérarchie
\emph{exacte $\succ$ Monte-Carlo $\succ$ asymptotique}. Le champ \code{nature\_p}
de la sortie indique laquelle a été retenue, et les trois valeurs
(\code{p\_exacte}, \code{p\_asymptotique}, \code{p\_monte\_carlo}) restent
affichées séparément pour permettre la revue. Aucune p-value asymptotique n'est
retenue lorsqu'une p-value exacte ou de Monte-Carlo est disponible.

\medskip
\textbf{Ce que cette hiérarchie ne dit pas~: la règle de sortie.} La hiérarchie
\emph{ordonne des p-values qui existent}. Elle ne se lit pas à l'envers~:
l'absence de p-value exacte et de p-value de Monte-Carlo n'autorise pas un
repli sur l'asymptotique, ni, plus généralement, sur la loi \emph{nominale}
--- celle que la statistique suivrait si les paramètres n'avaient pas été
estimés. Lorsqu'une grandeur est \textbf{rivée par l'estimation}, c'est-à-dire
fixée par les conditions du premier ordre du maximum de vraisemblance plutôt
que librement observée, aucune des trois voies n'est ouverte, et il faut le
dire des trois et non de la seule qui a échoué en premier~: \emph{(i)} il
n'existe pas de loi exacte sous $H_0$ pour une quantité que $H_0$ ne laisse pas
varier~; \emph{(ii)} la loi simulée sous le modèle ajusté porte un atome en la
valeur observée, de sorte que la p-value de Monte-Carlo se décide au signe d'un
bruit numérique~; \emph{(iii)} la loi nominale ne décrit pas la grandeur
restituée, puisque la valeur observée n'est pas une observation. Le point
\emph{(iii)} est le plus important pour le relecteur, car c'est le repli le
plus tentant~: à $\hat\pi_t$ constant, la p-value nominale de ces grandeurs est
une \emph{constante de $T$} qui ne dépend plus des données, identique pour deux
jeux étrangers l'un à l'autre. Le « OK » qu'elle produirait ne pourrait être
perdu par aucune donnée~; c'est exactement la fausse assurance que la
restitution en diagnostic vise à éviter. Le calcul et les mesures figurent aux
fiches~\ref{fiche:32} et~\ref{fiche:33}. Ces lignes sont donc restituées comme
\emph{diagnostics}, sans p-value retenue et avec le motif. Cette règle de
sortie ne concerne ni le test de Student sur la pente ni le test de Fisher
global~: leur statistique est librement observée et leur loi est exacte sous
normalité des erreurs~; ils conservent leur p-value, pour la raison exposée à
la sous-section~\ref{sec:t8}.
\end{encadre}

\subsection{Deux sources d'erreur distinctes}
\label{sec:deuxerreurs}

Lorsqu'une p-value est estimée par $B$ simulations, deux erreurs se
superposent~; elles ont des causes, des ordres de grandeur et des remèdes
différents, et ne doivent en aucun cas être confondues.

\begin{center}
\small
\begin{tabular}{@{}p{3.3cm}p{5.8cm}p{5.8cm}@{}}
\toprule
 & \textbf{Erreur de Monte-Carlo} & \textbf{Erreur d'approximation statistique} \\
\midrule
Source & Nombre fini $B$ de réplications & Taille finie $T$ de l'échantillon \\
Grandeur affectée & L'estimateur $\hat p_{\text{MC}}$ de la p-value bootstrap
& L'écart entre la loi bootstrap et la vraie loi sous $H_0$ \\
Ordre de grandeur & $\sqrt{p(1-p)/B}$, soit $\mathcal{O}(B^{-1/2})$ &
Dépend du test~; typiquement $\mathcal{O}(T^{-1/2})$ ou $\mathcal{O}(T^{-1})$ \\
Réduction & Augmenter $B$ (gratuit, purement numérique) &
\textbf{Impossible}~: $T = 8$ est une donnée du problème \\
Mesure dans le moteur & champ \code{erreur\_MC} & non mesurable directement~;
approchée par l'étude de simulation de la section~11 \\
\bottomrule
\end{tabular}
\end{center}

Avec $B = 999$ et $\hat p \approx 0{,}10$, l'erreur de Monte-Carlo vaut environ
$0{,}009$. Il serait erroné d'en conclure que la p-value est « précise à
$\pm 0{,}01$ »~: cette précision porte sur l'estimation de la p-value
\emph{bootstrap}, pas sur son écart à la vraie p-value sous $H_0$. La granularité
de $\hat p_{\text{MC}}$ est en outre bornée inférieurement par $1/(B+1)$
(Dwass, 1957 ; Hope, 1968 ; voir la bibliographie complémentaire, section~\ref{sec:biblio}).

\subsection{Ce que l'on peut et ne peut pas affirmer à $T = 8$}
\label{sec:t8}

\begin{itemize}[itemsep=3pt]
  \item \textbf{Résultats exacts disponibles à $T = 8$.} Loi binomiale
  (Cox--Stuart), loi de permutation de Spearman, loi mahonienne du $S$ de
  Kendall, loi combinatoire du nombre de suites, loi exacte de Smirnov à deux
  échantillons, et les lois de Student et de Fisher de la régression auxiliaire
  \emph{sous l'hypothèse de normalité exacte des erreurs}. Ces p-values sont
  valides sans réserve de taille d'échantillon.
  \item \textbf{Résultats asymptotiques sans justification suffisante à
  $T = 8$.} Les statistiques de type multiplicateur de Lagrange
  (Breusch--Pagan, White), portmanteau (Ljung--Box, Box--Pierce),
  Jarque--Bera, et les approximations normales de Mann--Kendall,
  Wald--Wolfowitz et D'Agostino. La littérature disponible ne fournit, pour
  aucune d'entre elles, de borne d'erreur exploitable à $T = 8$.
  \item \textbf{Deux cas où aucune p-value de Monte-Carlo n'est possible.} Les
  tests de Student sur la pente et de Fisher global ont pour hypothèse nulle
  $b = 0$. Or le bootstrap paramétrique simule sous le \emph{modèle ajusté},
  pour lequel $\hat b \neq 0$~: le modèle de simulation appartient à $H_1$ et
  non à $H_0$. Une p-value de Monte-Carlo y serait dépourvue de sens. Ces deux
  tests conservent donc leur loi de Fisher/Student, qui est \emph{exacte} sous
  normalité des erreurs~: la limitation porte sur l'hypothèse de normalité, non
  sur la taille d'échantillon.
  \item \textbf{Un troisième cas, de nature différente~: la grandeur rivée par
  l'estimation.} Le centrage et la variance unitaire des résidus standardisés
  (fiches~\ref{fiche:32} et~\ref{fiche:33}) sont fixés par les conditions du
  premier ordre du maximum de vraisemblance. La loi simulée sous le modèle
  ajusté existe bien, mais elle porte un \emph{atome} en la valeur observée~:
  la p-value de Monte-Carlo s'y décide au signe d'un bruit numérique. À la
  différence des deux cas précédents, il n'y a ici aucun repli possible~: ni
  loi exacte, ni loi nominale utilisable, cette dernière produisant à
  $\hat\pi_t$ constant une p-value qui ne dépend plus des données. Ces deux
  lignes sont restituées comme \emph{diagnostics}, sans p-value retenue et avec
  le motif.
  \item \textbf{Ce qu'aucune méthode ne restitue.} Ni le bootstrap ni les lois
  exactes n'augmentent la \emph{puissance}. À $T = 8$, l'ensemble du dispositif
  reste très peu puissant~: une p-value élevée n'est pas une validation
  d'hypothèse. Les diagnostics de robustesse de la section~10 conservent la
  priorité sur les tests d'hypothèse.
\end{itemize}

\newpage
% -----------------------------------------------------------------------------
\subsection{Tableau 1 : disponibilité et choix des p-values}
\label{tab:nature}
% -----------------------------------------------------------------------------

Ce tableau recense les tests \emph{effectivement présents} dans le moteur
(\code{R/engine.R}). La colonne « Asympt. utilisable à $T = 8$ ? » ne répond pas
mécaniquement~: elle distingue l'existence théorique d'une loi limite de la
qualité de son approximation à $T = 8$.

La colonne « Méthode » indique le périmètre d'emploi~: \textbf{LN} pour les
méthodes lognormales (risque de prime, section~B, et risque de réserve
n\textsuperscript{o}~1, section~C), \textbf{MW} pour la méthode du risque de
réserve n\textsuperscript{o}~2 (Merz--Wüthrich, section~D), et
\textbf{Les deux} lorsque le test est employé dans les deux dispositifs. Dans ce
dernier cas la statistique est identique, mais son support diffère~: résidus
normalisés $z_t$ d'un côté, résidus standardisés de Mack $r(i,j)$ de l'autre
(voir la section~\ref{sec:mw}).

\begin{landscape}
\begin{center}
\scriptsize
\setlength{\tabcolsep}{5pt}
\renewcommand{\arraystretch}{1.20}
\begin{longtable}{@{}p{2.9cm}p{1.5cm}p{1.1cm}p{1.3cm}p{1.3cm}p{1.3cm}p{2.7cm}p{4.7cm}p{4.1cm}@{}}
\toprule
\textbf{Test} & \textbf{Méthode} & \textbf{Stat.} & \textbf{Exacte ?} & \textbf{Asympt. ?} &
\textbf{MC ?} & \textbf{Voie retenue} &
\textbf{Asympt. utilisable à $T=8$ ?} & \textbf{Recommandation ACPR} \\
\midrule
\endfirsthead
\toprule
\textbf{Test} & \textbf{Méthode} & \textbf{Stat.} & \textbf{Exacte ?} & \textbf{Asympt. ?} &
\textbf{MC ?} & \textbf{Voie retenue} &
\textbf{Asympt. utilisable à $T=8$ ?} & \textbf{Recommandation ACPR} \\
\midrule
\endhead
\multicolumn{9}{@{}l@{}}{\textit{Suite page suivante}} \\
\endfoot
\bottomrule
\endlastfoot

Student, constante & LN & $t_a$ & Oui (sous normalité) & --- & Oui &
\textbf{Exacte} (MC en recoupement) &
Sans objet : loi exacte & Distribution exacte à privilégier. \emph{Mais} le non-rejet ne
prouve rien : compléter par le TOST \\

Équivalence (TOST) & LN & $\max$ des 2 unilatéraux & Oui si $\Delta$ fixée a priori & --- &
Sans objet & \textbf{Exacte} ($\Delta$ a priori) ou quasi-exacte ($\Delta = \theta\bar y$) &
Sans objet : loi exacte & \textbf{Seul test fournissant une preuve positive} de
proportionnalité ; fixer $\Delta$ a priori. Puissance de 46~\% seulement à $T=8$ \\

Student, pente & LN & $t_b$ & Oui (sous normalité) & --- & Non (voir texte) &
Exacte (loi $t$) & Sans objet : loi exacte & Distribution exacte à privilégier ;
vérifier la normalité \\

Fisher global & LN & $F$ & Oui (sous normalité) & --- & Non (voir texte) &
Exacte (loi $F$) & Sans objet : loi exacte & Idem ; redondant avec $t_b$ \\

RESET & LN & $F$ & \textbf{Non} & Oui & Oui & Monte-Carlo &
\textbf{Déconseillée} : $\hat y^2,\hat y^3$ dépendent de $y$, la loi $F$ n'est
pas exacte & Bootstrap recommandé \\

Spearman ($\times 2$) & LN & $S$ & Oui ($T\le 9$) & Oui & Oui & \textbf{Exacte} &
Inutile : l'exacte est calculable & Distribution exacte à privilégier \\

Mann--Kendall & LN & $Z$ / $S$ & Oui (loi mahonienne) & Oui & Oui & \textbf{Exacte} &
\textbf{Déconseillée} : approx. normale usuellement requise à $T\ge 10$ &
Distribution exacte à privilégier \\

Cox--Stuart & LN & $K$ & Oui (binomiale) & --- & Oui & \textbf{Exacte} &
Sans objet & Exacte ; mais $p_{\min}=0{,}125$ à $T=8$ : test inopérant \\

Breusch--Pagan studentisé (Koenker) & LN & LM & Non & Oui & Oui & Monte-Carlo &
\textbf{Déconseillée} : convergence $\chi^2$ lente, aucune borne à $T=8$ &
Bootstrap recommandé ; version de référence pour conclure \\

Breusch--Pagan original (1979) & LN & LM & Non & Oui, \emph{sous normalité} & Oui & Monte-Carlo &
\textbf{Déconseillée} : convergence lente \emph{et} hypothèse de normalité non acquise &
À lire en regard de la version studentisée ; un écart signale H3, pas H2 \\

White & LN & LM & Non & Oui & Oui & Monte-Carlo &
\textbf{Déconseillée} : 3 coefficients auxiliaires pour 8 points &
Bootstrap ; puissance quasi nulle \\

Goldfeld--Quandt & LN & $F$ & Approx. seulement & Oui & Oui & Monte-Carlo &
\textbf{Sous conditions} : exacte si régressions séparées, ici résidus globaux &
Bootstrap recommandé \\

Brown--Forsythe & LN & $W$ & Non & Oui & Oui & Monte-Carlo &
\textbf{Déconseillée} : loi $F$ approchée, 4 points par groupe &
Bootstrap recommandé \\

Smirnov 2 éch. & LN & $D$ & Oui (combinatoire) & Oui & Oui & \textbf{Exacte} &
Inutile & Exacte ; $n_1=n_2=4$ : $p_{\min}=0{,}029$ \\

Shapiro--Wilk & LN & $W$ & \textbf{Oui}, loi de $W$ évaluée par simulation directe (sans Royston)
& Approx. Royston & Oui & Monte-Carlo, avec p Royston \emph{et} p de loi nulle simulée &
\textbf{Sous conditions} pour Royston ; la loi nulle simulée s'en affranchit &
Loi nulle simulée à privilégier ; les deux concordent à $T=8$ \\

Shapiro--Francia & LN & $W'$ & Non & Approx. Royston & Oui & Monte-Carlo &
\textbf{Sous conditions} : validité annoncée $5\le T\le 5000$ &
Bootstrap recommandé \\

Anderson--Darling & LN & $A^2$ & Non (param. estimés) & Oui (Stephens, cas 3) & Oui &
Monte-Carlo, p de Stephens en complément &
\textbf{Sous conditions} : ajustement empirique, calibration vérifiée à $T=8$ &
Bootstrap retenu ; p de Stephens disponible \\

Cramér--von~Mises & LN & $W^2$ & Non & Oui (Stephens, cas 3) & Oui &
Monte-Carlo, p de Stephens en complément & Idem $A^2$ &
Bootstrap retenu ; p de Stephens disponible \\

KS contre $\mathcal{N}(0,1)$ & LN & $D$ & Non & Oui (Kolmogorov) & Oui & Monte-Carlo &
\textbf{Déconseillée} : 0 rejet sur 3\,000 simulations à 5 \%, test vidé de sa puissance &
Utiliser Lilliefors à la place \\

Lilliefors & LN & $D_{\mathrm{L}}$ & Non & Oui (Dallal--Wilkinson) & Oui &
p de Dallal--Wilkinson & \textbf{Oui} : calibration vérifiée à $T=8$ (5,2 \% à 5 \%) &
Version correcte du KS à paramètres estimés \\

Jarque--Bera & LN & JB & Non & Oui & Oui & Monte-Carlo &
\textbf{Déconseillée} : convergence $\chi^2_2$ très lente ; $b_2$ borné par
$\sim T$ & Bootstrap ; test peu informatif à $T=8$ \\

D'Agostino (asym.) & LN & $Z$ & Non & Oui (transf. $S_U$) & Oui & Monte-Carlo &
\textbf{Sous conditions} : défini dès $T\ge 8$, à la limite du domaine &
Bootstrap ; à interpréter avec prudence \\

Anscombe--Glynn & LN & $Z$ & Non & Oui & --- & \textbf{Non applicable} ($T<20$) &
Sans objet & Non utilisable à $T=8$ \\

Durbin--Watson & LN & DW & \textbf{Oui} : forme quadratique, méthode d'Imhof & --- & Oui &
\textbf{Exacte} (Imhof) & Sans objet : loi exacte &
Distribution exacte à privilégier ; statistique calculée sur résidus centrés \\

Ljung--Box (1, 2) & LN & $Q$ & Non & Oui & Oui & Monte-Carlo &
\textbf{Déconseillée} : approx. $\chi^2$ mal calibrée en petit échantillon &
Bootstrap recommandé \\

Box--Pierce & LN & $Q$ & Non & Oui & Oui & Monte-Carlo &
\textbf{Déconseillée} : encore moins bonne que Ljung--Box &
Bootstrap ; préférer Ljung--Box \\

Suites (W.--W.) & LN & $Z$ / $R$ & Oui (combinatoire) & Oui & Oui & \textbf{Exacte} &
\textbf{Déconseillée} : approx. normale requiert $n_1,n_2\ge 10$ &
Distribution exacte à privilégier \\

Centrage des résidus & LN & --- & --- & --- & --- &
\textbf{Diagnostic} & Sans objet & Aucune p-value~: $\bar z$ est rivé par une
identité de l'estimation. La p-value nominale du $t$ ne dépendrait pas des
données ($=1$ à $\hat\pi_t$ constant). Restitué sans verdict \\

Variance unitaire & LN & --- & --- & --- & --- &
\textbf{Diagnostic} & Sans objet & Idem pour $\operatorname{Var}(z)$~; la loi
simulée porte un atome en la valeur observée et la p-value nominale du $\chi^2$
ne dépendrait pas des données ($0{,}6651878$ à $T = 8$) \\

sup-F (rupture) & LN & $\sup F$ & Non & Théorique seulement (Andrews) & Oui & Monte-Carlo &
\textbf{Absence de justification} : loi de sup. de processus, convergence fonctionnelle &
\textbf{Seul test sans p-value non simulée} : bootstrap indispensable \\

OLS-CUSUM & LN & $\sup|S|$ & Non & Oui (Kolmogorov asympt.) & Oui & Monte-Carlo &
\textbf{Absence de justification} : 8 points ne forment pas une trajectoire
brownienne & Bootstrap indispensable ; p asympt. fournie pour mémoire \\

Grubbs & LN & $G$ & Borne de Bonferroni & --- & Oui & Monte-Carlo &
\textbf{Sous conditions} : borne conservatrice, non exacte &
Bootstrap ; borne utilisable comme majorant \\

Rosner (ESD) & LN & $R_i$ & Valeurs critiques tabulées & --- & --- &
Procédure de décision & Sans objet (pas de p-value) &
Indicatif seulement ; fiabilité annoncée pour $T\ge 25$ \\


\multicolumn{9}{@{}l@{}}{\textbf{\textcolor{financeblue}{Méthode Merz--Wüthrich --- M1 : proportionnalité des cumulés}}} \\

Tendance avec le cumulé, à colonne donnée & \textbf{MW} & $t$ pondéré (GLS,
effet fixe de colonne) & Non & Oui & Oui & Monte-Carlo &
\textbf{Déconseillée} : les cellules d'une colonne partagent $\hat f_j$ et
$\hat\sigma_j$ & Version agrégée puissante (100 \% sous l'alternative testée) ;
les tests par colonne localisent \\

Ordonnée à l'origine par colonne & \textbf{MW} & $X$ de Fisher & Non &
Oui, $\chi^2(2K)$ & Oui & Monte-Carlo &
\textbf{Déconseillée} : colonnes adjacentes non indépendantes &
\textbf{Test central de M1} ; calibration MC vérifiée (10,7 \% et 4,7 \%) \\

Homogénéité de $f_j$ entre années de survenance & \textbf{MW} & $X$ de Fisher &
Non & Oui, $\chi^2(2K)$ & Oui & Monte-Carlo & Idem &
Seconde exigence du texte : « pour toutes les années d'accident » \\

Absence de courbure & \textbf{MW} & $X$ de Fisher & Non & Oui, $\chi^2(2K)$ &
Oui & Monte-Carlo & Idem &
Puissance faible : 3 paramètres pour 4 à 6 observations \\

Famille $\alpha$ (pondération) & \textbf{MW} & amplitude relative & Non & Non &
Oui & Monte-Carlo & Sans objet : aucune loi analytique &
Diagnostic de robustesse ; chiffre l'enjeu du choix de pondération \\

\multicolumn{9}{@{}l@{}}{\textbf{\textcolor{financeblue}{Méthode Merz--Wüthrich --- M2 : variance proportionnelle au cumulé}}} \\

Breusch--Pagan sur résidus de Mack & \textbf{MW} & LM & Non & Oui & Oui &
Monte-Carlo & \textbf{Déconseillée} : convergence $\chi^2$ lente \emph{et}
dépendance intra-colonne & Bootstrap recommandé ; teste l'exposant 1 imposé par
le règlement \\

Exposant de variance par colonne & \textbf{MW} & $X$ de Fisher & Non &
Oui, $\chi^2(2K)$ & Oui & Monte-Carlo &
\textbf{Déconseillée} : colonnes non indépendantes &
Localise ce que Breusch--Pagan agrège ; un rejet est une limite du modèle, pas
un motif de changement de méthode \\

Variance unitaire des résidus de Mack & \textbf{MW} & --- & --- & --- & --- &
Diagnostic & Sans objet & Indicateur d'échelle de $\hat\sigma_j$ \\

\multicolumn{9}{@{}l@{}}{\textbf{\textcolor{financeblue}{Méthode Merz--Wüthrich --- M3 : indépendance}}} \\

Homogénéité des résidus entre années de survenance & \textbf{MW} & $H$ de
Kruskal--Wallis & Non & Oui, $\chi^2(k-1)$ & Oui & Monte-Carlo &
\textbf{Déconseillée} : groupes très déséquilibrés (de $J$ à 1 résidu) &
Traduction testable de D(2)(h)(i) ; complète le test calendaire \\

Effets d'année calendaire & \textbf{MW} & $Z$ & Non & Oui (moments \emph{exacts} de $Z$) & Oui &
Monte-Carlo & \textbf{Sous conditions} : les moments sont exacts, l'approximation
normale de leur somme ne l'est pas & Bootstrap de résidus ; test le plus
informatif sur un triangle \\

Corrélation entre années de développement & \textbf{MW} & $\rho$ agrégé & Non & Non &
Oui & Monte-Carlo & Sans objet : aucune loi analytique &
Bootstrap indispensable \\

Autocorrélation (Durbin--Watson) sur résidus de Mack & \textbf{MW} & DW & Non &
Non : loi d'Imhof inapplicable & Oui & Monte-Carlo &
Sans objet : aucune loi analytique & Bootstrap indispensable \\

Test des suites sur résidus de Mack & \textbf{MW} & $Z$ & Non (échangeabilité en
défaut) & Oui & Oui & Monte-Carlo &
\textbf{Déconseillée} : dépendance intra-colonne & Bootstrap recommandé \\

\multicolumn{9}{@{}l@{}}{\textbf{\textcolor{financeblue}{Méthode Merz--Wüthrich --- M4 : points aberrants}}} \\

Grubbs sur résidus de Mack & \textbf{MW} & $G$ & Borne de Bonferroni & --- & Oui &
Monte-Carlo & \textbf{Sous conditions} : suppose la normalité, non exigée par le
modèle & Bootstrap ; borne utilisable comme majorant \\

Rosner (ESD) sur résidus de Mack & \textbf{MW} & $R_i$ & Valeurs critiques
tabulées & --- & --- & Procédure de décision & Sans objet (pas de p-value) &
Indicatif ; $n = 28$ à $T = 8$, sous le seuil de fiabilité \\

\multicolumn{9}{@{}l@{}}{\textbf{\textcolor{financeblue}{Méthode Merz--Wüthrich --- M5 : normalité (diagnostic, non exigée)}}} \\

Shapiro--Wilk sur résidus de Mack & \textbf{MW} & $W$ & \textbf{Oui}, loi nulle
simulée & Approx. Royston & \textbf{Non} (voir texte) & Loi nulle simulée &
Sans objet : loi propre & Diagnostic seulement : normalité non exigée par D(2)(h) \\

Lilliefors sur résidus de Mack & \textbf{MW} & $D_{\mathrm{L}}$ & Non &
Oui (Dallal--Wilkinson) & Non & p de Dallal--Wilkinson &
\textbf{Sous conditions} & Variante secondaire ; diagnostic seulement \\

\multicolumn{9}{@{}l@{}}{\textbf{\textcolor{financeblue}{Méthode Merz--Wüthrich --- M6 : robustesse de l'estimation}}} \\

Concentration de la réserve sur la dernière année & \textbf{MW} & --- & --- & --- &
--- & Diagnostic & Sans objet & Équivalent du jackknife : mesure la concentration
du résultat sur la ligne la moins informative \\

Extrapolation de $\hat\sigma^2_{J-1}$ & \textbf{MW} & --- & --- & --- & --- &
Diagnostic & Sans objet & Valeur imposée par le règlement, non estimée :
à signaler dans le dossier \\
\end{longtable}
\end{center}
\end{landscape}

\noindent\emph{Lecture.} « Sans objet : loi exacte » signifie que la question de
la qualité de l'approximation asymptotique ne se pose pas, une loi exacte étant
employée. « Absence de justification suffisante » signifie que la littérature
consultée n'établit pas, pour $T = 8$, de borne d'erreur ni d'étude de
simulation permettant de conclure~: la p-value asymptotique n'est alors pas
retenue.

\newpage
% -----------------------------------------------------------------------------
\subsection{Tableau 2 : nature et vitesse des convergences}
\label{tab:conv}
% -----------------------------------------------------------------------------

Ce tableau sépare strictement les deux phénomènes~: la convergence en $T$ de la
loi de la statistique, et la convergence en $B$ de l'estimateur de la p-value.
Lorsque les deux coexistent, ils sont documentés sur deux lignes distinctes. Les
blocs (A), (B) et (D) concernent la méthode lognormale~; le bloc (C) est propre à
la méthode Merz--Wüthrich, où la profondeur pertinente n'est plus $T$ mais le
nombre d'observations $n_j = I - j$ de chaque colonne du triangle, qui décroît
jusqu'à 2 ou 3 sur les dernières années de développement.

\begin{landscape}
\begin{center}
\scriptsize
\setlength{\tabcolsep}{5pt}
\renewcommand{\arraystretch}{1.20}
\begin{longtable}{@{}p{2.7cm}p{3.4cm}p{2.6cm}p{2.5cm}p{3.5cm}p{3.0cm}p{3.6cm}@{}}
\toprule
\textbf{Test} & \textbf{Objet qui converge} & \textbf{Limite} &
\textbf{Nature} & \textbf{Vitesse} & \textbf{Hypothèses} &
\textbf{Conséquence à $T=8$} \\
\midrule
\endfirsthead
\toprule
\textbf{Test} & \textbf{Objet qui converge} & \textbf{Limite} &
\textbf{Nature} & \textbf{Vitesse} & \textbf{Hypothèses} &
\textbf{Conséquence à $T=8$} \\
\midrule
\endhead
\multicolumn{7}{@{}l@{}}{\textit{Suite page suivante}} \\
\endfoot
\bottomrule
\endlastfoot

\multicolumn{7}{@{}l@{}}{\textbf{(A) Convergence en $T$ de la loi de la
statistique --- non contrôlable, $T$ est une donnée}} \\
\midrule

Student / Fisher & \emph{Aucune} : loi exacte à $T$ fini & --- & --- & --- &
Erreurs i.i.d. $\mathcal{N}(0,\sigma^2)$, régresseurs fixes &
Valide, sous réserve de la normalité \\

Breusch--Pagan (Koenker), White & $\mathrm{LM} = TR^2_{\text{aux}}$ & $\chi^2(q)$ &
Convergence en loi & Non établie de façon exploitable ; TCL sous-jacent en
$\mathcal{O}(T^{-1/2})$ & Moments d'ordre 4 finis, indépendance &
Non exploitable ; bootstrap requis \\

Breusch--Pagan (1979) & $\mathrm{LM} = \tfrac12\,\mathrm{SCE}$ & $\chi^2(q)$ &
Convergence en loi & Idem, \emph{mais} la limite n'est atteinte que si
$\operatorname{Var}(u^2)=2\sigma^4$ & \textbf{Normalité des erreurs}, en sus
des précédentes & Non exploitable ; sur-rejet si queues lourdes \\

Ljung--Box, Box--Pierce & $Q$ & $\chi^2(h)$ & Convergence en loi &
Non établie ; niveau effectif connu comme dégradé en petit échantillon &
Stationnarité, moments d'ordre 2 & Non exploitable ; bootstrap requis \\

Jarque--Bera & $(\sqrt{b_1}, b_2)$ puis JB & $\chi^2(2)$ &
Convergence en loi & Réputée \emph{très} lente & Moments d'ordre 8 finis
pour la normalité asymptotique de $b_2$ & Non exploitable ; $b_2$ borné
par $\sim T$ \\

Mann--Kendall & $S$ ($U$-statistique) & $\mathcal{N}(0,1)$ après
normalisation & Convergence en loi & Berry--Esseen pour $U$-statistiques :
$\mathcal{O}(T^{-1/2})$ & Variance de la projection non nulle, moment d'ordre 3 &
Borne $\mathcal{O}(8^{-1/2})\approx 0{,}35$ : inexploitable. Loi exacte employée \\

Wald--Wolfowitz & $R$ & $\mathcal{N}(0,1)$ & Convergence en loi &
$\mathcal{O}(n^{-1/2})$ (TCL) ; recommandation d'usage $n_1,n_2\ge 10$ &
Continuité, indépendance & Loi exacte employée \\

Spearman & $\hat\rho_s$ & $\mathcal{N}$ / $t(T-2)$ & Convergence en loi &
$\mathcal{O}(T^{-1/2})$ & Indépendance, absence d'ex \ae quo &
Loi de permutation exacte employée \\

D'Agostino & $\sqrt{b_1}$ transformé & $\mathcal{N}(0,1)$ &
Convergence en loi & Transformation $S_U$ conçue pour corriger l'ordre
$T^{-1}$ ; validité annoncée $T\ge 8$ & Moments d'ordre 6 finis &
Limite du domaine ; bootstrap en complément \\

Anderson--Darling, CvM, KS & Processus empirique
$\sqrt{T}(\hat F - F_0)$ & Pont brownien (transformé) &
Convergence faible dans $D[0,1]$ & $\mathcal{O}(T^{-1/2})$ pour les
fonctionnelles régulières & $F_0$ continue, paramètres \emph{connus} &
Hypothèse violée (paramètres estimés) ; bootstrap requis \\

sup-F & $\{F(k)\}_k$ & Sup. d'un processus de Bessel &
Convergence faible & Fonctionnelle : lente, non quantifiée à $T$ petit &
Erreurs i.i.d., écrêtement $\tau$ fixé & Non exploitable ; bootstrap requis \\

OLS-CUSUM & $S(\lfloor uT\rfloor)$ & Pont brownien $B^\circ$ &
Convergence faible (principe d'invariance) & $\mathcal{O}(T^{-1/2})$ &
Erreurs i.i.d., variance finie & 8 points : non exploitable \\

Shapiro--Wilk / Francia & $W$, $W'$ & Approx. normalisée de Royston &
Approximation numérique, non asymptotique & Sans objet (ajustement empirique) &
Domaine de validité annoncé par l'auteur & Domaine couvert mais
transformation distincte à $T\le 11$ ; bootstrap en complément \\

Centrage, variance unitaire (diagnostics) & \emph{Aucun} : $\bar z$ et
$\operatorname{Var}(z)$ sont rivés par les conditions du premier ordre & --- &
--- & Sans objet & Optimum intérieur en $\gamma$ ; $\hat\pi_t$ constant pour la
forme simple $\sum_t z_t = 0$, $\sum_t z_t^2 = T$ &
Aucune p-value : la loi simulée sous le modèle ajusté porte un atome en la
valeur observée, et la loi nominale, qui ignore l'estimation, produit à
$\hat\pi_t$ constant une p-value qui ne dépend plus des données du tout : $1$
pour le $t$ de Student et $0{,}6651878$ pour le $\chi^2$, à $T = 8$. Restitués
sans verdict \\

\midrule
\multicolumn{7}{@{}l@{}}{\textbf{(B) Convergence en $B$ de l'estimateur
Monte-Carlo --- entièrement contrôlable}} \\
\midrule

Toutes statistiques bootstrapées & $\hat p_{\text{MC}} =
\frac{1+\#\{S^{(b)}\ge S_{\text{obs}}\}}{B+1}$ & $p^{*}$, p-value bootstrap
exacte & Convergence presque sûre (LFGN) et en loi (TCL) &
Écart-type $\sqrt{p(1-p)/B} = \mathcal{O}(B^{-1/2})$ ; granularité $1/(B+1)$ &
Réplications i.i.d. sous le modèle ajusté &
$B=999 \Rightarrow \pm 0{,}009$ à $p=0{,}10$. \textbf{Sans lien avec $T$} \\

Idem & Loi bootstrap $p^{*}$ & Vraie p-value sous $H_0$ &
Convergence bootstrap en $T$ & Consistance établie asymptotiquement ;
raffinement d'ordre supérieur pour statistiques pivotales &
Modèle correctement spécifié, estimateur consistant &
\textbf{Non quantifiable à $T=8$}. C'est ici, et non dans le choix de $B$, que
réside la limite du dispositif~: la correction en $1/(B+1)$ ne rend pas le test exact \\

\midrule
\multicolumn{7}{@{}l@{}}{\textbf{(C) Convergence propre à la méthode Merz--Wüthrich}} \\
\midrule

Effets d'année calendaire & $Z = \sum_k \min(L_k,S_k)$ & $\mathcal{N}(0,1)$
après centrage-réduction & Convergence en loi &
Non quantifiée ; les moments de chaque $Z_k$ sont en revanche \emph{exacts} &
Indépendance des diagonales, absence d'ex \ae quo dans les colonnes &
Le nombre de diagonales exploitables est $I-1$ : la somme porte sur peu de
termes, d'où le recours au bootstrap \\

Facteurs $\hat f_j$ & Estimateur des moindres carrés pondérés &
$f_j$ & Convergence en probabilité ; sans biais sous M1 &
$\mathcal{O}_P(n_j^{-1/2})$ avec $n_j = I-j$ observations &
M1 et M3 & Les colonnes de droite ne reposent que sur 2 ou 3 observations \\

Variances $\hat\sigma_j^2$ & Estimateur pondéré de la variance &
$\sigma_j^2$ & Convergence en probabilité &
$\mathcal{O}_P(n_j^{-1/2})$, avec $n_j - 1$ degrés de liberté &
M2, moments d'ordre 4 finis & $\hat\sigma^2_{J-1}$ n'est pas estimée mais
\textbf{extrapolée} par la règle du règlement \\

$\sqrt{\mathrm{MSEP}}$ & Erreur de prédiction à un an & Erreur vraie &
Convergence en probabilité & Non établie à $I$ fini &
M1 à M3, modèle correctement spécifié &
Non quantifiable ; la MSEP hérite de l'incertitude sur $\hat\sigma^2_{J-1}$ \\

Combinaisons de Fisher (ordonnée à l'origine, homogénéité, courbure,
exposant de variance) & $X = -2\sum\ln p_k$ & $\chi^2(2K)$ &
Convergence en loi & \textbf{Sans objet en toute rigueur} : la loi $\chi^2(2K)$
suppose l'indépendance des $K$ p-values, mise en défaut car les colonnes $j$ et
$j+1$ partagent $C(\cdot,j+1)$ & Indépendance des colonnes &
p-value du $\chi^2$ indicative seulement ; bootstrap retenu \\

Corrélation entre colonnes & $\hat\rho$ agrégé & \emph{aucune limite tabulable} &
--- & Sans objet : la géométrie du triangle interdit toute normalisation &
Continuité, absence d'ex \ae quo & Bootstrap seul recours ; 5 à 6 paires
exploitables à $T = 8$ \\

Durbin--Watson et suites sur résidus de Mack & DW, $Z$ &
$\mathcal{N}(0,1)$ / loi combinatoire & Convergence en loi &
Non établie : les résidus d'une colonne partagent $\hat f_j$ et $\hat\sigma_j$ &
Échangeabilité, mise en défaut ici & Lois classiques écartées ; bootstrap retenu \\

Bootstrap de résidus & $\hat p_{\mathrm{MC}}$ & p-value semi-paramétrique &
Convergence en $B$ (Monte-Carlo) et en $I$ (validité bootstrap) &
$\mathcal{O}(B^{-1/2})$ pour la première ; non établie pour la seconde &
Échangeabilité des résidus standardisés &
\textbf{Ne permet pas de tester la normalité} : la loi nulle simulée est la loi
empirique des résidus \\

\midrule
\multicolumn{7}{@{}l@{}}{\textbf{(D) Convergence de l'estimateur du paramètre (méthode lognormale)}} \\
\midrule

$(\hat\delta,\hat\gamma)$ (MV) & Estimateur du maximum de vraisemblance &
$(\delta_0,\gamma_0)$ & Convergence en probabilité ; normalité asymptotique &
$\mathcal{O}_P(T^{-1/2})$ sous conditions de régularité &
Modèle correct, identifiabilité, intérieur du domaine &
$\hat\delta$ souvent \textbf{au bord} : régularité en défaut, normalité
asymptotique non applicable \\

Test du rapport de vraisemblance sur $\delta$ & $\mathrm{LR}$ &
$\tfrac12\chi^2(0)+\tfrac12\chi^2(1)$ & Convergence en loi &
Non quantifiée à $T$ petit & Frontière du domaine (résultat de Chernoff) &
p-value $\chi^2(1)$ conservatrice d'un facteur 2 \\

\end{longtable}
\end{center}
\end{landscape}

\medskip
\begin{encadre}
\textbf{Point de méthode central.} Les deux blocs (A) et (B) répondent à des
questions différentes. Le bloc (B) montre que l'erreur numérique sur
$\hat p_{\text{MC}}$ est maîtrisée et réductible à volonté en augmentant $B$. Le
bloc (A) montre qu'à $T = 8$, aucune des lois asymptotiques mobilisées ne
dispose d'une justification exploitable. Le bootstrap paramétrique déplace le
problème sans le supprimer~: il remplace une approximation asymptotique par une
approximation \emph{de plug-in}, dont la validité reste asymptotique en $T$ et
n'est pas quantifiable à $T = 8$. C'est la raison pour laquelle la
section~\ref{sec:robustesse} (jackknife, largeur de l'intervalle bootstrap,
position de $\hat\delta$) doit primer sur la lecture des p-values dans le
dossier remis au superviseur.
\end{encadre}

\subsection{Choix de la base de résidus}
\label{sec:base-residus}

Plusieurs séries peuvent servir de support aux tests. Le choix n'est pas neutre~:
il détermine à la fois l'hypothèse effectivement testée et la validité des
p-values classiques.

\subsubsection*{Les trois séries candidates}

\begin{center}
\small
\rowcolors{2}{tablerow}{white}
\begin{tabular}{@{}p{2.6cm}p{4.6cm}p{7.6cm}@{}}
\rowcolor{tableheader}
\textcolor{financeblue}{\textbf{Série}} &
\textcolor{financeblue}{\textbf{Définition}} &
\textcolor{financeblue}{\textbf{Loi sous le modèle}} \\
Résidus normalisés & $z_t = \sqrt{\hat\pi_t}\,v_t$ &
i.i.d. $\mathcal{N}(0,1)$ : la seule série exactement échangeable \\
Résidus logarithmiques & $v_t = \ln r_t + \tfrac{1}{2\hat\pi_t} - \ln\hat\beta$ &
indépendants et normaux, mais de variance $1/\pi_t$ variable avec $x_t$ \\
Ratios centrés & $u_t = r_t - \bar r$, avec $r_t = y_t/x_t$ &
indépendants, mais ni normaux ni de même variance \\
\end{tabular}
\end{center}

\noindent
La différence $y_t - x_t$ n'est pas un résidu mais le \emph{résultat technique}
de l'exercice~: elle mélange le niveau de marge et l'erreur de prévision. Le
résidu correspondant sur l'échelle monétaire est $e_t = y_t - \hat\beta x_t$,
égal à $x_t\,(r_t - \hat\beta)$~: à un facteur de volume près, il coïncide avec
$u_t$. C'est donc $u_t$ qui est retenu comme base « brute », le passage au ratio
neutralisant l'effet de taille au premier ordre.

\subsubsection*{Amplitude effective de l'hétéroscédasticité résiduelle}

Sous le modèle, $\operatorname{Var}(r_t) = \sigma^2\big[\delta +
(1-\delta)\bar x/x_t\big]$. Sur un jeu de volumes allant de 50 à 200 avec
$\delta = 0{,}35$, l'écart-type théorique des $r_t$ varie de $0{,}120$ à
$0{,}076$, soit un facteur $1{,}6$ entre la première et la dernière année. Les
$u_t$ ne sont donc pas identiquement distribués~: toute loi de référence
supposant l'échangeabilité (permutation, combinatoire, ou normale i.i.d.) est
invalide sur cette base. Cette invalidité disparaît lorsque $\hat\pi_t$ est
constant --- c'est-à-dire lorsque $\hat\delta = 1$, \emph{ou} lorsque les
volumes $x_t$ le sont --- puisque les $r_t$ sont alors identiquement
distribués. Les trois séries ne deviennent pas pour autant proportionnelles~:
$z_t = \sqrt{\hat\pi}\,v_t$ l'est bien de $v_t$, à facteur constant, mais $u_t$
vit sur l'échelle linéaire quand $v_t$ vit sur l'échelle logarithmique, et
l'on ne passe de l'une à l'autre que par une transformation monotone non
linéaire. Ce point est repris plus bas~: à $\hat\pi_t$ constant, cette
transformation est tout ce qui sépare encore les deux bases, ce qui prive la
base $u_t$ de son objet.

\subsubsection*{Où la bascule est légitime, et où elle ne l'est pas}

\begin{center}
\small
\rowcolors{2}{tablerow}{white}
\begin{tabular}{@{}p{2.5cm}p{1.9cm}p{10.4cm}@{}}
\rowcolor{tableheader}
\textcolor{financeblue}{\textbf{Famille}} &
\textcolor{financeblue}{\textbf{Bascule}} &
\textcolor{financeblue}{\textbf{Justification}} \\
H1 --- linéarité & sans objet &
Les tests de tendance (Spearman, Mann--Kendall, Cox--Stuart) portent \emph{déjà}
sur $r_t$. Il n'existe pas de version sur $z_t$ qui aurait un sens~: la tendance
recherchée est celle du ratio lui-même. \\
H2 --- variance & \textbf{non} &
Changer de base changerait l'hypothèse nulle. Sur $z_t$, on teste s'il subsiste
une hétéroscédasticité \emph{après} pondération, ce qui doit être rejeté par les
données si le modèle est correct. Sur $u_t$, on testerait l'hétéroscédasticité
\emph{brute}, dont la présence est précisément ce que le modèle postule~: le
test serait attendu significatif et ne dirait rien de la qualité du calibrage. \\
H3 --- lognormalité & \textbf{non} &
Le modèle prédit que $z_t$ est normal, non que $r_t$ l'est ($r_t$ est
approximativement lognormal). Tester la normalité de $u_t$ reviendrait à tester
une hypothèse que le modèle ne formule pas. Sur $v_t$, la normalité est bien
prédite mais la variance $1/\hat\pi_t$ n'est pas constante, ce qui invalide les
tests de normalité usuels dès que $\hat\pi_t$ varie, c'est-à-dire dès que
$\hat\delta < 1$ \emph{et} que les volumes $x_t$ ne sont pas constants. \\
H4 --- indépendance & \textbf{oui} &
L'indépendance est une propriété que le modèle attribue aux $Y_t$
\emph{conditionnellement} aux $x_t$~: elle se transmet aux trois séries. Les
deux bases testent donc la même hypothèse. \\
Stabilité et valeurs aberrantes & \textbf{oui} &
Idem. Sur $u_t$, une valeur aberrante s'interprète directement comme un
boni/mali exceptionnel et une rupture comme un changement de régime du ratio,
sans dépendre de la qualité de l'ajustement. \\
\end{tabular}
\end{center}

\subsubsection*{Statut des p-values selon la base}

Sur la base $u_t$, les p-values exactes et asymptotiques sont \textbf{indicatives
seulement}~: elles supposent l'échangeabilité, mise en défaut par
l'hétéroscédasticité résiduelle. En revanche, \textbf{la p-value de Monte-Carlo
reste exactement valide}~: le bootstrap paramétrique simule $Y_t$ sous le modèle
ajusté, donc la loi de référence d'une statistique quelconque des données --- y
compris calculée sur $u_t$ --- est correctement reproduite, hétéroscédasticité
comprise. C'est cette p-value que le moteur retient sur cette base.

Vérification par simulation à $T = 8$, sous $H_0$, statistiques calculées sur
$u_t$ et p-values de Monte-Carlo (400 réplications, $B = 199$)~:

\begin{center}
\small
\rowcolors{2}{tablerow}{white}
\begin{tabular}{@{}p{5.0cm}cc@{}}
\rowcolor{tableheader}
\textcolor{financeblue}{\textbf{Statistique sur $u_t$}} &
\textcolor{financeblue}{\textbf{rejets à 10 \%}} &
\textcolor{financeblue}{\textbf{rejets à 5 \%}} \\
Durbin--Watson & 0,107 & 0,050 \\
Grubbs         & 0,115 & 0,055 \\
\end{tabular}
\end{center}

\noindent
Les deux valeurs sont compatibles avec les niveaux nominaux à l'intérieur de
l'intervalle de confiance à 95~\% ($\pm 0{,}029$ et $\pm 0{,}021$).

\subsubsection*{Ce contrôle est sans objet lorsque $\hat\pi_t$ est constant}

Six lignes de la table auditable sont calculées sur la base $u_t$~:
Durbin--Watson, Ljung--Box au retard~1, test des suites, rupture de niveau
(sup-F), stabilité cumulée (OLS-CUSUM) et Grubbs, chacune ayant son homologue
sur $z_t$. Le contrôle qu'elles portent consiste à vérifier que les conclusions
d'indépendance et de stabilité ne sont pas un artefact de la pondération
$\hat\pi_t$. \textbf{Là où il n'y a pas de pondération, il n'y a pas d'artefact
de pondération à détecter}, et ce contrôle perd son objet. C'est le cas dès que
$\hat\pi_t$ est constant, donc à $\hat\delta = 1$ ou à volumes $x_t$ constants
--- soit le cas courant à $T = 8$.

\emph{Conséquence algébrique exacte.} Lorsque $\hat\pi_t \equiv \hat\pi$, la
forme fermée de $\ln\hat\beta$ se réduit à
$\ln\hat\beta = \tfrac{1}{2\hat\pi} + \overline{\ln r}$, d'où
$v_t = \ln r_t - \overline{\ln r}$ et
\[
  z_t \;=\; \sqrt{\hat\pi}\,\big(\ln r_t - \overline{\ln r}\big).
\]
Les deux bases ne diffèrent alors plus que par la transformation logarithmique
$r \mapsto \ln r$, à un facteur d'échelle et un recentrage près. Cette
transformation est \emph{monotone}.

\emph{Conséquence exacte pour le test des suites.} La statistique des suites ne
dépend que des signes des écarts à la médiane, qu'une transformation monotone
préserve~: elle prend donc la \textbf{même valeur} sur les deux bases. Mesuré
sur le jeu de contrôle ($T = 8$, $\hat\delta = 1$)~: $Z = -0{,}7637626$ sur
$z_t$ comme sur $u_t$, et les huit signes de $r_t - \operatorname{med}(r)$ et de
$z_t - \operatorname{med}(z)$ coïncident un à un. Les deux lignes affichent
néanmoins des p-values différentes ($1{,}000$ contre $0{,}768$), non parce que
la statistique diffère, mais parce que la ligne sur $z_t$ retient la p-value
exacte combinatoire alors que celle sur $u_t$ retient une p-value de
Monte-Carlo.

\emph{Pour les cinq autres, un écart du second ordre.} En écrivant
$r_t = \bar r\,(1 + \varepsilon_t)$, le développement de $\ln$ donne
$\ln r_t - \overline{\ln r} = \varepsilon_t + \mathcal{O}(\varepsilon^2)$~: les
deux séries centrées coïncident au premier ordre, et leur différence est
d'ordre $\varepsilon^2$, soit du \emph{second} ordre en coefficient de
variation des $r_t$ (approximation numérique, développement limité). L'écart
tient donc à l'opposition \emph{logarithmique contre linéaire}, et non à
$\hat\pi_t$. Mesuré sur le jeu de contrôle, où le coefficient de variation vaut
$\mathrm{CV}(r) = 0{,}151$~: l'écart maximal entre $\ln r_t - \overline{\ln r}$
et $(r_t - \bar r)/\bar r$ vaut $0{,}0143$, à comparer à
$\mathrm{CV}(r)^2 = 0{,}0228$~;
$\operatorname{cor}(z, u) = 0{,}998$ et $\operatorname{cor}(z, \ln r) = 1$ à la
précision machine. Les six lignes sur base $u_t$ et leurs six homologues sur
$z_t$ y portent le même verdict \code{OK}~; pour les cinq paires autres que le
test des suites, les p-values retenues diffèrent d'au plus $0{,}091$ (sup-F,
$0{,}840$ contre $0{,}931$), l'écart le plus faible étant celui de
Durbin--Watson ($0{,}0011$). Une statistique qui
repose sur une quasi-annulation garde en revanche un écart relatif important
sur la valeur elle-même --- Ljung--Box au retard~1 vaut $9{,}2\cdot 10^{-3}$
sur $z_t$ et $1{,}3\cdot 10^{-4}$ sur $u_t$ --- sans que le verdict en dépende
à ce niveau de p-value.

\emph{Ce que cela change pour le dossier, et ce que cela ne change pas.} Les
six lignes restent exactes~: leur statistique, leur p-value de Monte-Carlo et
leur verdict sont justes, et elles gardent tout leur sens \emph{lorsque
$\hat\pi_t$ varie}, cas pour lequel elles ont été prévues. Ce qui était faux
est la promesse formulée sans condition~: présentées comme six vérifications
indépendantes, elles sont, à $\hat\pi_t$ constant, des \textbf{quasi-doublons}
de leurs homologues sur $z_t$ --- et un doublon exact pour le test des suites.
Six verdicts \code{OK} de la table auditable n'y apportent donc pas six
informations. Le moteur porte lui-même cette réserve~: dans ce cas, le champ
\code{detail} de chacune des six lignes commence par « \code{CONTROLE SANS
OBJET ICI} » et indique que la ligne est un quasi-doublon de son homologue sur
résidus standardisés. La situation n'a rien d'exceptionnel~: sur les $999$
réplications du bootstrap paramétrique du jeu de contrôle, $493$ --- soit
$49{,}3~\%$ --- ont un $\hat\pi_t^{*}$ constant (constat de simulation, graine
20260831).

\medskip
\begin{encadre}
\textbf{Usage recommandé.} La base $z_t$ reste celle du dossier~: c'est la seule
sur laquelle les p-values exactes et asymptotiques sont valides, et c'est elle
qui teste les hypothèses telles que l'annexe~XVII les formule. La base $u_t$ est
un \emph{contrôle de cohérence de la pondération}, \textbf{sous la condition
que cette pondération existe}~: elle vérifie que les conclusions
d'indépendance et de stabilité ne sont pas un artefact des poids $\hat\pi_t$.
Une divergence marquée entre les deux bases sur un même test
signale que le résultat dépend du modèle plutôt que des données, et doit être
documentée. \textbf{Lorsque $\hat\pi_t$ est constant, ce contrôle est sans
objet} et les six lignes correspondantes sont des quasi-doublons de leurs
homologues sur $z_t$~: elles ne doivent pas être lues comme six vérifications
supplémentaires (voir « Ce contrôle est sans objet lorsque $\hat\pi_t$ est
constant », plus haut dans cette sous-section, et le champ \code{detail} de ces
lignes, qui porte la réserve). Les deux bases sont calculées systématiquement par le moteur et
figurent toutes deux dans l'export~; la bascule de l'application ne fait que
filtrer l'affichage.
\end{encadre}

\subsection{Variantes principales et secondaires}
\label{sec:variantes}

Quatre tests existent dans le dispositif sous deux formes concurrentes. Pour
alléger la lecture, l'application n'affiche par défaut que la forme retenue
comme référence, les autres restant accessibles par une case à cocher et
présentes dans l'export.

\begin{center}
\small
\rowcolors{2}{tablerow}{white}
\begin{tabular}{@{}p{3.6cm}p{3.6cm}p{7.6cm}@{}}
\rowcolor{tableheader}
\textcolor{financeblue}{\textbf{Forme de référence}} &
\textcolor{financeblue}{\textbf{Forme secondaire}} &
\textcolor{financeblue}{\textbf{Motif du choix}} \\
Breusch--Pagan studentisé (Koenker) & Breusch--Pagan original (1979) &
La forme originale suppose la normalité des erreurs, hypothèse non acquise
à $T = 8$~; elle sur-rejette fortement sur queues lourdes. \\
Shapiro--Wilk (Royston) & Shapiro--Wilk à loi nulle simulée &
Les deux concordent à $T = 8$ (4,95~\% contre 5,05~\%). La forme de Royston est
retenue comme référence parce qu'elle est celle de \code{shapiro.test()}, donc
la plus immédiatement vérifiable. \\
Lilliefors & Kolmogorov--Smirnov contre $\mathcal{N}(0,1)$ &
La seconde applique une loi de référence valable à paramètres connus~; à
paramètres estimés elle est vidée de sa puissance. \\
Ljung--Box & Box--Pierce &
La pondération $(T+2)/(T-k)$ de Ljung--Box corrige le biais de petit
échantillon que Box--Pierce ignore. \\
\end{tabular}
\end{center}

\noindent
Le classement en forme de référence et forme secondaire est une convention
d'\emph{affichage}, consignée dans le champ \code{variante} de la table des
tests. Il ne modifie ni les calculs, ni les verdicts, ni le contenu de l'export.

\subsection{Tests dont l'usage reste fragile à $T = 8$}

Les tests suivants figurent dans le dispositif mais leur interprétation doit
être explicitement réservée à $T = 8$. Chacun est accompagné d'un complément
méthodologique, indiqué en troisième colonne.

\begin{center}
\small
\rowcolors{2}{tablerow}{white}
\begin{tabular}{@{}p{3.4cm}p{5.4cm}p{5.9cm}@{}}
\rowcolor{tableheader}
\textcolor{financeblue}{\textbf{Test}} &
\textcolor{financeblue}{\textbf{Fragilité à $T = 8$}} &
\textcolor{financeblue}{\textbf{Complément méthodologique}} \\
\rowcolor{tableheader}
\multicolumn{3}{@{}l@{}}{\textcolor{financeblue}{\textbf{Méthode lognormale}}} \\
Cox--Stuart & $m = 4$ paires $\Rightarrow$
$p_{\min} = 2\cdot 2^{-4} = 0{,}125$~: le rejet est \emph{impossible} à tout
seuil usuel & Conservé (loi exacte), mais signalé comme inopérant~; la détection
de tendance repose sur Mann--Kendall exact \\
Jarque--Bera & $b_2$ borné par $\sim T$~; convergence $\chi^2_2$ très lente &
p-value de Monte-Carlo ajoutée~; test à lire comme purement indicatif \\
White & 3 coefficients auxiliaires estimés sur 8 points & p-value de
Monte-Carlo ajoutée~; Breusch--Pagan à privilégier \\
Anscombe--Glynn & Non défini pour $T < 20$ & Marqué \emph{non applicable}~;
l'analyse de queue repose sur $A^2$ et le QQ-plot à enveloppe \\
Rosner (ESD) & Fiabilité des valeurs critiques annoncée pour $T \ge 25$ &
Conservé comme signal d'attention~; aucune exclusion de donnée ne doit en
découler sans justification métier \\
sup-F, OLS-CUSUM & Convergence fonctionnelle~; 8 points ne forment pas une
trajectoire limite & p-value de Monte-Carlo seule retenue~; à lire
conjointement avec la distance de Cook \\
Student / Fisher (pente) & Loi exacte, mais \emph{conditionnelle} à la
normalité des erreurs, non testable indépendamment à $T = 8$ &
Aucune p-value de Monte-Carlo possible (voir texte)~; la normalité est
examinée séparément en section~6 \\
\rowcolor{tableheader}
\multicolumn{3}{@{}l@{}}{\textcolor{financeblue}{\textbf{Méthode Merz--Wüthrich}}} \\
Corrélation entre années de développement & Cinq à six paires de colonnes
seulement comptent au moins trois observations à $T = 8$ &
Aucune loi analytique n'existe~; bootstrap seul recours, puissance très faible \\
Rosner (ESD) sur résidus de Mack & Fiabilité annoncée pour $n \ge 25$~; un
triangle $8\times8$ fournit 28 résidus, soit la limite du domaine &
Conservé comme signal d'attention~; aucune exclusion de cellule sur ce fondement \\
Extrapolation de $\hat\sigma^2_{J-1}$ & La quantité n'est pas estimée mais
imposée par une règle du règlement, appuyée sur deux colonnes &
Diagnostic dédié~; à mentionner explicitement dans le dossier \\
Tests de normalité (M5) & Non exigés par le modèle, et exclus du bootstrap de
résidus dont la nulle est la loi empirique observée &
Loi nulle propre pour Shapiro--Wilk~; lecture comme diagnostic seulement \\
\end{tabular}
\end{center}


\newpage

\partie{Partie III}{Fiches de tests par hypothèse}

% =============================================================================
\section{Hypothèse H1 --- linéarité et proportionnalité}
\label{sec:h1}
% =============================================================================

\begin{encadre}
\textbf{Énoncé réglementaire} (annexe~XVII, B/C(2)(f)(i)). L'espérance des pertes agrégées est
\emph{linéairement proportionnelle} aux primes acquises (ou à la provision d'ouverture)~:
$\mathbb{E}[Y_t] = \beta x_t$, sans terme constant, avec $\beta$ constant dans le temps.
\end{encadre}

\begin{fiche}{Test de Student sur la constante}
\refl{W.\,S. Gosset (« Student »), \emph{The probable error of a mean}, \emph{Biometrika}, 6(1),
1908, p.~1--25.}
\Cadre
Régression auxiliaire $y_t = a + b x_t + \varepsilon_t$. Hypothèses non testées~: les
$\varepsilon_t$ sont indépendants, de variance constante $\sigma_\varepsilon^2$, de loi
$\mathcal{N}(0,\sigma_\varepsilon^2)$, et les $x_t$ sont exogènes et non aléatoires. Nécessite
$T \ge 3$ et $S_{xx} > 0$.
\Hzero
$H_0 : a = 0$, absence de composante de sinistralité indépendante du volume. C'est la traduction
testable de la \emph{proportionnalité stricte} exigée par le règlement.
\Stat
\[
  t_a \;=\; \frac{\hat a}{\widehat{\mathrm{se}}(\hat a)},
  \qquad
  \widehat{\mathrm{se}}(\hat a) = \hat s \sqrt{\frac{1}{T} + \frac{\bar x^{\,2}}{S_{xx}}},
  \qquad \hat s^2 = \frac{\mathrm{SCR}}{T-2}.
\]
\Loi
$t_a \sim t(T-2)$, loi exacte sous normalité des erreurs.
\Pval
$p = 2\big(1 - F_{t,T-2}(|t_a|)\big)$. Cette p-value est \textbf{exacte} sous normalité des
erreurs~: elle est donc retenue en priorité par le moteur, la p-value de Monte-Carlo n'étant
calculée qu'à titre de recoupement (les deux concordent étroitement sur le jeu de référence).
\Usage
On souhaite \emph{ne pas} rejeter $H_0$, ce qui constitue la faiblesse logique de ce test~:
\textbf{l'absence de preuve n'est pas une preuve d'absence}. Un non-rejet ne démontre pas la
proportionnalité, alors que c'est précisément ce que le règlement demande d'établir. Limite
aggravante dans le contexte USP~: $x_t$ est un volume toujours très éloigné de zéro et de faible
amplitude relative. La constante est donc extrapolée loin du domaine des données et
$\widehat{\mathrm{se}}(\hat a)$ est mécaniquement très grande (le terme $\bar x^2/S_{xx}$ domine),
si bien que le test ne rejette pratiquement jamais. Le test d'équivalence de la fiche suivante
renverse cette charge de la preuve.
\end{fiche}

\begin{fiche}{Équivalence de la constante à zéro (TOST)}
\refl{D.\,J. Schuirmann, \emph{A comparison of the two one-sided tests procedure and the power
approach for assessing the equivalence of average bioavailability}, \emph{Journal of
Pharmacokinetics and Biopharmaceutics}, 15(6), 1987, p.~657--680 ; W.\,J. Westlake,
\emph{Response to T.\,B.\,L. Kirkwood: bioequivalence testing --- a need to rethink},
\emph{Biometrics}, 37(3), 1981, p.~589--594.}
\Cadre
Même régression auxiliaire et mêmes hypothèses de validité que la fiche précédente. On se donne
en outre une \emph{marge d'équivalence} $\Delta > 0$, qui n'est pas une quantité statistique mais
un seuil économique~: en deçà de $\Delta$, une composante fixe de sinistralité est jugée sans
portée pratique. Le moteur propose deux modes~: $\Delta$ fixée a priori
(\code{delta\_equiv}), ou $\Delta = \theta\,\bar y$ (\code{theta\_equiv}, 10~\% par défaut),
c'est-à-dire une fraction de la perte annuelle moyenne.
\Hzero
\[
  H_0 : |a| \ge \Delta \quad\text{(la constante n'est \emph{pas} négligeable)}
  \qquad\text{contre}\qquad
  H_1 : |a| < \Delta .
\]
\textbf{Les rôles sont inversés par rapport au test précédent}~: c'est le \emph{rejet} de $H_0$
qui établit la proportionnalité pratique. L'hypothèse à démontrer est passée en $H_1$, ce qui est
la construction correcte lorsque l'on veut prouver une absence d'effet.
\Stat
Décomposition en deux tests unilatéraux (procédure d'intersection-union)~:
\[
  t_{\text{bas}} = \frac{\hat a + \Delta}{\widehat{\mathrm{se}}(\hat a)}
  \quad (H_0^{-} : a \le -\Delta),
  \qquad
  t_{\text{haut}} = \frac{\hat a - \Delta}{\widehat{\mathrm{se}}(\hat a)}
  \quad (H_0^{+} : a \ge +\Delta).
\]
On rejette $H_0$ au niveau $\alpha$ si et seulement si les deux sous-hypothèses sont rejetées,
soit $t_{\text{bas}} > t_{1-\alpha,\,T-2}$ \emph{et} $t_{\text{haut}} < -t_{1-\alpha,\,T-2}$.
\Loi
Chaque statistique unilatérale suit \emph{exactement} une loi de Student à $T-2$ degrés de
liberté sous normalité des erreurs. Aucun résultat asymptotique n'intervient.
\Pval
Règle du maximum, propre aux procédures d'intersection-union~:
\[
  p_{\text{TOST}} = \max\Big(1 - F_{t,T-2}(t_{\text{bas}}),\ F_{t,T-2}(t_{\text{haut}})\Big).
\]
Cette p-value est exacte \textbf{si $\Delta$ est fixée a priori}. Avec la marge relative
$\Delta = \theta\bar y$, la marge dépend des données et l'exactitude n'est plus qu'approchée~; le
moteur signale alors la nature de la p-value comme \emph{quasi-exacte} et non comme exacte. Pour
un dossier ACPR, il est recommandé d'arrêter $\Delta$ à l'avance et de le documenter.
\Usage
C'est le seul test du dispositif qui puisse fournir une \emph{preuve positive} de la
proportionnalité exigée par l'annexe~XVII. Sa limite est celle de tous les tests à $T = 8$~: la
puissance est faible. Une étude de simulation sous $a = 0$ avec $\theta = 10\,\%$ donne une
probabilité de conclure à l'équivalence de 46~\% à $T = 8$, 80~\% à $T = 15$ et 98~\% à
$T = 30$. Autrement dit, à $T = 8$, l'échec à démontrer l'équivalence est le résultat le plus
probable \emph{même lorsque la proportionnalité est parfaitement vérifiée}~: ce résultat ne doit
donc pas être présenté comme un indice défavorable, mais comme une limite d'information.

\medskip
\begin{encadre}
\textbf{Formulation écartée.} Poser $H_0 : a > 0$ et $H_0 : a < 0$ comme deux tests séparés sans
marge ne répond pas au besoin~: rejeter $H_0 : a \ge 0$ prouve $a < 0$, donc une preuve
\emph{contre} la proportionnalité. Quant à l'union $H_0 : a \neq 0$, elle n'est pas testable~:
la valeur $a = 0$ est adhérente à cet ensemble, la borne supérieure de la probabilité de rejet
sur $H_0$ y vaut 1, et aucun test de niveau $\alpha < 1$ n'existe. L'introduction d'une marge
$\Delta > 0$ strictement positive est ce qui rend le problème bien posé.
\end{encadre}
\end{fiche}

\begin{fiche}{Test de Student sur la pente}
\refl{Student (1908), \emph{op.\ cit.}}
\Cadre
Identique à la fiche précédente (mêmes hypothèses de validité non testées).
\Hzero
$H_0 : b = 0$, absence de tout lien linéaire entre le volume et les pertes.
\Stat
\[
  t_b \;=\; \frac{\hat b}{\widehat{\mathrm{se}}(\hat b)},
  \qquad
  \hat b = \frac{S_{xy}}{S_{xx}},
  \qquad
  \widehat{\mathrm{se}}(\hat b) = \frac{\hat s}{\sqrt{S_{xx}}}.
\]
\Loi
$t_b \sim t(T-2)$.
\Pval
$p = 2\big(1 - F_{t,T-2}(|t_b|)\big)$.
\Usage
\textbf{Ce test s'interprète en sens inverse des autres tests de cette section}~: on souhaite ici
\emph{rejeter} $H_0$. Une pente non significative signifierait que le volume n'explique pas les
pertes, ce qui priverait de fondement le modèle réglementaire sur ce segment. Le script traduit
cette inversion dans ses verdicts ($p$ faible $\Rightarrow$ \code{OK}). Limites~: avec $T = 10$
la puissance reste modeste~; et la significativité de la pente ne dit rien sur la
\emph{proportionnalité} (absence de constante), qui relève de la fiche précédente. Les deux tests
sont complémentaires et non substituables.
\end{fiche}

\begin{fiche}{Test de Fisher de significativité globale}
\refl{R.\,A. Fisher, \emph{On the mathematical foundations of theoretical statistics},
\emph{Philosophical Transactions of the Royal Society A}, 222, 1922, p.~309--368 ;
\emph{Statistical Methods for Research Workers}, Oliver \& Boyd, Édimbourg, 1925.}
\Cadre
Identique aux deux fiches précédentes. On note $k$ le nombre de régresseurs hors constante
(ici $k=1$).
\Hzero
$H_0 : b = 0$ (plus généralement, nullité conjointe de tous les coefficients de pente).
\Stat
\[
  F \;=\; \frac{(\mathrm{SCT} - \mathrm{SCR})/k}{\mathrm{SCR}/(T-k-1)}
    \;=\; \frac{R^2/k}{(1-R^2)/(T-k-1)} .
\]
\Loi
$F \sim \mathcal{F}(k,\,T-k-1)$, soit ici $\mathcal{F}(1, T-2)$.
\Pval
$p = 1 - F_{\mathcal{F},k,T-k-1}(F)$.
\Usage
On souhaite, comme pour la pente, \emph{rejeter} $H_0$. En régression simple ce test est
strictement équivalent au test de Student sur la pente ($F = t_b^2$, p-values identiques)~: il est
conservé pour la lisibilité du rapport et redeviendrait informatif si l'on ajoutait des
régresseurs (indicatrices d'année, variable d'inflation). Même limite de puissance à $T$ faible.
\end{fiche}

\begin{fiche}{Coefficient de détermination $R^2$}
\refl{Sortie standard de \code{lm()}. Sur la loi sous $H_0$~: R.\,A. Fisher, \emph{The general
sampling distribution of the multiple correlation coefficient}, \emph{Proceedings of the Royal
Society A}, 121, 1928, p.~654--673.}
\Cadre
Même régression auxiliaire. \textbf{Il ne s'agit pas d'un test} mais d'un indicateur descriptif ;
les rubriques sont renseignées par cohérence de présentation.
\Hzero
Aucune en tant qu'indicateur. Si l'on souhaite néanmoins tester $R^2$, l'hypothèse pertinente est
celle du test de Fisher ci-dessus, $H_0 : b = 0$.
\Stat
\[
  R^2 = 1 - \frac{\mathrm{SCR}}{\mathrm{SCT}} = \frac{S_{xy}^2}{S_{xx}S_{yy}},
  \qquad
  R^2_{\text{aj}} = 1 - (1-R^2)\frac{T-1}{T-k-1}.
\]
\Loi
Sous $H_0 : b = 0$ et normalité des erreurs,
$R^2 \sim \mathcal{B}\!\left(\tfrac{k}{2},\ \tfrac{T-k-1}{2}\right)$ (loi bêta), d'espérance
$k/(T-1)$.
\Pval
Celle du test de Fisher,
$p = 1 - F_{\mathcal{F},k,T-k-1}\!\big(\tfrac{R^2/k}{(1-R^2)/(T-k-1)}\big)$.
Le script n'affiche pas de p-value sur cette ligne et la marque comme indicateur.
\Usage
La loi bêta donne la clé de lecture essentielle~: sous $H_0$, l'espérance de $R^2$ vaut
$k/(T-1)$, soit $11\,\%$ pour $T=10$. Autrement dit, un $R^2$ apparemment « correct » est
mécaniquement facile à atteindre sur un petit échantillon et ne doit jamais être lu comme une
validation du modèle. Le script signale $R^2 < 0{,}5$ comme point d'attention, et il convient de
privilégier $R^2_{\text{aj}}$, qui pénalise le nombre de paramètres.
\end{fiche}

\begin{fiche}{Test RESET de Ramsey}
\refl{J.\,B. Ramsey, \emph{Tests for specification errors in classical linear least-squares
regression analysis}, \emph{Journal of the Royal Statistical Society, Series B}, 31(2), 1969,
p.~350--371.}
\Cadre
On part du modèle contraint réglementaire $y_t = \beta x_t + \varepsilon_t$ (sans constante),
d'où $\hat y_t = \hat\beta x_t$, puis l'on estime le modèle augmenté
$y_t = \beta x_t + \gamma_2 \hat y_t^2 + \gamma_3 \hat y_t^3 + \varepsilon_t$.
Hypothèses non testées~: erreurs indépendantes, homoscédastiques, normales. Nécessite
$T > k + q = 3$.
\Hzero
$H_0 : \gamma_2 = \gamma_3 = 0$, soit « la forme linéaire est correctement spécifiée ».
\Stat
Avec $\mathrm{SCR}_0$ et $\mathrm{SCR}_1$ les sommes de carrés résiduels des modèles contraint et
augmenté, $q = 2$ restrictions et $k = 1$ paramètre contraint~:
\[
  F \;=\; \frac{(\mathrm{SCR}_0 - \mathrm{SCR}_1)/q}{\mathrm{SCR}_1/(T-k-q)} .
\]
\Loi
$F \sim \mathcal{F}(q,\,T-k-q)$, soit $\mathcal{F}(2, T-3)$.
\Pval
$p = 1 - F_{\mathcal{F},q,T-k-q}(F)$.
\Usage
On souhaite \emph{ne pas} rejeter. Le test est générique~: il détecte une non-linéarité sans en
indiquer la forme. Limites à $T$ faible~: avec $T = 5$ il ne reste que 2 degrés de liberté au
dénominateur, ce qui le rend presque inopérant~; et $\hat y^2$, $\hat y^3$ sont fortement
colinéaires lorsque l'amplitude des volumes est faible, ce qui réduit encore la puissance. À
utiliser en complément, non en substitution, de l'examen graphique.
\end{fiche}

\begin{fiche}{Corrélation de rang de Spearman (ratio contre volume, ratio contre temps)}
\refl{C.~Spearman, \emph{The proof and measurement of association between two things},
\emph{American Journal of Psychology}, 15(1), 1904, p.~72--101.}
\Cadre
Couple de séries $(u_t, v_t)$~: $(r_t, x_t)$ pour la variante « ratio contre volume »,
$(r_t, t)$ pour la variante « ratio contre temps ». On note $R_t$, $Q_t$ les rangs respectifs et
$d_t = R_t - Q_t$. Hypothèses non testées~: observations indépendantes, variables continues (peu
ou pas d'ex \ae quo).
\Hzero
$H_0$ : indépendance des deux séries, donc absence de toute association monotone.
\Stat
En l'absence d'ex \ae quo,
\[
  \hat\rho_s = 1 - \frac{6\sum_{t=1}^{T} d_t^2}{T(T^2-1)} ,
  \qquad\text{le script rapportant } S = \sum_t d_t^2 .
\]
En présence d'ex \ae quo, $\hat\rho_s$ est le coefficient de Pearson calculé sur les rangs.
\Loi
Sous $H_0$, la loi exacte de $S$ est la loi de permutation~: les $T!$ appariements de rangs sont
équiprobables. Pour $T$ modéré on utilise l'approximation
\[
  t = \hat\rho_s\sqrt{\frac{T-2}{1-\hat\rho_s^2}} \;\sim\; t(T-2).
\]
\Pval
Approximation employée par le script (appel avec \code{exact = FALSE})~:
$p = 2\big(1 - F_{t,T-2}(|t|)\big)$. La p-value exacte s'obtient par énumération des permutations.
\Usage
Test non paramétrique robuste aux valeurs extrêmes et à la non-normalité, bon complément des
tests paramétriques précédents. Deux limites~: (i) il ne détecte que les associations
\emph{monotones} -- une relation en U passera inaperçue~; (ii) l'approximation de Student est
médiocre pour $T < 10$, où la loi exacte de permutation devrait être préférée. Sur la variante
« ratio contre volume », un rejet signale un effet d'échelle non capté par la proportionnalité
stricte.
\end{fiche}

\begin{fiche}{Test de tendance de Mann--Kendall}
\refl{H.\,B. Mann, \emph{Nonparametric tests against trend}, \emph{Econometrica}, 13(3), 1945,
p.~245--259 ; M.\,G. Kendall, \emph{Rank Correlation Methods}, 4\textsuperscript{e} éd., Griffin,
Londres, 1975.}
\Cadre
Série chronologique $r_1,\dots,r_T$. Hypothèses non testées~: sous $H_0$ les $r_t$ sont
indépendants et identiquement distribués, la variable est continue. On note $g$ le nombre de
groupes d'ex \ae quo et $t_g$ leurs effectifs.
\Hzero
$H_0$ : absence de tendance monotone, les $r_t$ étant i.i.d. (toutes les permutations
chronologiques sont équiprobables).
\Stat
\[
  S = \sum_{i=1}^{T-1}\sum_{j=i+1}^{T} \operatorname{sign}(r_j - r_i),
  \qquad
  \operatorname{Var}(S) = \frac{T(T-1)(2T+5) - \sum_g t_g(t_g-1)(2t_g+5)}{18},
\]
puis la statistique standardisée avec correction de continuité
\[
  Z = \begin{cases}
    (S-1)/\sqrt{\operatorname{Var}(S)} & \text{si } S > 0,\\[2pt]
    0 & \text{si } S = 0,\\[2pt]
    (S+1)/\sqrt{\operatorname{Var}(S)} & \text{si } S < 0.
  \end{cases}
\]
\Loi
Sous $H_0$, $S$ suit une loi exacte tabulée, symétrique autour de~0 ; asymptotiquement
$Z \xrightarrow{\ \mathcal{L}\ } \mathcal{N}(0,1)$.
\Pval
$p = 2\big(1 - \Phi(|Z|)\big)$.
\Usage
Test de référence pour détecter une dérive progressive du ratio S/P (dérive tarifaire, inflation
des sinistres, évolution de la politique de provisionnement), ce qui contredirait directement la
constance de $\beta$. Limites~: l'approximation normale n'est acceptable qu'à partir de
$T \approx 10$, en-deçà la loi exacte est nécessaire. Surtout, le test suppose
\emph{l'indépendance} des observations~: en présence d'autocorrélation positive il rejette
beaucoup trop souvent, ce qui impose de le lire conjointement avec les tests de la section~7.
\end{fiche}

\begin{fiche}{Test de tendance par signes de Cox--Stuart}
\refl{D.\,R. Cox, A.~Stuart, \emph{Some quick sign tests for trend in location and dispersion},
\emph{Biometrika}, 42(1/2), 1955, p.~80--95.}
\Cadre
On pose $c = \lceil T/2 \rceil$ et l'on forme les paires $(r_t, r_{t+c})$. Soit
$D_t = r_{t+c} - r_t$ et $m$ le nombre de paires avec $D_t \neq 0$. Hypothèses non testées~:
indépendance des observations, continuité.
\Hzero
$H_0$ : absence de tendance, formalisée par $\mathbb{P}(D_t > 0) = 1/2$ pour toute paire.
\Stat
$K = \#\{t : D_t > 0\}$, nombre de différences positives parmi les $m$ paires non nulles.
\Loi
$K \sim \mathcal{B}(m,\,1/2)$, loi binomiale \emph{exacte} (aucune approximation asymptotique).
\Pval
Test binomial bilatéral exact~:
\[
  p = \sum_{j\,:\,\mathbb{P}(K=j)\,\le\,\mathbb{P}(K=k_{\text{obs}})} \binom{m}{j}\left(\tfrac12\right)^{m},
\]
somme des probabilités des issues au moins aussi improbables que celle observée.
\Usage
Avantage décisif~: la loi est exacte, sans hypothèse de normalité ni approximation asymptotique.
\textbf{Limite rédhibitoire à petit $T$}~: le test n'utilise que le signe de $m \approx T/2$
différences. Pour $T = 10$, $m = 5$ et la p-value bilatérale minimale atteignable vaut
$2 \times (1/2)^5 = 0{,}0625 > 0{,}05$~: \emph{le test ne peut structurellement jamais rejeter au
seuil de 5~\%}, quelle que soit la force de la tendance. Il faut $m \ge 6$ (donc $T \ge 12$) pour
qu'un rejet à 5~\% devienne possible. À $T$ faible, ce test est purement indicatif et
Mann--Kendall doit lui être préféré.
\end{fiche}

\newpage
\clearpage
% =============================================================================
\section{Hypothèse H2 --- structure de variance quadratique}
% =============================================================================

\begin{encadre}
\textbf{Énoncé réglementaire} (annexe~XVII, B/C(2)(f)(ii)). La variance des pertes agrégées est
$\operatorname{Var}(Y_t) = \sigma^2\big[(1-\delta)\bar x\, x_t + \delta x_t^2\big]$, mélange
convexe d'une composante linéaire en volume (risque mutualisable) et d'une composante quadratique
(risque systématique). Le paramètre $\delta \in [0,1]$ est estimé conjointement à $\gamma$.
\end{encadre}

Les tests de cette section portent sur les résidus \emph{déjà normalisés} $z_t$~: si la
pondération $\hat\pi_t$ capture correctement l'hétéroscédasticité, il ne doit subsister
\emph{aucune} relation entre $z_t^2$ et $x_t$. Un rejet signale donc une structure de variance
mal spécifiée, et non une hétéroscédasticité en soi (qui est attendue sur les données brutes).

\begin{fiche}{Test de Breusch--Pagan, version studentisée de Koenker}
\refl{T.\,S. Breusch, A.\,R. Pagan, \emph{A simple test for heteroscedasticity and random
coefficient variation}, \emph{Econometrica}, 47(5), 1979, p.~1287--1294 ; version robuste
$T R^2$ due à R.~Koenker, \emph{A note on studentizing a test for heteroscedasticity},
\emph{Journal of Econometrics}, 17(1), 1981, p.~107--112.}
\Cadre
Régression auxiliaire $z_t^2 = c_0 + c_1 x_t + \eta_t$, de coefficient de détermination
$R^2_{\text{aux}}$. Hypothèses non testées~: indépendance des observations et validité de
l'approximation asymptotique.
\Hzero
$H_0 : c_1 = 0$ : la variance des résidus normalisés ne dépend pas du volume, donc la pondération
réglementaire $\pi_t$ est correctement spécifiée.
\Stat
\[
  \mathrm{LM} \;=\; T \cdot R^2_{\text{aux}} .
\]
\Loi
$\mathrm{LM} \xrightarrow{\ \mathcal{L}\ } \chi^2(q)$, $q$ étant le nombre de régresseurs
auxiliaires hors constante, ici $q = 1$.
\Pval
$p = 1 - F_{\chi^2_q}(\mathrm{LM})$.
\Usage
On souhaite \emph{ne pas} rejeter. Limites sérieuses à $T$ faible~: la convergence vers la loi du
khi-deux est lente et la taille effective du test s'écarte notablement du niveau nominal pour
$T < 30$. De plus, la forme testée est \emph{linéaire} en $x_t$~: une hétéroscédasticité
résiduelle de forme différente peut échapper au test, d'où la variante de White. Enfin, la
version originale de Breusch--Pagan exige la normalité des erreurs~; la forme $T R^2$ retenue ici
(Koenker) s'en affranchit et est donc préférable dans ce contexte.
\end{fiche}

\begin{fiche}{Test de Breusch--Pagan, version originale de 1979}
\refl{T.\,S. Breusch, A.\,R. Pagan, \emph{A simple test for heteroscedasticity and random
coefficient variation}, \emph{Econometrica}, 47(5), 1979, p.~1287--1294.}
\Cadre
Mêmes données et mêmes notations que la fiche précédente. \textbf{Hypothèse de validité
supplémentaire, et déterminante~: la normalité des erreurs.} C'est elle qui justifie
$\operatorname{Var}(u_t^2) = 2\sigma^4$, d'où le facteur $\tfrac12$ de la statistique. Cette
hypothèse n'est pas testée ici~: elle relève de H3, examinée séparément en partie~III.
\Hzero
$H_0 : c_1 = 0$, la variance des résidus normalisés ne dépend pas du volume.
\Stat
En posant $g_t = z_t^2 / \hat\sigma^2$ avec $\hat\sigma^2 = T^{-1}\sum_t z_t^2$, on régresse $g$
sur $x$ et l'on retient la moitié de la somme des carrés expliqués~:
\[
  \mathrm{LM} = \tfrac{1}{2}\sum_{t=1}^{T}\big(\hat g_t - \bar g\big)^2 .
\]
\Loi
$\mathrm{LM} \xrightarrow{\ \mathcal{L}\ } \chi^2(q)$ sous $H_0$ \emph{et} normalité des erreurs,
avec $q = 1$ régresseur auxiliaire.
\Pval
$p = 1 - F_{\chi^2_1}(\mathrm{LM})$, p-value non simulée. Une p-value de Monte-Carlo est calculée
en parallèle et reste celle retenue pour le verdict à $T = 8$.
\Usage
Cette version est celle qui est publiée et citée dans la littérature~; elle figure dans l'outil à
ce titre, et pour permettre la confrontation avec la version studentisée. \textbf{Elle n'est pas
robuste à la non-normalité}~: si les erreurs ont des queues plus lourdes que la normale, alors
$\operatorname{Var}(u_t^2) > 2\sigma^4$, le facteur $\tfrac12$ sous-estime la variance de la
statistique et le test sur-rejette massivement. L'ampleur du phénomène a été mesurée par
simulation.

\begin{center}
\small
\rowcolors{2}{tablerow}{white}
\begin{tabular}{@{}p{6.4cm}cc@{}}
\rowcolor{tableheader}
\textcolor{financeblue}{\textbf{Loi des erreurs sous $H_0$}} &
\textcolor{financeblue}{\textbf{BP 1979}} &
\textcolor{financeblue}{\textbf{Koenker 1981}} \\
Normale, $T = 8$    & 0,021 & 0,043 \\
Normale, $T = 30$   & 0,039 & 0,047 \\
Normale, $T = 200$  & 0,051 & 0,051 \\
Student à 3 ddl (queues lourdes), $T = 200$ & \textbf{0,380} & 0,041 \\
Exponentielle (asymétrique), $T = 200$      & \textbf{0,270} & 0,038 \\
\end{tabular}
\end{center}

\noindent
Fréquence de rejet au seuil nominal de 5~\% sous $H_0$ vraie, 4\,000 réplications. Sous erreurs
normales, les deux versions convergent vers le niveau nominal, la version originale étant
seulement conservatrice en petit échantillon. Sous erreurs à queues lourdes, la version originale
rejette près de huit fois trop souvent, alors que la version studentisée reste calibrée.

\medskip
\begin{encadre}
\textbf{Règle de lecture.} Les deux versions doivent être lues \emph{ensemble}. Une divergence
marquée entre elles ne signale pas une hétéroscédasticité, mais un écart à la normalité~: elle
renvoie à H3 et non à H2. Pour conclure sur la structure de variance, c'est la version
studentisée qui fait foi, la normalité n'étant jamais acquise à $T = 8$.
\end{encadre}
\end{fiche}

\begin{fiche}{Test de White}
\refl{H.~White, \emph{A heteroskedasticity-consistent covariance matrix estimator and a direct
test for heteroskedasticity}, \emph{Econometrica}, 48(4), 1980, p.~817--838.}
\Cadre
Régression auxiliaire enrichie $z_t^2 = c_0 + c_1 x_t + c_2 x_t^2 + \eta_t$. Mêmes hypothèses de
validité que Breusch--Pagan. Nécessite $T > 3$.
\Hzero
$H_0 : c_1 = c_2 = 0$.
\Stat
$\mathrm{LM} = T \cdot R^2_{\text{aux}}$, avec le $R^2$ de la régression enrichie.
\Loi
$\mathrm{LM} \xrightarrow{\ \mathcal{L}\ } \chi^2(q)$ avec $q = 2$.
\Pval
$p = 1 - F_{\chi^2_2}(\mathrm{LM})$.
\Usage
Plus général que Breusch--Pagan (il capte une hétéroscédasticité quadratique) au prix d'un degré
de liberté supplémentaire, ce qui le rend peu adapté au régime USP~: avec $T = 5$ à $8$, estimer
trois coefficients auxiliaires laisse trop peu d'information et la puissance devient très faible.
Le test est par ailleurs un test de mauvaise spécification \emph{générale} plutôt que
d'hétéroscédasticité stricto sensu~: un rejet peut aussi traduire une erreur de forme
fonctionnelle relevant de H1.
\end{fiche}

\begin{fiche}{Test de Goldfeld--Quandt}
\refl{S.\,M. Goldfeld, R.\,E. Quandt, \emph{Some tests for homoscedasticity}, \emph{Journal of
the American Statistical Association}, 60(310), 1965, p.~539--547.}
\Cadre
Résidus triés selon $x_t$ croissant, puis scindés en deux sous-groupes de taille
$h = \lfloor T/2 \rfloor$ (observation centrale écartée si $T$ impair). Hypothèses non testées~:
normalité et indépendance des résidus~; la variance, si elle varie, est monotone en $x_t$.
\Hzero
$H_0 : \sigma_1^2 = \sigma_2^2$, égalité des variances entre le sous-groupe des faibles volumes et
celui des forts volumes.
\Stat
\[
  F \;=\; \frac{\frac{1}{h}\sum_{t \in \text{groupe haut}} z_t^2}
               {\frac{1}{h}\sum_{t \in \text{groupe bas}} z_t^2}.
\]
\Loi
$F \sim \mathcal{F}(h, h)$ sous $H_0$ et normalité. \emph{Approximation}~: les $z_t$ provenant
d'un ajustement global et non de deux régressions séparées, les degrés de liberté sont légèrement
surestimés.
\Pval
Bilatérale~: $p = 2\min\big(F_{\mathcal{F},h,h}(F),\, 1 - F_{\mathcal{F},h,h}(F)\big)$.
\Usage
Test simple et lisible dont le principal défaut est de \emph{jeter la moitié de l'information}
utile~: il ne compare que deux blocs et ignore la structure interne à chacun. Avec $T = 10$, on
compare deux blocs de 5 observations~: un rapport de variances de 2 ou 3 ne sera pas détecté.
Sensible à la non-normalité, contrairement à Brown--Forsythe.
\end{fiche}

\begin{fiche}{Test de Brown--Forsythe}
\refl{M.\,B. Brown, A.\,B. Forsythe, \emph{Robust tests for the equality of variances},
\emph{Journal of the American Statistical Association}, 69(346), 1974, p.~364--367 ; variante
robuste de H.~Levene, \emph{Robust tests for equality of variances}, in \emph{Contributions to
Probability and Statistics}, Stanford University Press, 1960.}
\Cadre
Deux groupes $g \in \{1,2\}$ définis par $x_t \le \operatorname{med}(x)$ et
$x_t > \operatorname{med}(x)$, d'effectifs $n_1, n_2$, $N = n_1+n_2$, $k=2$. On calcule les écarts
absolus à la médiane du groupe, $d_{gi} = |z_{gi} - \operatorname{med}_g(z)|$, de moyennes
$\bar d_g$ et $\bar d$. Hypothèse non testée~: indépendance des observations. La normalité n'est
\emph{pas} requise, c'est l'apport de la méthode.
\Hzero
$H_0$ : égalité des dispersions entre les deux groupes.
\Stat
\[
  W \;=\; \frac{(N-k)\sum_{g=1}^{k} n_g\,(\bar d_g - \bar d)^2}
                {(k-1)\sum_{g=1}^{k}\sum_{i=1}^{n_g} (d_{gi} - \bar d_g)^2}.
\]
\Loi
$W \sim \mathcal{F}(k-1,\,N-k)$ approximativement, soit ici $\mathcal{F}(1, T-2)$.
\Pval
$p = 1 - F_{\mathcal{F},k-1,N-k}(W)$.
\Usage
L'usage de la médiane (plutôt que de la moyenne, comme chez Levene) confère une bonne robustesse
aux distributions asymétriques ou à queues lourdes, ce qui est précieux en assurance non-vie.
Limite~: comme Goldfeld--Quandt, il ne compare que deux blocs et la loi de Fisher n'est
qu'approximative. À $T \le 10$, à considérer comme un simple indicateur.
\end{fiche}

\begin{fiche}{Test de Smirnov à deux échantillons}
\refl{N.\,V. Smirnov, \emph{Sur les écarts de la courbe de distribution empirique},
\emph{Recueil Mathématique (Matematicheskii Sbornik)}, 6(48), 1939, p.~3--26.}
\Cadre
Deux sous-échantillons de résidus normalisés~: $z^{(1)}$ (années à faible volume, effectif $n_1$)
et $z^{(2)}$ (années à fort volume, effectif $n_2$), de fonctions de répartition empiriques
$\hat F_1$, $\hat F_2$. Hypothèses non testées~: indépendance intra et inter échantillons,
continuité des lois. Le script n'active ce test que si $T \ge 8$ et $n_1, n_2 \ge 3$.
\Hzero
$H_0 : F_1 = F_2$, égalité complète des lois -- hypothèse plus forte que la seule égalité des
variances.
\Stat
\[
  D_{n_1,n_2} \;=\; \sup_{u \in \mathbb{R}} \big| \hat F_1(u) - \hat F_2(u) \big| .
\]
\Loi
Sous $H_0$, la loi exacte de $D_{n_1,n_2}$ est combinatoire (dénombrement de chemins), calculée
exactement par \code{ks.test} en l'absence d'ex \ae quo. Asymptotiquement, avec
$\lambda = \sqrt{n_1 n_2/(n_1+n_2)}$~:
\[
  \mathbb{P}\big(\lambda D_{n_1,n_2} \le u\big) \;\longrightarrow\;
  K(u) = 1 - 2\sum_{j=1}^{\infty} (-1)^{j-1} e^{-2 j^2 u^2}.
\]
\Pval
Exacte par énumération pour $n_1 n_2$ petit ; sinon
\[
  p \approx 2\sum_{j\ge1} (-1)^{j-1}\exp\!\big(-2 j^2 \lambda^2 D_{n_1,n_2}^2\big).
\]
\Usage
Complément utile aux tests de variance~: il détecte aussi un décalage de niveau ou un changement
de forme entre petits et gros volumes. Deux limites~: (i) puissance faible à effectifs réduits --
avec $n_1 = n_2 = 5$, la plus petite p-value atteignable est $2/\binom{10}{5} \approx 0{,}008$,
le rejet reste donc possible mais exige un écart maximal~; (ii) les deux échantillons ne sont pas
strictement indépendants puisque les $z_t$ partagent les mêmes paramètres estimés
$(\hat\beta, \hat\delta, \hat\gamma)$, ce qui invalide légèrement la loi de référence. Le
\emph{QQ-plot à deux échantillons} associé (section~9) reste le diagnostic le plus lisible.
\end{fiche}

\begin{fiche}{Position de $\hat\delta$ dans son domaine (diagnostic)}
\refl{Annexe XVII, sections B et C, paragraphe~6 (contrainte $0 \le \delta \le 1$).}
\Cadre
$\hat\delta$ est obtenu par minimisation sous contrainte de bornes de la log-vraisemblance
négative. \textbf{Il ne s'agit pas d'un test} mais d'un diagnostic d'identifiabilité.
\Hzero
Sans objet en tant que diagnostic. Le test formel correspondant est le test du rapport de
vraisemblance $H_0 : \delta = 0$ ou $\delta = 1$, documenté en section~10.
\Stat
La valeur $\hat\delta$ elle-même, avec le drapeau
$\mathbb{1}\{\hat\delta < 10^{-6}\ \text{ou}\ \hat\delta > 1 - 10^{-6}\}$.
\Loi
Sans objet.
\Pval
Sans objet~; le script marque la ligne \code{ALERTE} si la solution est au bord, \code{OK} sinon.
\Usage
Une solution \emph{au bord} signifie que les données ne permettent pas d'identifier un mélange
interne des deux composantes de variance~: la vraisemblance est monotone en $\delta$ sur $[0,1]$
et l'optimum est poussé sur une borne. Ce n'est pas une erreur mais un signal fort à documenter,
car (i) l'écart-type asymptotique de $\hat\delta$ n'est alors pas défini de manière usuelle,
(ii) le paramètre final peut varier fortement selon la borne retenue, et (iii) la théorie
asymptotique standard des tests sur $\delta$ ne s'applique plus (voir la correction de Chernoff
en section~10). Sur des historiques de 5 à 12 points, la solution au bord est le cas de figure le
plus fréquent en pratique.
\end{fiche}

\newpage
\clearpage
% =============================================================================
\section{Hypothèse H3 --- lognormalité des pertes agrégées}
% =============================================================================

\begin{encadre}
\textbf{Énoncé réglementaire} (annexe~XVII, B/C(2)(f)(iii)). Les pertes agrégées (ou le montant
de liquidation) suivent une loi lognormale. Cette hypothèse se teste de façon équivalente comme
la normalité des résidus normalisés $z_t$.
\end{encadre}

\begin{encadre}
\textbf{Point commun à toute cette section.} Les $z_t$ ne sont pas des observations brutes mais
des résidus d'un modèle dont trois paramètres $(\hat\beta, \hat\delta, \hat\gamma)$ ont été
estimés sur les mêmes données. La loi de référence de tous les tests de normalité en est modifiée
(elle devient stochastiquement plus concentrée), ce qui rend ces tests \emph{conservateurs} si
l'on utilise les p-values tabulées pour paramètres connus. D'où le recalcul systématique d'une
p-value de Monte-Carlo (section~\ref{sec:mc}).
\end{encadre}

\begin{fiche}{Test de Shapiro--Wilk}
\refl{S.\,S. Shapiro, M.\,B. Wilk, \emph{An analysis of variance test for normality (complete
samples)}, \emph{Biometrika}, 52(3/4), 1965, p.~591--611 ; normalisation de J.\,P. Royston,
\emph{Approximating the Shapiro-Wilk W-test for non-normality}, \emph{Statistics and Computing},
2, 1992, p.~117--119.}
\Cadre
Échantillon $z_1,\dots,z_T$ supposé (hypothèse non testée) i.i.d., sans ex \ae quo. Valable pour
$3 \le T \le 5000$. On note $m = (m_1,\dots,m_T)'$ le vecteur des espérances des statistiques
d'ordre d'un échantillon normal centré réduit et $V$ leur matrice de covariance.
\Hzero
$H_0$ : les $z_t$ sont issus d'une loi normale $\mathcal{N}(\mu,\sigma^2)$, $\mu$ et $\sigma$ non
spécifiés.
\Stat
\[
  W \;=\; \frac{\Big(\sum_{i=1}^{T} a_i\, z_{(i)}\Big)^{2}}{\sum_{t=1}^{T} (z_t - \bar z)^2},
  \qquad
  a' = \frac{m' V^{-1}}{\sqrt{m' V^{-1} V^{-1} m}} .
\]
$W \in\, ]0,1]$, une valeur proche de 1 étant favorable à la normalité.
\Loi
Aucune forme analytique fermée. Royston (1992) fournit une transformation normalisante~: pour
$12 \le T \le 5000$, $\ln(1-W)$ est approximativement $\mathcal{N}(\mu_W, \sigma_W^2)$, où $\mu_W$
et $\sigma_W$ sont des polynômes en $\ln T$~; pour $4 \le T \le 11$, la transformation utilisée
est $-\ln(\varsigma - \ln(1-W))$ avec $\varsigma$ dépendant de $T$.
\Pval
$p = 1 - \Phi\!\left(\dfrac{\ln(1-W) - \mu_W}{\sigma_W}\right)$ pour $T \ge 12$, avec la
transformation adaptée pour $T \le 11$.
\Usage
Test de référence pour les petits échantillons et généralement le plus puissant contre un large
éventail d'alternatives. Ses limites en régime USP restent néanmoins sévères~: à $T = 8$, sa
puissance contre une loi à queues lourdes réaliste n'est que de quelques dizaines de pour cent,
si bien qu'un $p$ élevé n'apporte quasiment aucune information. Il est par ailleurs sensible aux
ex \ae quo (fréquents si les montants sont arrondis) et à l'estimation préalable des paramètres.
\textbf{Lire prioritairement \code{p\_mc}.}
\end{fiche}

\begin{fiche}{Test de Shapiro--Francia}
\refl{S.\,S. Shapiro, R.\,S. Francia, \emph{An approximate analysis of variance test for
normality}, \emph{Journal of the American Statistical Association}, 67(337), 1972, p.~215--216 ;
p-values selon J.\,P. Royston, \emph{A toolkit for testing for non-normality in complete and
censored samples}, \emph{The Statistician}, 42(1), 1993, p.~37--43.}
\Cadre
Mêmes hypothèses que Shapiro--Wilk. Valable pour $5 \le T \le 5000$. On utilise les scores normaux
approchés de Blom, $\tilde m_i = \Phi^{-1}\!\big((i - 3/8)/(T + 1/4)\big)$.
\Hzero
$H_0$ : les $z_t$ sont issus d'une loi normale.
\Stat
\[
  W' \;=\; \operatorname{corr}\big(z_{(i)},\, \tilde m_i\big)^2
  \;=\; \frac{\Big(\sum_i \tilde m_i z_{(i)}\Big)^2}
              {\Big(\sum_i \tilde m_i^2\Big)\Big(\sum_t (z_t-\bar z)^2\Big)} ,
\]
soit exactement le carré du coefficient de corrélation du QQ-plot normal.
\Loi
Pas de forme fermée. Royston (1993)~: en posant $u = \ln T$ et $v = \ln u$,
\[
  \ln(1-W') \;\approx\; \mathcal{N}(\mu, \sigma^2),\quad
  \mu = -1{,}2725 + 1{,}0521\,(v-u),\quad
  \sigma = 1{,}0308 - 0{,}26758\,(v + 2/u).
\]
\Pval
$p = 1 - \Phi\!\left(\dfrac{\ln(1-W') - \mu}{\sigma}\right)$.
\Usage
Simplification de Shapiro--Wilk (remplacement de $V^{-1}$ par l'identité), souvent plus puissante
contre les alternatives \emph{à queues lourdes} et légèrement moins contre les alternatives
asymétriques. Son intérêt est d'être la traduction numérique exacte du QQ-plot, ce qui facilite la
justification en dossier. Mêmes limites de puissance qu'à la fiche précédente.
\end{fiche}

\begin{fiche}{Test d'Anderson--Darling}
\refl{T.\,W. Anderson, D.\,A. Darling, \emph{A test of goodness of fit}, \emph{Journal of the
American Statistical Association}, 49(268), 1954, p.~765--769 ; approximations de
R.\,B. D'Agostino, M.\,A. Stephens (dir.), \emph{Goodness-of-Fit Techniques}, Marcel Dekker,
1986.}
\Cadre
Échantillon i.i.d. (hypothèse non testée), $p_{(i)} = \Phi(z_{(i)})$. Statistique de distance à la
fonction de répartition empirique (famille EDF).
\Hzero
$H_0$ : les $z_t$ suivent la loi $\mathcal{N}(0,1)$ de référence.
\Stat
\[
  A^2 \;=\; -T - \frac{1}{T}\sum_{i=1}^{T} (2i-1)\Big[\ln p_{(i)} + \ln\big(1 - p_{(T+1-i)}\big)\Big],
\]
forme discrète de
$A^2 = T\int \frac{(\hat F(u)-F_0(u))^2}{F_0(u)(1-F_0(u))}\,dF_0(u)$~: la pondération
$[F_0(1-F_0)]^{-1}$ amplifie les écarts dans les \emph{queues}.
\Loi
Sous $H_0$ à paramètres connus, $A^2$ converge vers une loi limite tabulée indépendante de $F_0$.
Lorsque les paramètres sont estimés, on utilise la statistique modifiée
$A^{*2} = A^2\big(1 + 0{,}75/T + 2{,}25/T^2\big)$, de loi différente.
\Pval
Deux p-values sont rapportées. La première est \emph{non simulée}~: approximations par
morceaux de D'Agostino \& Stephens (1986) pour le cas où moyenne et variance sont estimées, par
exemple $p \approx \exp(1{,}2937 - 5{,}709\,A^{*2} + 0{,}0186\,A^{*4})$ sur la plage
$0{,}60 < A^{*2} < 13$, avec d'autres branches ailleurs. Ce sont des ajustements empiriques et non
un résultat exact~; leur calibration a été vérifiée par simulation (5,2~\% de rejets au seuil de
5~\% à $T = 8$). La seconde est la p-value de Monte-Carlo, qui reste celle retenue pour le verdict
car elle tient compte du modèle effectivement ajusté.
\Usage
Test le plus pertinent pour un usage actuariel~: la pondération en queue le rend sensible
précisément là où se joue le capital réglementaire. Limite propre à ce contexte~: cette
sensibilité est aussi sa fragilité, car sur 8 à 12 points les deux ou trois observations extrêmes
déterminent presque entièrement $A^2$. À croiser impérativement avec les tests de valeurs
aberrantes de la section~8.
\end{fiche}

\begin{fiche}{Test de Cramér--von Mises}
\refl{H.~Cramér, \emph{On the composition of elementary errors}, \emph{Skandinavisk
Aktuarietidskrift}, 11, 1928 ; R.~von Mises, \emph{Wahrscheinlichkeit, Statistik und Wahrheit},
Springer, 1928 ; forme et approximations selon M.\,A. Stephens, \emph{EDF statistics for goodness
of fit and some comparisons}, \emph{JASA}, 69(347), 1974, p.~730--737.}
\Cadre
Identique à Anderson--Darling.
\Hzero
$H_0$ : les $z_t$ suivent $\mathcal{N}(0,1)$.
\Stat
\[
  W^2 \;=\; \frac{1}{12T} + \sum_{i=1}^{T}\left(p_{(i)} - \frac{2i-1}{2T}\right)^{2}
  \;=\; T\int \big(\hat F - F_0\big)^2\, dF_0 .
\]
\Loi
Loi limite de Cramér--von Mises sous $H_0$ à paramètres connus ; statistique modifiée
$W^{*2} = W^2(1 + 0{,}5/T)$ lorsque les paramètres sont estimés.
\Pval
Approximations par morceaux de Stephens pour le cas à paramètres estimés, rapportées comme
p-value non simulée (calibration vérifiée par simulation~: 4,9~\% de rejets au seuil de
5~\% à $T = 8$)~; la p-value de Monte-Carlo reste celle retenue pour le verdict.
\Usage
Version non pondérée d'Anderson--Darling~: elle accorde le même poids à toute la distribution, et
est donc \emph{moins} sensible aux queues, ce qui la rend plus stable mais moins pertinente pour
un usage prudentiel. Utile en recoupement~: un rejet par $A^2$ non confirmé par $W^2$ oriente vers
un problème localisé en queue plutôt que vers une non-normalité d'ensemble.
\end{fiche}

\begin{fiche}{Test de Kolmogorov--Smirnov contre $\mathcal{N}(0,1)$}
\refl{A.\,N. Kolmogorov, \emph{Sulla determinazione empirica di una legge di distribuzione},
\emph{Giornale dell'Istituto Italiano degli Attuari}, 4, 1933, p.~83--91 ; H.\,W. Lilliefors,
\emph{On the Kolmogorov-Smirnov test for normality with mean and variance unknown},
\emph{JASA}, 62(318), 1967, p.~399--402.}
\Cadre
Identique aux deux fiches précédentes.
\Hzero
$H_0$ : les $z_t$ suivent $\mathcal{N}(0,1)$.
\Stat
\[
  D \;=\; \sup_{u} \big|\hat F(u) - \Phi(u)\big|
    \;=\; \max_{1\le i\le T} \max\!\left(\frac{i}{T} - p_{(i)},\; p_{(i)} - \frac{i-1}{T}\right).
\]
\Loi
À paramètres connus, $\sqrt{T}\,D$ converge vers la loi de Kolmogorov~:
\[
  \mathbb{P}\big(\sqrt{T}\,D \le u\big) \longrightarrow K(u) = 1 - 2\sum_{j=1}^{\infty}(-1)^{j-1}e^{-2j^2u^2}.
\]
À paramètres \emph{estimés}, la loi correcte est celle de Lilliefors, stochastiquement plus
concentrée et sans forme analytique.
\Pval
$p \approx 2\sum_{j\ge1}(-1)^{j-1}\exp(-2j^2 T D^2)$, formule de Kolmogorov, rapportée à titre de
p-value non simulée. La p-value de Monte-Carlo reste celle retenue pour le verdict.
\Usage
Test historiquement dominant mais \emph{le moins puissant} de cette section~: le supremum ne
retient qu'un seul point d'écart et sous-pondère les queues. Piège classique, quantifié ici~:
appliquer la loi de Kolmogorov alors que les paramètres ont été estimés rend le
test extrêmement conservateur. Une simulation à $T = 8$ donne \textbf{zéro rejet sur 3\,000
réplications} au seuil de 5~\%, ce qui produirait une fausse impression de validation. Le test de
Lilliefors (fiche suivante) est la réponse correcte à cette situation.
\end{fiche}

\begin{fiche}{Test de Lilliefors}
\refl{H.\,W. Lilliefors, \emph{On the Kolmogorov-Smirnov test for normality with mean and variance
unknown}, \emph{Journal of the American Statistical Association}, 62(318), 1967, p.~399--402 ;
p-value~: G.\,E. Dallal, L.~Wilkinson, \emph{An analytic approximation to the distribution of
Lilliefors' test statistic for normality}, \emph{The American Statistician}, 40(4), 1986,
p.~294--296, prolongée au-delà de $0{,}10$ par les polynômes de Stephens (1974).}
\Cadre
Échantillon i.i.d. (hypothèse non testée). \textbf{Différence essentielle avec la fiche
précédente}~: la moyenne et l'écart-type ne sont pas ceux du modèle mais sont \emph{estimés sur
l'échantillon lui-même}, ce qui correspond à la situation réelle lorsque l'on teste la normalité
sans connaître les paramètres.
\Hzero
$H_0$ : les $z_t$ sont issus d'une loi normale de moyenne et de variance quelconques.
\Stat
\[
  D_{\mathrm{L}} = \sup_u \left| \hat F(u) - \Phi\!\left(\frac{u - \bar z}{s_z}\right) \right|,
\]
soit la statistique de Kolmogorov--Smirnov calculée après standardisation empirique.
\Loi
Loi de Lilliefors, \emph{stochastiquement plus concentrée} que celle de Kolmogorov puisque
l'ajustement des paramètres rapproche mécaniquement $\hat F$ de la loi de référence. Elle n'a pas
de forme analytique fermée.
\Pval
Approximation analytique de Dallal \& Wilkinson (1986), valable pour $p < 0{,}10$~:
\[
  p \approx \exp\!\Big(-7{,}01256\,K^2(T + 2{,}78019) + 2{,}99587\,K\sqrt{T + 2{,}78019}
  - 0{,}122119 + \tfrac{0{,}974598}{\sqrt{T}} + \tfrac{1{,}67997}{T}\Big),
\]
prolongée au-delà par les polynômes de Stephens. Il s'agit d'une p-value \textbf{non simulée},
disponible dès $T \ge 5$. Une p-value de Monte-Carlo est calculée en parallèle.
\Usage
C'est la version correcte du test de Kolmogorov--Smirnov lorsque les paramètres sont estimés.
La calibration a été vérifiée par simulation à
$T = 8$~: 10,6~\% de rejets au seuil nominal de 10~\% et 5,2~\% au seuil de 5~\%, à comparer aux
0~\% du test de Kolmogorov mal appliqué. Sa puissance reste néanmoins celle de la famille
Kolmogorov--Smirnov, c'est-à-dire faible devant Shapiro--Wilk et Anderson--Darling.
\end{fiche}

\begin{fiche}{Shapiro--Wilk sans la normalisation de Royston}
\refl{Shapiro \& Wilk (1965), \emph{op.\ cit.} ; loi nulle évaluée par simulation directe.
Sur la légitimité de l'invariance~: la statistique $W$ est invariante par translation et par
changement d'échelle, propriété établie dans l'article original.}
\Cadre
Même statistique $W$ que la fiche « Shapiro--Wilk »~; seule la \emph{loi de référence} change.
\Hzero
$H_0$ : les $z_t$ sont issus d'une loi normale.
\Stat
$W$, identique à la fiche précédente.
\Loi
$W$ étant invariante par translation et par changement d'échelle, sa loi sous $H_0$ ne dépend
d'\emph{aucun paramètre}~: elle ne dépend que de $T$. Elle peut donc être évaluée numériquement
en simulant des échantillons $\mathcal{N}(0,1)$ de taille $T$. L'outil emploie
20\,000 tirages, mis en cache par valeur de $T$ et pilotés par une graine dédiée.
\Pval
$\hat p = \big(1 + \#\{W^{(b)} \le W_{\text{obs}}\}\big) / (B_{\text{null}} + 1)$, rejet en queue
basse.
\Usage
Cette variante répond à une objection de principe~: la p-value de \code{shapiro.test()} repose sur
la transformation normalisante de Royston (1992), qui est un \emph{ajustement empirique}, non un
résultat exact. La version proposée ici évalue directement la loi exacte de $W$, à l'erreur de
Monte-Carlo près, en $\mathcal{O}(B_{\text{null}}^{-1/2})$.

\medskip
\textbf{Point important à ne pas confondre.} Cette simulation n'est \emph{pas} un bootstrap
paramétrique~: elle ne dépend pas du modèle USP ajusté, mais seulement de $T$. Elle relève donc
de la catégorie « exacte, évaluée numériquement » et non de la catégorie « Monte-Carlo » du
tableau~\ref{tab:nature}. En pratique, les deux versions concordent étroitement~: sur 2\,000
réplications à $T = 8$, la fréquence de rejet au seuil de 5~\% vaut 4,95~\% pour la loi simulée
contre 5,05~\% pour la normalisation de Royston, ce qui valide a posteriori l'approximation de
Royston à cette taille d'échantillon.
\end{fiche}

\begin{fiche}{Test de Jarque--Bera}
\refl{C.\,M. Jarque, A.\,K. Bera, \emph{Efficient tests for normality, homoscedasticity and
serial independence of regression residuals}, \emph{Economics Letters}, 6(3), 1980, p.~255--259.}
\Cadre
Échantillon i.i.d., moments d'ordre 4 finis (hypothèses non testées). Test du multiplicateur de
Lagrange contre la famille de Pearson.
\Hzero
$H_0$ : asymétrie et aplatissement conformes à la loi normale, soit $\mathbb{E}[\sqrt{b_1}] = 0$
et $\mathbb{E}[b_2] = 3$ conjointement.
\Stat
\[
  \mathrm{JB} \;=\; \frac{T}{6}\left(b_1 + \frac{(b_2-3)^2}{4}\right),
  \qquad
  \sqrt{b_1} = \frac{m_3}{m_2^{3/2}},\quad b_2 = \frac{m_4}{m_2^{2}} .
\]
\Loi
$\mathrm{JB} \xrightarrow{\ \mathcal{L}\ } \chi^2(2)$, comme somme des carrés de deux
statistiques asymptotiquement normales et indépendantes.
\Pval
$p = 1 - F_{\chi^2_2}(\mathrm{JB})$.
\Usage
\textbf{À manier avec une grande prudence en régime USP.} La convergence vers $\chi^2(2)$ est
notoirement lente~: pour $T < 50$, la taille effective s'écarte fortement du niveau nominal.
Surtout, l'aplatissement empirique est \emph{mécaniquement borné}, $b_2$ ne pouvant dépasser une
quantité de l'ordre de $T$~: avec $T = 8$, l'échantillon est structurellement incapable de
manifester une queue lourde et le test ne peut pratiquement pas rejeter pour excès
d'aplatissement. La p-value \code{p\_mc} est ici indispensable~; en son absence, ce test n'apporte
à peu près rien sur un historique réglementaire minimal.
\end{fiche}

\begin{fiche}{Test d'asymétrie de D'Agostino}
\refl{R.\,B. D'Agostino, \emph{Transformation to normality of the null distribution of $g_1$},
\emph{Biometrika}, 57(3), 1970, p.~679--681.}
\Cadre
Échantillon i.i.d. ; nécessite $T \ge 8$ pour que la transformation normalisante soit définie, le
script marquant le test non applicable en-deçà.
\Hzero
$H_0$ : coefficient d'asymétrie de la population nul.
\Stat
En posant $\sqrt{b_1} = m_3/m_2^{3/2}$~:
\[
  Y = \sqrt{b_1}\,\sqrt{\frac{(T+1)(T+3)}{6(T-2)}},
  \qquad
  \beta_2 = \frac{3(T^2+27T-70)(T+1)(T+3)}{(T-2)(T+5)(T+7)(T+9)},
\]
\[
  W^2 = -1 + \sqrt{2(\beta_2-1)},\qquad
  \alpha = \sqrt{\frac{2}{W^2-1}},\qquad
  \varsigma = \frac{1}{\sqrt{\ln W}},
\]
\[
  Z \;=\; \varsigma \,\ln\!\left(\frac{Y}{\alpha} + \sqrt{\left(\frac{Y}{\alpha}\right)^2 + 1}\right)
    \;=\; \varsigma\,\operatorname{arcsinh}\!\left(\frac{Y}{\alpha}\right).
\]
\Loi
$Z \sim \mathcal{N}(0,1)$ approximativement sous $H_0$, la transformation de Johnson $S_U$
ci-dessus corrigeant l'asymétrie de la loi exacte de $\sqrt{b_1}$ à $T$ fini.
\Pval
$p = 2\big(1 - \Phi(|Z|)\big)$.
\Usage
Utile pour \emph{localiser} la cause d'un rejet global de normalité~: une asymétrie significative
oriente vers un mélange de sinistres graves et attritionnels mal séparés, diagnostic actionnable
au regard de l'exigence de l'annexe~XVII sur le retraitement des événements déjà couverts par les
sous-modules catastrophe. Limite~: la correction de petit échantillon vaut dès $T \ge 8$, mais la
puissance ne devient réellement exploitable que vers $T \ge 20$.
\end{fiche}

\begin{fiche}{Test d'aplatissement d'Anscombe--Glynn}
\refl{F.\,J. Anscombe, W.\,J. Glynn, \emph{Distribution of the kurtosis statistic $b_2$ for
normal samples}, \emph{Biometrika}, 70(1), 1983, p.~227--234.}
\Cadre
Échantillon i.i.d. ; nécessite $T \ge 20$. Sur un historique réglementaire de 5 à 15 années, ce
test est donc presque toujours \emph{non applicable} et le script l'indique explicitement.
\Hzero
$H_0$ : aplatissement de la population égal à 3 (valeur normale).
\Stat
Avec $b_2 = m_4/m_2^2$,
\[
  \mathbb{E}[b_2] = \frac{3(T-1)}{T+1},
  \qquad
  \operatorname{Var}(b_2) = \frac{24T(T-2)(T-3)}{(T+1)^2(T+3)(T+5)},
  \qquad
  x = \frac{b_2 - \mathbb{E}[b_2]}{\sqrt{\operatorname{Var}(b_2)}},
\]
\[
  \sqrt{\beta_1} = \frac{6(T^2-5T+2)}{(T+7)(T+9)}\sqrt{\frac{6(T+3)(T+5)}{T(T-2)(T-3)}},
  \qquad
  A = 6 + \frac{8}{\sqrt{\beta_1}}\left(\frac{2}{\sqrt{\beta_1}} + \sqrt{1 + \frac{4}{\beta_1}}\right),
\]
\[
  Z \;=\; \frac{\Big(1 - \dfrac{2}{9A}\Big) - \left(\dfrac{1 - 2/A}{1 + x\sqrt{2/(A-4)}}\right)^{1/3}}
                {\sqrt{2/(9A)}} .
\]
\Loi
$Z \sim \mathcal{N}(0,1)$ approximativement (transformation de Wilson--Hilferty).
\Pval
$p = 2\big(1 - \Phi(|Z|)\big)$.
\Usage
C'est le test le plus directement pertinent pour un actuaire (l'épaisseur de queue conditionne le
quantile à 99,5~\%), mais aussi le moins souvent disponible~: il exige $T \ge 20$, davantage que
ce qu'exige la crédibilité pleine de l'annexe~XVII. Lorsqu'il n'est pas applicable, l'analyse de
queue doit reposer sur le QQ-plot, sur Anderson--Darling et sur les tests de valeurs aberrantes.
\end{fiche}

\newpage
\clearpage
% =============================================================================
\section{Hypothèse H4 --- indépendance et validité du maximum de vraisemblance}
\label{sec:h4}
% =============================================================================

\begin{encadre}
\textbf{Énoncé réglementaire} (annexe~XVII, B/C(2)(f)(iv)). L'estimation par maximum de
vraisemblance est appropriée, ce qui suppose notamment l'indépendance des $Y_t$ conditionnellement
aux $x_t$, sans autocorrélation résiduelle.
\end{encadre}

\begin{fiche}{Test de Durbin--Watson}
\refl{J.~Durbin, G.\,S. Watson, \emph{Testing for serial correlation in least squares regression,
I}, \emph{Biometrika}, 37(3/4), 1950, p.~409--428 ; \emph{II}, \emph{Biometrika}, 38(1/2), 1951,
p.~159--178.}
\Cadre
Série de résidus $z_1,\dots,z_T$ ordonnée chronologiquement. Hypothèses non testées~: résidus
issus d'un modèle correctement spécifié, homoscédastiques et gaussiens~; régresseurs exogènes et
non aléatoires.
\Hzero
$H_0 : \rho = 0$ dans le schéma $z_t = \rho z_{t-1} + \nu_t$, soit l'absence d'autocorrélation
d'ordre~1.
\Stat
\[
  \mathrm{DW} \;=\; \frac{\sum_{t=2}^{T}(z_t - z_{t-1})^2}{\sum_{t=1}^{T} z_t^2}
  \;\approx\; 2\,(1 - \hat\rho_1).
\]
$\mathrm{DW} \in [0,4]$~: proche de 2 sous $H_0$, proche de 0 en autocorrélation positive forte,
proche de 4 en autocorrélation négative forte.
\Loi
Sous $H_0$, les résidus normalisés sont $\mathcal{N}(0, I_T)$~: $\mathrm{DW}$ est alors le
\emph{rapport de deux formes quadratiques} en variables normales,
$\mathrm{DW} = z'Az / z'z$ où $A$ est la matrice des différences secondes. Ses valeurs propres
sont connues analytiquement, $\lambda_j = 2\big(1 - \cos(\pi j / T)\big)$, $j = 0,\dots,T-1$
(vérification numérique faite). La loi de $\mathrm{DW}$ est donc \textbf{exactement calculable},
sans tabulation ni approximation.
\Pval
On écrit $\mathbb{P}(\mathrm{DW} \le c) = \mathbb{P}\big(z'(A - cI)z \le 0\big)$, queue d'une
forme quadratique évaluée par la méthode d'Imhof (1961)~:
\[
  \mathbb{P}(Q > 0) = \frac{1}{2} - \frac{1}{\pi}\int_0^{\infty}
  \frac{\sin\theta(u)}{u\,\rho(u)}\,\mathrm{d}u,
  \quad
  \theta(u) = \tfrac12\sum_j \arctan(h_j u),
  \quad
  \rho(u) = \prod_j \big(1 + h_j^2u^2\big)^{1/4},
\]
avec $h_j = \lambda_j - c$. La p-value bilatérale est
$2\min\big(\mathbb{P}(\mathrm{DW} \le d),\, \mathbb{P}(\mathrm{DW} \ge d)\big)$.
Le centrage retirant une dimension, l'intégration porte sur les $T-1$ valeurs propres non nulles
de $A$ projetée sur l'orthogonal du vecteur constant.
\Usage
Deux points d'implémentation méritent d'être signalés. \emph{Le centrage}~: la statistique est
calculée sur des résidus centrés, car la définition classique porte sur des résidus de moindres
carrés, de moyenne nulle par construction. Les $z_t$ satisfont la contrainte
$\sum_t\sqrt{\hat\pi_t}z_t = 0$ mais pas $\bar z = 0$~; sans centrage, le dénominateur serait
majoré de $T\bar z^2$ et $\mathrm{DW}$ biaisée vers le bas. Le centrage rend en outre la
statistique cohérente avec l'autocorrélation employée par les tests portmanteau.
\emph{La p-value}~: elle est exacte, et non asymptotique~; la validation par simulation donne
10,1~\% de rejets au seuil nominal de 10~\% et 5,2~\% au seuil de 5~\% (1\,500 réplications à
$T = 8$). Les bornes $d_L/d_U$, établies pour des résidus de moindres carrés, ne sont pas
utilisées. Le test ne détecte par ailleurs qu'une dépendance d'ordre~1.
\end{fiche}

\begin{fiche}{Tests portemanteau de Ljung--Box et de Box--Pierce}
\refl{G.\,E.\,P. Box, D.\,A. Pierce, \emph{Distribution of residual autocorrelations in
autoregressive-integrated moving average time series models}, \emph{JASA}, 65(332), 1970,
p.~1509--1526 ; G.\,M. Ljung, G.\,E.\,P. Box, \emph{On a measure of lack of fit in time series
models}, \emph{Biometrika}, 65(2), 1978, p.~297--303.}
\Cadre
Série $z_1,\dots,z_T$, autocorrélations empiriques
$\hat\rho_k = \sum_{t=k+1}^{T}(z_t-\bar z)(z_{t-k}-\bar z)\big/\sum_{t=1}^{T}(z_t-\bar z)^2$.
Le script teste $h = 1$ et, si $T \ge 8$, $h = 2$. Hypothèses non testées~: stationnarité et
existence des moments d'ordre~2.
\Hzero
$H_0 : \rho_1 = \dots = \rho_h = 0$, absence d'autocorrélation jusqu'au retard $h$.
\Stat
\[
  Q_{\mathrm{BP}} = T\sum_{k=1}^{h}\hat\rho_k^2 ,
  \qquad
  Q_{\mathrm{LB}} = T(T+2)\sum_{k=1}^{h}\frac{\hat\rho_k^2}{T-k}.
\]
La pondération $(T+2)/(T-k)$ de Ljung--Box corrige le biais de $Q_{\mathrm{BP}}$ à $T$ fini.
\Loi
$Q \xrightarrow{\ \mathcal{L}\ } \chi^2(h)$ sous $H_0$ (et $\chi^2(h-p-q)$ si les résidus
proviennent d'un ARMA$(p,q)$ ajusté, ce qui n'est pas le cas ici).
\Pval
$p = 1 - F_{\chi^2_h}(Q)$.
\Usage
Ljung--Box est à préférer systématiquement à Box--Pierce, dont l'approximation est plus grossière
à $T$ faible~; le script affiche les deux en recoupement. Limites~: l'approximation $\chi^2$ n'est
fiable qu'à partir de $T \approx 30$, et la règle d'usage $h \le T/5$ impose de se restreindre à
$h \le 2$ ou $3$ sur un historique réglementaire. Un rejet au retard 2 non confirmé au retard 1
doit être interprété avec prudence~: sur 12 points, il s'agit souvent d'un artefact.
\end{fiche}

\begin{fiche}{Test des suites de Wald--Wolfowitz}
\refl{A.~Wald, J.~Wolfowitz, \emph{On a test whether two samples are from the same population},
\emph{Annals of Mathematical Statistics}, 11(2), 1940, p.~147--162.}
\Cadre
On transforme la série en une suite binaire de signes
$s_t = \operatorname{sign}(z_t - \operatorname{med}(z))$, en écartant les valeurs nulles. Soit
$n_1$ le nombre de signes positifs, $n_2$ le nombre de négatifs, $n = n_1+n_2$, et $R$ le nombre
de « suites » (séquences maximales de signes identiques). Hypothèse non testée~: continuité de la
variable.
\Hzero
$H_0$ : la suite des signes est un arrangement aléatoire, les $\binom{n}{n_1}$ arrangements étant
équiprobables.
\Stat
\[
  \mathbb{E}[R] = \frac{2n_1n_2}{n} + 1,
  \qquad
  \operatorname{Var}(R) = \frac{2n_1n_2(2n_1n_2 - n)}{n^2(n-1)},
  \qquad
  Z = \frac{R - \mathbb{E}[R]}{\sqrt{\operatorname{Var}(R)}} .
\]
\Loi
Loi exacte combinatoire pour $R$ (dénombrement des compositions) ; asymptotiquement
$Z \xrightarrow{\ \mathcal{L}\ } \mathcal{N}(0,1)$.
\Pval
Approximation normale utilisée par le script~: $p = 2\big(1 - \Phi(|Z|)\big)$.
\Usage
Test non paramétrique très général~: il détecte aussi bien une autocorrélation (trop peu de
suites) qu'une alternance systématique (trop de suites), sans hypothèse de forme. Limite
importante~: l'approximation normale n'est recommandée que pour $n_1, n_2 \ge 10$, soit
$T \ge 20$. En-deçà, la loi exacte devrait être utilisée, et la granularité de $R$ (entier compris
entre 2 et $n$) limite fortement le nombre de p-values atteignables.
\end{fiche}

\begin{fiche}{Condition du premier ordre (diagnostic numérique)}
\refl{Diagnostic numérique ; conditions d'optimalité du maximum de vraisemblance.}
\Cadre
En annulant la dérivée de l'objectif par rapport à $\ln\beta$, on obtient la contrainte que
l'estimation impose mécaniquement aux résidus \emph{quel que soit} le couple $(\delta,\gamma)$~:
\[
  \sum_{t=1}^{T} \hat\pi_t\, v_t = 0
  \qquad\Longleftrightarrow\qquad
  \sum_{t=1}^{T} \sqrt{\hat\pi_t}\; z_t = 0 .
\]
\textbf{Diagnostic, non test.}
\Hzero
Sans objet.
\Stat
Écart relatif à la condition d'optimalité~:
$\displaystyle \mathrm{FOC} = \Big|\sum_t \hat\pi_t v_t\Big| \Big/ \sum_t \hat\pi_t$.
\Loi
Sans objet.
\Pval
Aucune. Le moteur pose \code{ECHEC} si $\mathrm{FOC} > 10^{-6}$.
\Usage
\textbf{Point de vigilance essentiel.} C'est la moyenne \emph{pondérée} $\sum_t\hat\pi_t v_t$ qui
est contrainte à zéro, et non la moyenne simple $\bar z$ des résidus normalisés. Les deux ne
coïncident que si $\hat\pi_t$ est constant, c'est-à-dire si $\hat\delta = 1$ ou si les volumes
$x_t$ sont constants. Sur un jeu de données où $\hat\delta$ est intérieur au domaine, on observe
typiquement $\mathrm{FOC} \approx 10^{-17}$ (précision machine) alors que
$\bar z \approx 10^{-2}$~: la moyenne simple des résidus n'est donc \emph{pas} un contrôle de
convergence. Elle n'est pas pour autant une quantité libre~: l'écart entre les deux moyennes est
de l'ordre de la dispersion des poids $\hat\pi_t$, ce qui interdit d'en faire un test et conduit à
la restituer comme diagnostic (fiche~\ref{fiche:32}).

\textbf{Portée exacte de ce diagnostic.} $\mathrm{FOC}$ ne mesure pas non plus la convergence.
\code{usp\_noyau()} calcule $\ln\hat\beta$ en forme fermée,
$\ln\hat\beta = \big(T/2 + \sum_t\hat\pi_t\ln r_t\big)\big/\sum_t\hat\pi_t$~; en reportant dans
$v_t = \ln r_t + 1/(2\hat\pi_t) - \ln\hat\beta$, on obtient $\sum_t\hat\pi_t v_t = 0$
\emph{identiquement}, pour tout couple $(\delta,\gamma)$ et non seulement à l'optimum. Vérifié
numériquement sur des couples arbitraires --- $(0{,}5\,;-1)$, $(0\,;-3)$, $(1\,;0)$,
$(0{,}3\,;-2{,}5)$, $(1\,;3)$ --- où $\mathrm{FOC}$ reste inférieur à $10^{-15}$. Le
\code{ECHEC} de la rubrique~5 ne peut donc se déclencher qu'en cas d'anomalie arithmétique
(débordement, valeur non finie), jamais parce que l'optimisation se serait arrêtée trop tôt. La
convergence en $\gamma$, elle, est signalée par l'écart de $\sum_t z_t^2$ à $T$
(fiche~\ref{fiche:33}), seul indicateur de cette nature dans la restitution.
\end{fiche}

\begin{fiche}{Centrage des résidus standardisés (diagnostic)}
\refl{Diagnostic ; conditions du premier ordre du maximum de vraisemblance
(fiche~\ref{fiche:31}).}
\Cadre
Résidus standardisés $z_t = \sqrt{\hat\pi_t}\,v_t$ et leur moyenne simple $\bar z$.
\textbf{Diagnostic, non test.} La relation $\sum_t\sqrt{\hat\pi_t}\,z_t = 0$ vaut
\emph{toujours}, et pas seulement à l'optimum~: $\ln\hat\beta$ étant obtenu en forme fermée,
$\ln\hat\beta = \big(T/2 + \sum_t\hat\pi_t\ln r_t\big)\big/\sum_t\hat\pi_t$, elle est une
\textbf{identité algébrique} vérifiée pour tout couple $(\delta,\gamma)$, convergé ou non
(fiche~\ref{fiche:31}). $\bar z$ est donc rivé par l'estimation en toute circonstance. Cette
contrainte pondérée ne se réduit à $\sum_t z_t = 0$ que lorsque
$\hat\pi_t$ est \emph{constant}, c'est-à-dire lorsque $\hat\delta = 1$ ou lorsque les volumes
$x_t$ le sont. Le seul fait que $\hat\delta$ soit au bord de $[0,1]$ ne suffit pas~: à
$\hat\delta = 0$ avec des volumes variables, $\hat\pi_t$ varie et l'égalité $\sum_t z_t = 0$
ne tient pas. Le moteur distingue ces trois situations dans le commentaire attaché à la ligne.
\Hzero
Sans objet. L'hypothèse $\mathbb{E}[z_t] = 0$, sous laquelle cette grandeur était auparavant
restituée comme un test, ne peut pas être mise en défaut par les données~: $\bar z$ est déterminé
par l'estimation, non librement observé.
\Stat
Aucune statistique de test. La grandeur restituée est l'\emph{estimation} $\bar z$, portée par le
champ \code{estim} de la table des tests~; les champs \code{stat} et \code{loi} sont vides.
\Loi
Sans objet, pour deux raisons distinctes.
\emph{Loi nominale.} La loi $t(T-1)$ de $\bar z\sqrt{T}/s_z$ est optimiste~: les $z_t$ dépendent
de trois paramètres estimés $(\hat\beta,\hat\delta,\hat\gamma)$ et satisfont la contrainte
linéaire ci-dessus, de sorte que le nombre effectif de degrés de liberté est inférieur à $T-1$.
Cet argument, exact, reste faible~: il suggère une p-value seulement trop optimiste, donc
utilisable avec prudence. \textbf{Le fait qui tranche est plus radical.} Lorsque $\hat\pi_t$ est
constant, $\bar z = 0$ est une identité algébrique~; la valeur affichée est un zéro machine, le
rapport $\bar z\sqrt{T}/s_z$ l'est aussi, et la p-value bilatérale vaut $1$ \emph{quelles que
soient les données}. Ce n'est pas une p-value élevée~: c'est une \textbf{constante de $T$}, qui
ne contient aucune information sur l'échantillon. Conséquence algébrique exacte, vérifiée
numériquement sur deux jeux sans rapport l'un avec l'autre, à $T = 8$~:

\begin{center}
\small
\rowcolors{2}{tablerow}{white}
\begin{tabular}{@{}p{5.6cm}ccc@{}}
\rowcolor{tableheader}
\textcolor{financeblue}{\textbf{Jeu de données}} &
\textcolor{financeblue}{\textbf{$\hat\delta$}} &
\textcolor{financeblue}{\textbf{$t$ observé}} &
\textcolor{financeblue}{\textbf{p nominale bilatérale}} \\
Jeu de contrôle du dossier & 1 & $-5{,}3\cdot 10^{-16}$ & $1{,}0000000000$ \\
Jeu étranger, volumes $x_t$ constants\textsuperscript{(a)} & 0 &
$-8{,}7\cdot 10^{-17}$ & $1{,}0000000000$ \\
\end{tabular}
\end{center}

\noindent
\textsuperscript{(a)} $x_t \equiv 100$ et $y_t = 55,\,70,\,61,\,84,\,48,\,92,\,66,\,73$, jeu
construit pour cette vérification et sans rapport avec le dossier.

\noindent
Le second jeu n'a aucun point commun avec le premier et conduit même à
$\hat\delta = 0$~; il donne pourtant la même p-value, parce que la constance de $\hat\pi_t$ y
résulte non de $\hat\delta$ mais de celle des volumes. Sur les 493 réplications à
$\delta^{*} = 1$ de la simulation décrite ci-dessous, la p-value nominale vaut $1$ sur toutes,
sans exception. Un verdict fondé sur
cette p-value serait un \code{OK} qu'\textbf{aucune donnée ne peut faire perdre} --- l'assurance
d'un contrôle qui n'en est pas un, que la rubrique~6 écarte. La même vérification, conduite sur la
variance unitaire dont la loi nominale est un $\chi^2$ et non un $t$, donne le même constat
(fiche~\ref{fiche:33})~: ce n'est donc pas une particularité de la loi de Student.

\emph{Loi simulée.} La loi de $\bar z$ sous le modèle ajusté n'est pas davantage exploitable, non
parce qu'elle serait constante, mais parce qu'elle est un \emph{mélange}. Constat de simulation
sur le jeu de contrôle ($T = 8$, $\hat\delta = 1$, $B = 999$ réplications, graine 20260831)~:
493 réplications donnent $\delta^{*} = 1$ et une moyenne $\bar z^{*}$ d'écart-type
$2{,}2\cdot 10^{-16}$, soit un zéro machine~; 460 donnent $\delta^{*} = 0$ et un écart-type de
$1{,}7\cdot 10^{-2}$, composante diffuse~; 46 un $\delta^{*}$ intérieur.

Ce n'est pas le mélange en lui-même qui prive la p-value de sens, mais le fait que la loi simulée
porte un \textbf{atome en la valeur observée}. La valeur observée est elle-même un zéro machine
($\bar z = -2{,}0\cdot 10^{-16}$) et les 493 réplications de la première composante le sont aussi
($|\bar z^{*}| \le 7{,}8\cdot 10^{-16}$)~: une p-value bilatérale
$2\min\big(\mathbb{P}(\bar z^{*}\ge\bar z),\,\mathbb{P}(\bar z^{*}\le\bar z)\big)$ se décide alors
par le classement d'ex \ae quo numériques. Mesuré~: sur ces 493 réplications, 420 vérifient
$\bar z^{*} \ge \bar z$ et 74 vérifient $\bar z^{*} \le \bar z$, uniquement d'après le signe d'un
bruit de l'ordre de $10^{-16}$ --- d'où la dépendance du résultat à la plateforme de calcul. La
somme $420 + 74 = 494$ excède $493$ d'une unité~: une réplication tombe exactement sur la valeur
observée et est comptée dans les deux queues, les deux inégalités de la p-value bilatérale étant
larges.

Se restreindre à la composante diffuse ne sauverait pas la vérification~: la p-value recalculée
sur les seules 506 réplications à $\hat\pi_t^{*}$ non constant vaut $0{,}966$. La valeur observée
est par construction au centre de cette composante, de sorte qu'un test conditionnel serait
\emph{inopérant} et non simplement mal calibré. Ces deux chiffres sont des constats de simulation
sur ce jeu de données, non des propriétés démontrées du modèle.
\Pval
Aucune. La grandeur a été retirée des statistiques simulées par
\code{.stats\_bootstrapables()}~: ni p-value exacte, ni p-value de Monte-Carlo, ni p-value
asymptotique n'est calculée ni affichée pour cette ligne, et aucune p-value n'est retenue.
\Usage
\textbf{Pourquoi aucun verdict.} Une grandeur que l'estimation rive ne peut pas contredire le
modèle~; lui attacher un verdict \code{OK} donnerait au lecteur du dossier l'assurance d'un
contrôle qui n'en est pas un. La ligne est donc restituée comme diagnostic, avec le verdict
\code{INFO}, et accompagnée du motif. Ce qu'elle apporte à $T = 8$ est très limité, et il faut le
dire sans détour. Lorsque $\hat\pi_t$ est constant, $\bar z$ n'est que du bruit d'arrondi et ne
renseigne sur rien --- en particulier \textbf{pas} sur la qualité de l'arrêt de l'optimiseur,
puisque l'identité rappelée en rubrique~1 ne dépend pas de la convergence~: mesuré, $\bar z$ vaut
$-1{,}7\cdot 10^{-17}$ à $(\hat\delta,\gamma) = (1,\,0)$ et $1{,}6\cdot 10^{-16}$ à $(1,\,3)$,
deux points très éloignés de l'optimum. Lorsque $\hat\pi_t$ varie, $\bar z$ est l'écart entre
centrage pondéré et centrage non pondéré~: l'identité
$\sum_t\sqrt{\hat\pi_t}\,z_t = 0$ donne la borne exacte
\[
  |\bar z| \;\le\; \frac{\max_t\big|\sqrt{\hat\pi_t} - \bar c\big|}{\bar c}
  \cdot \frac{1}{T}\sum_t |z_t| ,
  \qquad \bar c = \frac{1}{T}\sum_t \sqrt{\hat\pi_t} .
\]
$\bar z$ est donc \emph{de l'ordre de} la dispersion des poids, sans en être une fonction~: le
coefficient qui les relie est aléatoire. Mesuré sur les 506 réplications à $\hat\pi_t^{*}$ non
constant, la corrélation entre $|\bar z^{*}|$ et l'amplitude relative de $\hat\pi_t^{*}$ n'est que
de $0{,}26$, et le rapport des deux s'étend de $0{,}006$ à $0{,}12$ entre les quantiles 5~\% et
95~\%. Dans les deux cas, $\bar z$ n'est pas une propriété libre des résidus. Un biais
systématique de niveau non capté par
$\hat\beta$ reste recherché ailleurs, par des vérifications qui portent, elles, un verdict~: les
tests de tendance de l'hypothèse~H1 (section~\ref{sec:h1}) pour une dérive du ratio S/P, les
fiches~\ref{fiche:34} et~\ref{fiche:35} pour une rupture de régime.
\end{fiche}

\begin{fiche}{Variance unitaire des résidus standardisés (diagnostic)}
\refl{Diagnostic ; conditions du premier ordre du maximum de vraisemblance
(fiche~\ref{fiche:31}).}
\Cadre
Résidus standardisés $z_t = \sqrt{\hat\pi_t}\,v_t$ et leur variance empirique
$\operatorname{Var}(z)$. \textbf{Diagnostic, non test.} La condition du premier ordre en $\gamma$
rive l'échelle des $z_t$ comme celle en $\ln\beta$ en rive le niveau~: lorsque $\hat\pi_t$ est
constant ($\hat\delta = 1$, ou volumes $x_t$ constants), les deux conditions donnent ensemble
$\sum_t z_t^2 = T$, soit $\operatorname{Var}(z) = T/(T-1)$ --- c'est-à-dire $1{,}143$ à $T = 8$,
et non~1. À la différence de celle en $\ln\beta$ (fiche~\ref{fiche:31}), cette seconde relation
n'est \emph{pas} une identité algébrique~: elle suppose la dérivée en $\gamma$ effectivement
annulée, donc l'optimisation aboutie. Hors du cas $\hat\pi_t$ constant, seule la contrainte
pondérée subsiste et $\operatorname{Var}(z)$ est de l'ordre de l'écart entre échelle pondérée et
échelle non pondérée. L'argument suppose en outre l'optimum \emph{intérieur} en $\gamma$~: si
$\hat\gamma$ butait sur une borne de son domaine $[-12,\,3]$, la dérivée ne serait pas annulée et
l'égalité tomberait. Cette réserve est théorique~: $\sigma = \hat\beta\,e^{\gamma}$, de sorte que
$\gamma = 3$ signifierait $\sigma \approx 20\,\hat\beta$ et $\gamma = -12$,
$\sigma \approx 6\cdot 10^{-6}\,\hat\beta$, deux situations sans réalité pour un ratio S/P~; sur
le jeu de contrôle, aucune des 999 réplications n'atteint une borne ($\hat\gamma^{*}$ reste entre
$-3{,}34$ et $-1{,}37$).
\Hzero
Sans objet. L'hypothèse $\operatorname{Var}(z_t) = 1$, sous laquelle cette grandeur était
auparavant restituée comme un test, n'est pas celle que l'estimation laisse libre~: la valeur de
référence imposée par les conditions du premier ordre est $T/(T-1)$, et non~1.
\Stat
Aucune statistique de test~; la quantité $C = (T-1)\,s_z^2$ n'est plus calculée. La grandeur
restituée est l'\emph{estimation} $\operatorname{Var}(z)$, portée par le champ \code{estim} de la
table des tests.
\Loi
Sans objet. La loi $\chi^2(T-1)$ de $C$ est nominale~: elle ignore l'estimation préalable de
$(\hat\beta,\hat\delta,\hat\gamma)$ et la contrainte d'échelle qui en résulte. Comme au centrage
(fiche~\ref{fiche:32}), l'argument décisif n'est pas que cette loi serait optimiste, mais qu'à
$\hat\pi_t$ constant elle produit une p-value \textbf{indépendante des données}. Les conditions du
premier ordre donnent alors $\sum_t z_t^2 = T$ et $\bar z = 0$, d'où $C = (T-1)s_z^2
= \sum_t z_t^2 - T\bar z^2 = T$ à la tolérance d'arrêt près, et la p-value bilatérale
$2\min\big(F_{\chi^2(T-1)}(T),\,1 - F_{\chi^2(T-1)}(T)\big)$ ne dépend plus que de $T$~: elle vaut
$0{,}6651878$ à $T = 8$. Cette conséquence est algébrique, mais --- à la différence de celle du
centrage --- elle n'est exacte qu'à la tolérance d'arrêt de l'optimiseur près, puisqu'elle repose
sur la condition en $\gamma$ et non sur la forme fermée de $\ln\hat\beta$ (rubrique~1). Vérifiée
numériquement~: sur le jeu de contrôle comme sur un jeu étranger à volumes constants
(fiche~\ref{fiche:32}), la p-value nominale vaut $0{,}6651888$, les deux valeurs coïncidant
jusqu'à la huitième décimale~; sur les 493 réplications à $\delta^{*} = 1$, elle reste comprise
entre $0{,}665100$ et $0{,}665224$, l'amplitude résiduelle n'étant due qu'à cette même tolérance
d'arrêt, discutée plus bas. Le verdict qui en découlerait serait donc, lui aussi, un \code{OK}
qu'aucune donnée ne peut faire perdre.

La loi simulée sous le modèle ajusté porte, comme celle du centrage (fiche~\ref{fiche:32}), un atome en la valeur
observée. Constat de simulation sur les mêmes $B = 999$ réplications (graine 20260831)~: les 493
à $\hat\pi_t^{*}$ constant donnent une valeur rivée à $T/(T-1)$, d'écart-type
$3{,}5\cdot 10^{-6}$~; les 506 autres une valeur diffuse, d'écart-type $4{,}9\cdot 10^{-3}$.

Une différence avec le centrage mérite d'être relevée, car elle change la nature du bruit
comparé~: ici, ce n'est pas l'arrondi machine mais la \textbf{tolérance d'arrêt de l'optimiseur}.
Sur les 493 réplications concernées, $\operatorname{Var}(z^{*}) - T/(T-1)$ s'étend de
$-3{,}0\cdot 10^{-5}$ à $+7{,}1\cdot 10^{-5}$, de médiane $-7{,}7\cdot 10^{-7}$, quand la valeur
observée vaut $-7{,}8\cdot 10^{-7}$. Or les deux ne sont pas produites au même réglage~:
l'ajustement observé emploie \code{factr = 1e5} et le réajustement des réplications
\code{factr = 1e7}. La p-value de Monte-Carlo autrefois retenue ($0{,}864$) comparait donc deux
tolérances d'optimiseur différentes. Se restreindre à la composante diffuse ne sauverait pas
davantage la vérification qu'au centrage~: la p-value recalculée sur les seules 506 réplications à
$\hat\pi_t^{*}$ non constant vaut $0{,}998$.
\Pval
Aucune. La grandeur a été retirée des statistiques simulées par
\code{.stats\_bootstrapables()}~; aucune p-value n'est calculée ni retenue pour cette ligne.
\Usage
\textbf{Pourquoi aucun verdict.} Comme pour le centrage, la grandeur est rivée par l'estimation~:
elle est restituée comme diagnostic, avec le verdict \code{INFO} et le motif. Sur le jeu de
contrôle ($T = 8$, $\hat\delta = 1$), $\operatorname{Var}(z) = 1{,}142856$ contre
$T/(T-1) = 1{,}142857$~: l'écart, $-6{,}8\cdot 10^{-7}$ en relatif, provient de la tolérance
d'arrêt de l'optimiseur et non des données.

\textbf{Ce que ce diagnostic apporte, et qu'aucun autre n'apporte.} L'écart de $\sum_t z_t^2$ à
$T$ est, à ce jour, le \emph{seul} indicateur de la convergence en $\gamma$ de toute la
restitution~: la condition en $\ln\beta$ étant une identité (fiche~\ref{fiche:31}), elle ne peut
rien signaler, tandis que celle en $\gamma$ n'est satisfaite qu'à l'optimum. Les ordres de
grandeur mesurés situent ce que « notable » veut dire~: à l'optimum, $10^{-6}$ en relatif avec
\code{factr = 1e5}, et jusqu'à $6\cdot 10^{-5}$ sur les réplications réajustées avec
\code{factr = 1e7}~; à $\hat\delta = 1$ mais $\gamma$ forcé hors de l'optimum,
$\sum_t z_t^2 - T$ vaut $-7{,}8$ pour $\gamma = 0$ et $-7{,}97$ pour $\gamma = 3$. Un écart de cet
ordre est donc un défaut de convergence, non une propriété des données. Lorsque $\hat\pi_t$ varie,
l'écart est en revanche de l'ordre de la dispersion des poids, sans en être une fonction, et ne
se lit plus comme un contrôle numérique. Le contrôle de l'échelle de $\sigma$ proprement dit
relève de l'intervalle bootstrap et du jackknife (section~\ref{sec:robustesse}), qui portent, eux,
sur $\sigma_{\mathrm{USP}}$.
\end{fiche}

\newpage
\clearpage
% =============================================================================
\section{Diagnostics de stabilité, ruptures et points aberrants}
\label{sec:stabilite}
% =============================================================================

Ces diagnostics ne correspondent pas directement à l'une des quatre hypothèses réglementaires,
mais conditionnent la validité pratique des tests précédents~: un point aberrant ou une rupture de
régime peut à lui seul provoquer ou masquer un rejet ailleurs.

\begin{fiche}{Test de rupture sup-F (Quandt / Chow / Andrews)}
\refl{R.\,E. Quandt, \emph{Tests of the hypothesis that a linear regression system obeys two
separate regimes}, \emph{JASA}, 55(290), 1960, p.~324--330 ; G.\,C. Chow, \emph{Tests of equality
between sets of coefficients in two linear regressions}, \emph{Econometrica}, 28(3), 1960,
p.~591--605 ; loi limite~: D.\,W.\,K. Andrews, \emph{Tests for parameter instability and
structural change with unknown change point}, \emph{Econometrica}, 61(4), 1993, p.~821--856.}
\Cadre
Série $z_1,\dots,z_T$ et ensemble des points de rupture candidats
$\mathcal{K} = \{k : \max(2, \lfloor \tau T\rfloor) \le k \le \min(T-2,\, T - \lfloor \tau T\rfloor)\}$
avec un écrêtement $\tau = 0{,}15$ (indispensable~: sans lui la statistique diverge aux
extrémités). Les bornes $2$ et $T-2$ garantissent au moins deux observations de chaque côté,
sans quoi la variance intra-segment n'est pas définie.
Hypothèses non testées~: erreurs indépendantes, homoscédastiques, gaussiennes.
\Hzero
$H_0$ : absence de rupture de niveau, soit $\mathbb{E}[z_t]$ constant sur toute la période.
\Stat
Pour un point de rupture $k$ fixé, en notant $\bar z_1^{(k)}$ et $\bar z_2^{(k)}$ les moyennes
avant et après~:
\[
  \mathrm{SCR}(k) = \sum_{t\le k}\big(z_t - \bar z_1^{(k)}\big)^2 + \sum_{t> k}\big(z_t - \bar z_2^{(k)}\big)^2,
  \qquad
  F(k) = \frac{\mathrm{SCT} - \mathrm{SCR}(k)}{\mathrm{SCR}(k)/(T-2)},
\]
puis $\displaystyle \sup F = \max_{k\in\mathcal{K}} F(k)$.
\Loi
\textbf{Pas} $\mathcal{F}(1,T-2)$~: le point de rupture étant \emph{estimé} et non fixé a priori,
la loi de $\sup F$ est celle d'un supremum de processus, tabulée par Andrews (1993) et dépendant
de $\tau$. Ses quantiles sont sensiblement supérieurs à ceux de la loi de Fisher.
\Pval
Approximation analytique de B.\,E. Hansen (\emph{Approximate asymptotic p values for
structural-change tests}, \emph{Journal of Business \& Economic Statistics}, 15(1), 1997,
p.~60--67). Le script utilise la p-value de Monte-Carlo.
\Usage
Erreur classique à éviter~: comparer $\sup F$ aux quantiles de $\mathcal{F}(1,T-2)$, ce qui
surestime massivement le nombre de ruptures détectées, le maximum sur une dizaine de candidats
étant mécaniquement grand sous $H_0$. La simulation règle ce problème. Un rejet est très
informatif~: il signale un changement de politique de souscription, de provisionnement ou de
périmètre, et peut justifier de restreindre la profondeur $T$ retenue.
\end{fiche}

\begin{fiche}{Test de stabilité OLS-CUSUM}
\refl{R.\,L. Brown, J.~Durbin, J.\,M. Evans, \emph{Techniques for testing the constancy of
regression relationships over time}, \emph{Journal of the Royal Statistical Society, Series B},
37(2), 1975, p.~149--192 ; version fondée sur les résidus MCO~: W.~Ploberger, W.~Krämer,
\emph{The CUSUM test with OLS residuals}, \emph{Econometrica}, 60(2), 1992, p.~271--285.}
\Cadre
Série de résidus $z_t$, $\hat\sigma$ leur écart-type empirique. Hypothèses non testées~:
indépendance et homoscédasticité sous $H_0$.
\Hzero
$H_0$ : constance des paramètres du modèle sur toute la période.
\Stat
\[
  S(j) = \frac{1}{\hat\sigma\sqrt{T}}\sum_{t=1}^{j}\big(z_t - \bar z\big),
  \qquad
  \mathrm{CUSUM} = \max_{1\le j\le T}\big|S(j)\big| .
\]
\Loi
Sous $H_0$, $S(\lfloor uT \rfloor) \Rightarrow B^{\circ}(u)$, pont brownien standard sur $[0,1]$.
La statistique converge donc vers $\sup_{u\in[0,1]}|B^\circ(u)|$, de loi de Kolmogorov.
\Pval
Asymptotiquement
$p = 2\sum_{j=1}^{\infty}(-1)^{j-1}\exp\!\big(-2j^2\,\mathrm{CUSUM}^2\big)$.
Le script retient la p-value de Monte-Carlo, l'approximation brownienne étant inutilisable à
$T \le 15$.
\Usage
Complémentaire du sup-F~: le CUSUM détecte plutôt des dérives \emph{graduelles}, là où le sup-F
détecte des ruptures \emph{brutales}. Sa limite majeure ici est que la convergence vers le pont
brownien suppose $T$ grand~: sur 10 points, le processus cumulé n'a aucune chance de ressembler à
une trajectoire brownienne et seule la simulation donne un niveau correct.
\end{fiche}

\begin{fiche}{Test de Grubbs}
\refl{F.\,E. Grubbs, \emph{Sample criteria for testing outlying observations}, \emph{Annals of
Mathematical Statistics}, 21(1), 1950, p.~27--58 ; \emph{Procedures for detecting outlying
observations in samples}, \emph{Technometrics}, 11(1), 1969, p.~1--21.}
\Cadre
Échantillon $z_1,\dots,z_T$, $T \ge 3$. Hypothèse non testée \emph{et} déterminante~: sous $H_0$,
les observations sont i.i.d. \emph{normales}.
\Hzero
$H_0$ : toutes les observations proviennent de la même loi normale (aucune valeur aberrante).
L'alternative est la présence d'\emph{exactement une} valeur aberrante.
\Stat
\[
  G \;=\; \frac{\max_{t}\big|z_t - \bar z\big|}{s_z}.
\]
\Loi
Sous $H_0$, la loi exacte de $G$ se déduit de celle de Student~: en posant
\[
  t^\star = \sqrt{\frac{T(T-2)\,G^2}{(T-1)^2 - T\,G^2}},
\]
on a $t^\star \sim t(T-2)$ pour une observation fixée~; $G$ étant un maximum sur $T$ observations,
la loi de référence intègre une correction de multiplicité.
\Pval
Borne de Bonferroni utilisée par le script, tronquée à 1~:
\[
  p = \min\Big(1,\; 2T\big(1 - F_{t,T-2}(t^\star)\big)\Big).
\]
\Usage
Deux limites structurelles. (i) \emph{Masquage}~: en présence de deux valeurs aberrantes ou plus,
celles-ci gonflent $s_z$ et le test ne détecte plus rien, d'où la nécessité du test de Rosner.
(ii) La statistique est bornée~: $G \le (T-1)/\sqrt{T}$, soit $2{,}85$ pour $T=10$~; sur de très
petits échantillons, aucune observation ne peut être « suffisamment » extrême. Enfin, la
sensibilité à l'hypothèse de normalité est totale~: sur une vraie loi à queue lourde, le test
signale des aberrations qui n'en sont pas.
\end{fiche}

\begin{fiche}{Procédure ESD généralisée de Rosner}
\refl{B.~Rosner, \emph{Percentage points for a generalized ESD many-outlier procedure},
\emph{Technometrics}, 25(2), 1983, p.~165--172.}
\Cadre
Généralisation itérative de Grubbs. On fixe a priori un nombre maximal $k$ de valeurs aberrantes
recherchées (le script prend $k = \max(1, \lfloor T/4 \rfloor)$) et un niveau $\alpha$, qui est
celui passé à \code{usp\_run()} (0,10 par défaut) et non une valeur figée. Hypothèse non
testée~: normalité de la population sous-jacente.
\Hzero
$H_0$ : aucune valeur aberrante. On teste séquentiellement les alternatives « il y a exactement
$i$ valeurs aberrantes », $i=1,\dots,k$.
\Stat
Pour $i = 1,\dots,k$, sur l'échantillon $\mathcal{S}_i$ obtenu en retirant les $i-1$ observations
déjà écartées~:
\[
  R_i = \frac{\max_{z\in\mathcal{S}_i}\big|z - \bar z_{\mathcal{S}_i}\big|}{s_{\mathcal{S}_i}} .
\]
\Loi
Valeurs critiques exactes, avec $n_i = T-i+1$ et
$t_{\text{crit}} = F^{-1}_{t,\,n_i-2}\big(1 - \tfrac{\alpha}{2 n_i}\big)$~:
\[
  \lambda_i \;=\; \frac{(n_i - 1)\, t_{\text{crit}}}{\sqrt{\big(n_i - 2 + t_{\text{crit}}^2\big)\, n_i}} .
\]
\Pval
\textbf{Aucune p-value}~: il s'agit d'une procédure de décision à niveau $\alpha$ fixé. Le nombre
retenu de valeurs aberrantes est $\hat k = \max\{i : R_i > \lambda_i\}$, et $0$ si aucun
dépassement.
\Usage
Résout le problème de masquage de Grubbs et constitue la référence pour la détection multiple.
Limites~: la littérature recommande $T \ge 25$ pour une fiabilité correcte des valeurs critiques,
ce qui n'est jamais atteint en régime USP~; le choix de $k$ est arbitraire et influence le
résultat~; l'hypothèse de normalité reste critique. À $T \le 15$, la procédure est un simple
signal d'attention, non une justification d'exclusion de données -- laquelle doit toujours reposer
sur une analyse métier (sinistre exceptionnel identifié, erreur de saisie avérée).
\end{fiche}

\begin{fiche}{Distance de Cook (diagnostic)}
\refl{R.\,D. Cook, \emph{Detection of influential observation in linear regression},
\emph{Technometrics}, 19(1), 1977, p.~15--18.}
\Cadre
Régression contrainte $y_t = \beta x_t + \varepsilon_t$ à $k$ paramètres ($k=1$), de résidus
$e_t$, de variance résiduelle estimée $\hat s^2$ et de leviers $h_t$ (fiche suivante).
\textbf{Diagnostic d'influence, non test.}
\Hzero
Sans objet.
\Stat
\[
  D_t \;=\; \frac{e_t^2}{k\,\hat s^2}\cdot\frac{h_t}{(1-h_t)^2}
  \;=\; \frac{\sum_{i}\big(\hat y_i - \hat y_{i(-t)}\big)^2}{k\,\hat s^2},
\]
où $\hat y_{i(-t)}$ est la prévision obtenue en retirant l'observation $t$.
\Loi
Aucune loi exacte. On rapproche parfois $D_t$ des quantiles de $\mathcal{F}(k, T-k)$, mais il
s'agit d'une heuristique sans fondement rigoureux.
\Pval
Aucune. Le script applique le seuil conventionnel $D_t > 4/T$.
\Usage
Mesure combinée du levier et de l'ampleur du résidu~: elle répond à la question « de combien le
modèle changerait-il si je retirais cette année ? ». C'est le diagnostic le plus directement utile
en régime USP, où une seule année peut déterminer le paramètre. Limites~: le seuil $4/T$ est
purement conventionnel (d'autres auteurs retiennent $1$), et la distance porte sur la régression
auxiliaire, non sur le modèle de vraisemblance complet -- le jackknife de la section~10 fournit la
mesure exacte pour ce dernier.
\end{fiche}

\begin{fiche}{Leviers (\emph{hat values}, diagnostic)}
\refl{D.\,C. Hoaglin, R.\,E. Welsch, \emph{The hat matrix in regression and ANOVA},
\emph{The American Statistician}, 32(1), 1978, p.~17--22.}
\Cadre
Matrice de projection $H = X(X'X)^{-1}X'$ de la régression contrainte. \textbf{Diagnostic, non
test.}
\Hzero
Sans objet.
\Stat
$h_t = H_{tt}$, soit ici, en régression simple sans constante,
$\displaystyle h_t = \frac{x_t^2}{\sum_{i=1}^{T} x_i^2}$, avec $\sum_t h_t = k$ et
$0 \le h_t \le 1$.
\Loi
Aucune.
\Pval
Aucune. Seuil conventionnel~: $h_t > 2k/T$.
\Usage
Le levier ne dépend que des $x_t$~: il identifie les années dont le \emph{volume} est atypique,
indépendamment de la sinistralité observée. Une année à fort levier n'est pas problématique en
soi, mais elle le devient si son résidu est également grand -- c'est exactement ce que combine la
distance de Cook. En contexte USP, un levier élevé signale généralement une année de forte
croissance ou une intégration de portefeuille, à rapprocher du contrôle d'amplitude du volume de
la section~3.
\end{fiche}

\newpage
\clearpage
% =============================================================================
\section{Diagnostics de robustesse de l'estimation}
\label{sec:robustesse}
% =============================================================================

\begin{fiche}{Test du rapport de vraisemblance sur les cas limites de $\delta$}
\refl{S.\,S. Wilks, \emph{The large-sample distribution of the likelihood ratio for testing
composite hypotheses}, \emph{Annals of Mathematical Statistics}, 9(1), 1938, p.~60--62 ;
correction de frontière~: H.~Chernoff, \emph{On the distribution of the likelihood ratio},
\emph{Annals of Mathematical Statistics}, 25(3), 1954, p.~573--578 ; D.\,O. Stram, J.\,W. Lee,
\emph{Variance components testing in the longitudinal mixed effects model}, \emph{Biometrics},
50(4), 1994, p.~1171--1177.}
\Cadre
On note $\ell(\delta,\gamma)$ la log-vraisemblance et $\mathcal{O}(\delta,\gamma) = -2\ell$ à une
constante près, soit la fonction objectif minimisée par le script. Hypothèses non testées~:
régularité du modèle, identifiabilité, échantillon suffisant pour la théorie asymptotique.
\Hzero
$H_0 : \delta = \delta_0$ avec $\delta_0 \in \{0, 1\}$, c'est-à-dire structure de variance purement
linéaire en volume ($\delta_0 = 0$) ou purement quadratique ($\delta_0 = 1$).
\Stat
Statistique du rapport de vraisemblance, obtenue par profilage sur $\gamma$~:
\[
  \mathrm{LR}(\delta_0) \;=\; \min_{\gamma}\mathcal{O}(\delta_0,\gamma)
                        \;-\; \min_{\delta\in[0,1],\,\gamma}\mathcal{O}(\delta,\gamma) .
\]
\Loi
Sous les conditions de régularité de Wilks,
$\mathrm{LR} \xrightarrow{\ \mathcal{L}\ } \chi^2(1)$.
\textbf{Mais $\delta_0 = 0$ et $\delta_0 = 1$ sont sur la frontière du domaine $[0,1]$}~: la
condition de régularité est violée et le résultat de Chernoff (1954) s'applique, la loi limite
devenant le mélange
\[
  \tfrac12\,\chi^2(0) + \tfrac12\,\chi^2(1),
\]
où $\chi^2(0)$ désigne la masse de Dirac en~0.
\Pval
p-value naïve rapportée par le script~: $p = 1 - F_{\chi^2_1}(\mathrm{LR})$.
p-value corrigée pour la frontière~:
$p_{\text{Chernoff}} = \tfrac12\big(1 - F_{\chi^2_1}(\mathrm{LR})\big) = p/2$.
\Usage
\textbf{Point d'attention majeur.} La p-value affichée est \emph{conservatrice d'un facteur~2}
lorsque $\hat\delta$ est au bord~: la vraie p-value est la moitié de celle rapportée. Le script
affiche volontairement la version naïve, plus prudente, mais la note de justification doit
mentionner la correction. Sur le fond, ce test répond à la question pratique décisive~: la
structure de mélange imposée par le règlement est-elle réellement identifiée par les données, ou
un cas limite serait-il statistiquement aussi défendable ? Une $\mathrm{LR}$ faible pour les deux
bornes signifie que $\delta$ n'est pas identifié et que le paramètre final repose sur une
convention plutôt que sur l'information contenue dans l'historique.
\end{fiche}

\begin{fiche}{Jackknife (sensibilité au retrait d'une année)}
\refl{M.\,H. Quenouille, \emph{Approximate tests of correlation in time series}, \emph{Journal of
the Royal Statistical Society, Series B}, 11(1), 1949, p.~68--84 ; J.\,W. Tukey, \emph{Bias and
confidence in not quite large samples}, \emph{Annals of Mathematical Statistics}, 29(2), 1958,
p.~614.}
\Cadre
Pour chaque $t$, on réestime intégralement le modèle sur l'échantillon privé de l'observation $t$,
d'où $\hat\sigma^{(-t)}$ et $\sigma_{\mathrm{USP}}^{(-t)}$. Le facteur de crédibilité $c$ et la
correction $\sqrt{(T+1)/(T-1)}$ sont maintenus à ceux de l'échantillon complet, l'objectif étant
d'isoler l'effet de l'observation retirée et non de rejouer les règles de l'annexe~XVII sur $T-1$
années. \textbf{Diagnostic, non test.}
\Hzero
Sans objet.
\Stat
\[
  \Delta_{\max} \;=\; \max_{1\le t\le T}
  \frac{\big|\sigma_{\mathrm{USP}}^{(-t)} - \sigma_{\mathrm{USP}}\big|}{\sigma_{\mathrm{USP}}} .
\]
L'estimateur jackknife du biais,
$\widehat{\mathrm{biais}} = (T-1)\big(\overline{\sigma^{(-\cdot)}} - \hat\sigma\big)$, est
également calculable à partir de la table produite.
\Loi
Aucune loi de référence n'est utilisée pour $\Delta_{\max}$.
\Pval
Aucune. Seuils conventionnels du script~: \code{ALERTE} au-delà de 10~\%, \code{ECHEC} au-delà de
20~\%.
\Usage
C'est, à $T$ faible, le diagnostic le plus informatif de tout le dispositif. Il répond directement
à la question que posera le superviseur~: le paramètre proposé résiste-t-il au retrait d'une seule
année ? Un $\Delta_{\max}$ de 25~\% signifie que le calibrage est déterminé par une observation
unique, ce qui fragilise le dossier bien davantage qu'un test de normalité en échec. Limite~: le
jackknife suppose une influence approximativement linéaire de chaque observation, hypothèse
discutable en présence d'une solution au bord sur $\delta$, où le retrait d'un point peut faire
basculer l'optimum d'une borne à l'autre de façon discontinue.
\end{fiche}

\begin{fiche}{Intervalle de confiance par bootstrap paramétrique}
\refl{Efron (1979), \emph{op.\ cit.} ; B.~Efron, R.\,J. Tibshirani, \emph{An Introduction to the
Bootstrap}, Chapman \& Hall, 1993.}
\Cadre
On réutilise les $B$ réplications décrites en section~\ref{sec:mc}, dont on extrait
$\sigma_{\mathrm{USP}}^{(b)} = c\,\hat\sigma^{(b)}\sqrt{(T+1)/(T-1)} + (1-c)\,\sigma_{\text{std}}$.
\textbf{Diagnostic, non test.}
\Hzero
Sans objet (procédure d'estimation d'incertitude, non de décision).
\Stat
Quantiles empiriques de la distribution bootstrap~:
$\big[\hat q_{\alpha/2},\, \hat q_{1-\alpha/2}\big]$ (intervalle de percentile), et largeur
relative $L = \big(\hat q_{0{,}95} - \hat q_{0{,}05}\big)/\sigma_{\mathrm{USP}}$.
\Loi
La distribution bootstrap approxime la loi d'échantillonnage de $\hat\sigma_{\mathrm{USP}}$
conditionnellement au modèle ajusté.
\Pval
Aucune. Seuils conventionnels du script sur $L$~: \code{ALERTE} au-delà de 50~\%, \code{ECHEC}
au-delà de 80~\%.
\Usage
Quantifie l'incertitude d'échantillonnage du paramètre final, ce que ne fait aucun test
d'hypothèse. Il est fréquent d'obtenir, sur un historique de 10 ans, un intervalle à 90~\% dont la
largeur dépasse 50~\% de la valeur centrale~: cela ne disqualifie pas le calibrage (le mélange par
crédibilité en amortit une partie) mais doit être présenté explicitement. Limites~: l'intervalle
de percentile n'est corrigé ni du biais ni de l'asymétrie -- les variantes BCa seraient plus
exactes -- et il reste conditionnel à la validité du modèle paramétrique, donc optimiste si le
modèle est mal spécifié.
\end{fiche}

\begin{fiche}{Convergence multi-démarrages}
\refl{Diagnostic numérique ; sur l'optimisation sous contraintes de bornes~: R.\,H. Byrd, P.~Lu,
J.~Nocedal, C.~Zhu, \emph{A limited memory algorithm for bound constrained optimization},
\emph{SIAM Journal on Scientific Computing}, 16(5), 1995, p.~1190--1208 (algorithme L-BFGS-B).}
\Cadre
L'optimisation de $\mathcal{O}(\delta,\gamma)$ est relancée depuis une grille de points de départ
($9 \times 6 = 54$ démarrages dans la configuration par défaut). \textbf{Diagnostic, non test.}
\Hzero
Sans objet.
\Stat
Proportion de démarrages atteignant le minimum global à $10^{-6}$ près~:
\[
  \hat\kappa = \frac{1}{S}\sum_{s=1}^{S} \mathbb{1}\Big\{\big|\mathcal{O}_s^{\star} -
  \min_{s'}\mathcal{O}_{s'}^{\star}\big| < 10^{-6}\Big\},
\]
complétée par le code de retour de l'optimiseur.
\Loi
Sans objet.
\Pval
Aucune. \code{ALERTE} si $\hat\kappa < 0{,}5$ ou si le code de retour est non nul.
\Usage
Garantit que le minimum rapporté est global et non local, ce qui n'a rien d'automatique~: la
fonction objectif de l'annexe~XVII peut présenter un plateau très plat en $\delta$ lorsque
l'amplitude des volumes $x_t$ est faible (car $\pi_t$ dépend alors peu de $\delta$), situation
dans laquelle l'optimiseur s'arrête en des points différents selon son initialisation. Un
$\hat\kappa$ faible est un signal d'alerte à traiter \emph{avant} toute interprétation des tests,
au même titre que l'écart de $\sum_t z_t^2$ à $T$ (fiche~\ref{fiche:33})~; la nullité de la
moyenne pondérée des résidus, elle, est une identité et ne signale rien (fiche~\ref{fiche:31}).
\end{fiche}

\newpage
\clearpage
% =============================================================================
\section{Diagnostics graphiques complémentaires}
% =============================================================================

Les graphiques ne comportent ni hypothèse nulle ni p-value. Toutes les quantités
qu'ils représentent sont calculées dans \code{R/engine.R}
(\code{engine\_plots\_data()} et \code{mw\_plots\_data()}) et exposées dans
\code{res\$plots\_data}~; la couche d'affichage ne fait que les tracer. Ils sont
documentés ici car, aux profondeurs réglementaires, leur pouvoir informatif
dépasse souvent celui des tests formels.

\subsection{Couverture exhaustive}

Les vingt-six représentations produites par l'outil sont recensées ci-dessous.
Celles qui appellent un commentaire méthodologique font l'objet d'une fiche~; les
autres sont d'interprétation directe et sont décrites dans le tableau.

\begin{center}
\small
\rowcolors{2}{tablerow}{white}
\begin{tabular}{@{}p{4.7cm}p{1.3cm}p{1.9cm}p{6.3cm}@{}}
\rowcolor{tableheader}
\textcolor{financeblue}{\textbf{Fonction}} &
\textcolor{financeblue}{\textbf{Méthode}} &
\textcolor{financeblue}{\textbf{Rattachement}} &
\textcolor{financeblue}{\textbf{Objet}} \\
\multicolumn{4}{@{}l@{}}{\textit{\textcolor{financeblue}{Données et ajustement}}} \\
\code{plot\_ajustement()} & LN & --- & Nuage $(x_t, y_t)$ et droite
$\hat\beta x$~: contrôle visuel direct de la proportionnalité (H1). \\
\code{plot\_ratio()} & LN & H1 & Ratio $r_t$ dans le temps, avec le niveau
$\hat\beta$~: une dérive contredit la constance de $\beta$. \\
\code{plot\_residus()} & LN & H2 & Résidus $z_t$ contre le volume~: structure
résiduelle après pondération. \\
\code{plot\_mw\_facteurs()} & MW & --- & Facteurs $\hat f_j$ par année de
développement, avec la référence $f = 1$~: visualise la cadence de liquidation. \\
\code{plot\_mw\_reserve()} & MW & M6 & Réserve par année de survenance~: montre
la concentration sur les années récentes. \\
\multicolumn{4}{@{}l@{}}{\textit{\textcolor{financeblue}{Structure de variance}}} \\
\code{plot\_spread()} & LN & H2 & Fiche dédiée. \\
\code{plot\_qq2ech()} & LN & H2 & Fiche dédiée. \\
\code{plot\_mw\_regressions()} & MW & M1 & Fiche dédiée. Petits multiples
$C(i,j+1)$ contre $C(i,j)$, une vignette par année de développement. \\
\code{plot\_mw\_origine()} & MW & M1 & Fiche dédiée. Ordonnée à l'origine
estimée par colonne, avec son intervalle à 95~\%. \\
\code{plot\_mw\_alpha()} & MW & M1 & Fiche dédiée. Facteur estimé selon
$\alpha = 0$, 1 et 2. \\
\code{plot\_mw\_residus\_C()} & MW & M1, M2 & Résidus de Mack contre $C(i,j)$~:
teste conjointement la proportionnalité de la moyenne et celle de la variance. \\
\code{plot\_mw\_residus\_dev()} & MW & M2 & Résidus par année de développement~:
localise la colonne dont la variance est mal spécifiée. \\
\multicolumn{4}{@{}l@{}}{\textit{\textcolor{financeblue}{Normalité}}} \\
\code{plot\_qqnorm()} & LN & H3 & Fiche dédiée. \\
\code{plot\_mw\_qq()} & MW & M5 & QQ-plot des résidus de Mack. Diagnostic
seulement~: la normalité n'est pas une hypothèse du modèle. \\
\multicolumn{4}{@{}l@{}}{\textit{\textcolor{financeblue}{Indépendance}}} \\
\code{plot\_mw\_residus\_acc()} & MW & M3 & Résidus par année de survenance~:
détecte une ligne atypique du triangle. \\
\code{plot\_mw\_residus\_cal()} & MW & M3 & Résidus par année calendaire~:
contrepartie graphique du test des effets calendaires, la plus directement
lisible des trois vues de résidus. \\
\end{tabular}
\end{center}

\noindent
Suite du recensement~: diagnostics d'influence, estimation et calibration.

\begin{center}
\small
\rowcolors{2}{tablerow}{white}
\begin{tabular}{@{}p{4.7cm}p{1.3cm}p{1.9cm}p{6.3cm}@{}}
\rowcolor{tableheader}
\textcolor{financeblue}{\textbf{Fonction}} &
\textcolor{financeblue}{\textbf{Méthode}} &
\textcolor{financeblue}{\textbf{Rattachement}} &
\textcolor{financeblue}{\textbf{Objet}} \\
\multicolumn{4}{@{}l@{}}{\textit{\textcolor{financeblue}{Influence et leviers}}} \\
\code{plot\_influence\_levier()} & LN & --- & Fiche dédiée. \\
\code{plot\_influence\_cook()} & LN & --- & Fiche dédiée. \\
\code{plot\_influence\_sigma()} & LN & --- & Fiche dédiée. \\
\code{plot\_mw\_levier()} & MW & --- & Fiche dédiée. \\
\code{plot\_mw\_dfbeta()} & MW & --- & Fiche dédiée. \\
\code{plot\_mw\_contributions()} & MW & M6 & Fiche dédiée. \\
\multicolumn{4}{@{}l@{}}{\textit{\textcolor{financeblue}{Estimation et calibration}}} \\
\code{plot\_profil\_delta()} & LN & --- & Profil de vraisemblance en $\delta$~:
un plateau signale une structure de variance non identifiée. \\
\code{plot\_surface\_objectif()} & LN & --- & Surface $-2\ln L$ sur la grille
$(\delta,\gamma)$~: montre si l'optimum au bord est un vrai minimum. \\
\code{plot\_coupe\_delta()} & LN & --- & Coupe de la surface à $\hat\gamma$~:
lecture bidimensionnelle de la même platitude, avec son amplitude chiffrée. \\
\code{plot\_boot\_sigma()} & LN + MW & --- & Distribution bootstrap de
$\hat\sigma$~: matérialise l'incertitude d'échantillonnage du paramètre. \\
\code{plot\_calibration()} & LN + MW & --- & Valeurs candidates et paramètre
retenu~: situe $\sigma_{\mathrm{USP}}$ entre le standard et l'estimé seul. \\
\end{tabular}
\end{center}

\subsection{Fiches détaillées}

\begin{fiche}{QQ-plot normal}
\refl{M.\,B. Wilk, R.~Gnanadesikan, \emph{Probability plotting methods for the analysis of data},
\emph{Biometrika}, 55(1), 1968, p.~1--17.}
\Cadre
Nuage des couples $\big(\Phi^{-1}\!\big((i-\tfrac38)/(T+\tfrac14)\big),\, z_{(i)}\big)$,
$i=1,\dots,T$, avec droite de référence passant par les premier et troisième quartiles. Rattaché
à l'hypothèse H3.
\Hzero
Sans objet. Le test numérique correspondant est Shapiro--Francia, dont la statistique $W'$ est
exactement le carré du coefficient de corrélation de ce nuage.
\Stat
Sans objet (représentation graphique).
\Loi
Sans objet.
\Pval
Sans objet.
\Usage
Lecture~: un alignement sur la droite conforte la normalité~; une courbure en S traduit des queues
plus lourdes (points extrêmes au-dessus à droite et en dessous à gauche) ou plus légères~; une
courbure convexe ou concave traduit une asymétrie. Limite~: sur 8 à 12 points, l'œil surestime
facilement les écarts -- il est utile de superposer une enveloppe de simulation obtenue par
bootstrap paramétrique pour calibrer la lecture.
\end{fiche}

\begin{fiche}{QQ-plot à deux échantillons (faible volume contre fort volume)}
\refl{Wilk \& Gnanadesikan (1968), \emph{op.\ cit.}}
\Cadre
Les résidus sont scindés en deux groupes selon $x_t$ inférieur ou supérieur à
$\operatorname{med}(x)$. On porte en abscisse les quantiles empiriques du groupe « faible volume »
et en ordonnée ceux du groupe « fort volume », évalués aux mêmes probabilités avec
$n = \min(n_1,n_2)$. Droite de référence~: la première bissectrice. Rattaché à l'hypothèse H2.
\Hzero
Sans objet. Le test numérique correspondant est le test de Smirnov à deux échantillons
(section~5).
\Stat
Sans objet.
\Loi
Sans objet.
\Pval
Sans objet.
\Usage
Diagnostic visuel direct de l'hypothèse de structure de variance~: si la pondération $\pi_t$ est
correcte, les résidus normalisés des petits et des gros volumes suivent la même loi et les points
s'alignent sur la bissectrice. Une \emph{pente} supérieure à 1 signale une dispersion résiduelle
plus forte sur les gros volumes (donc un $\delta$ sous-estimé), une pente inférieure à 1
l'inverse~; un décalage vertical d'ensemble signale une différence de niveau relevant plutôt de
H1. Limite~: avec 5 points par groupe, la pente apparente est très instable.
\end{fiche}

\begin{fiche}{Graphiques d'influence et de levier}
\refl{R.\,D. Cook (1977), \emph{op.\ cit.} ; D.\,C. Hoaglin, R.\,E. Welsch (1978),
\emph{op.\ cit.} ; pour la représentation conjointe, J.\,M. Chambers,
W.\,S. Cleveland, B.~Kleiner, P.\,A. Tukey, \emph{Graphical Methods for Data
Analysis}, Wadsworth, 1983.}
\Cadre
Trois représentations, rattachées aux diagnostics d'influence de la
section~\ref{sec:robustesse}. Pour la méthode lognormale, la régression de
référence est le modèle réglementaire contraint $y_t = \beta x_t$, à $k = 1$
paramètre, de leviers $h_t = x_t^2 / \sum_i x_i^2$ (de somme $k$) et de résidus
standardisés $e_t / (s\sqrt{1-h_t})$.
\Hzero
Sans objet. Les tests numériques correspondants sont la distance de Cook et les
leviers, documentés en section~\ref{sec:robustesse}.
\Stat
Sans objet (représentations graphiques). Les iso-distances de Cook tracées dans
le plan $(h, r)$ sont les courbes
\[
  r = \pm\sqrt{\frac{D\,k\,(1-h)}{h}}, \qquad D \in \{0{,}5\,;\,1\},
\]
qui découlent directement de l'identité $D_t = r_t^2\,h_t / \big(k(1-h_t)\big)$,
vérifiée numériquement à $10^{-10}$ près par le moteur.
\Loi
Sans objet.
\Pval
Sans objet.
\Usage
Le premier graphique croise levier et résidu standardisé~: il sépare visuellement
les observations \emph{atypiques en volume} (levier élevé, résidu faible) de
celles qui sont \emph{réellement influentes} (au-delà d'une iso-courbe de Cook).
Un levier élevé seul n'est pas problématique~; c'est sa combinaison avec un
résidu important qui l'est. Le deuxième donne la distance de Cook par année avec
son seuil $4/T$. Le troisième, propre au projet, est le plus directement
exploitable en dossier~: il représente l'écart relatif sur $\sigma_{\mathrm{USP}}$
induit par le retrait de chaque année, issu du jackknife. À $T = 8$, il n'est pas
rare qu'une seule année déplace le paramètre de plus de 10~\%, ce que la seule
lecture des distances de Cook ne montre pas.

\medskip
\textbf{Transposition à la méthode Merz--Wüthrich.} Le facteur $\hat f_j$ étant
une moyenne pondérée des facteurs individuels de poids $C(i,j)$, le levier exact
de la cellule $(i,j)$ vaut $h(i,j) = C(i,j)/S_j$, de somme~1 par colonne
(vérifié numériquement). L'influence exacte se mesure par le DFBETA
\[
  \hat f_j^{(-i)} = \frac{\sum_{i'} C(i',j+1) - C(i,j+1)}{\sum_{i'} C(i',j) - C(i,j)},
\]
obtenu sans réajustement itératif. S'y ajoute la décomposition de la réserve et de
la variance de processus de la MSEP par année de survenance. Cette dernière est une
\emph{décomposition} et non un leave-one-out~: le terme d'erreur d'estimation de la
MSEP comporte des termes croisés, partagés entre paires d'années, qui ne sont pas
attribuables à une année isolée, ce que le graphique ne prétend pas faire.
\end{fiche}

\begin{fiche}{Graphique dispersion-niveau (\emph{spread-location})}
\refl{J.\,W. Tukey, \emph{Exploratory Data Analysis}, Addison-Wesley, 1977 ; lissage local de
W.\,S. Cleveland, \emph{Robust locally weighted regression and smoothing scatterplots},
\emph{JASA}, 74(368), 1979, p.~829--836.}
\Cadre
Nuage $\big(x_t,\ \sqrt{|z_t|}\big)$ avec courbe de lissage local (\emph{lowess}). Rattaché à
l'hypothèse H2.
\Hzero
Sans objet.
\Stat
Sans objet. La transformation $\sqrt{|\cdot|}$ stabilise la variance de la mesure de dispersion et
rend la courbe plus lisible que $|z_t|$ ou $z_t^2$.
\Loi
Sans objet.
\Pval
Sans objet.
\Usage
Une courbe lissée horizontale confirme que la pondération capture l'hétéroscédasticité~; une pente
marquée signale une structure de variance résiduelle non modélisée, en cohérence avec un rejet de
Breusch--Pagan ou de White. Limite~: le lissage local n'a aucune stabilité en-dessous d'une
dizaine de points et sa forme dépend fortement du paramètre de fenêtre~; à $T$ faible, seule une
pente très prononcée doit être retenue.
\end{fiche}

\begin{fiche}{Visualisation de la régression chain-ladder (Merz--Wüthrich)}
\label{sec:graph-mw}
\refl{T.~Mack (1993), \emph{op.\ cit.}, section~3, qui recommande explicitement
l'examen graphique des nuages $C(i,j+1)$ contre $C(i,j)$ avant tout test formel.}
\Cadre
Trois représentations, rattachées à l'hypothèse M1. La première est un ensemble
de petits multiples~: pour chaque année de développement $j$, le nuage des
couples $\big(C(i,j),\,C(i,j+1)\big)$, la droite $y = \hat f_j x$ --- celle
qu'impose le règlement, passant par l'origine --- et la droite ajustée avec
constante. La deuxième porte les ordonnées à l'origine estimées $\hat a_j$ avec
leur intervalle à 95~\%. La troisième compare les estimateurs
$\hat f_j^{(\alpha)}$ pour $\alpha = 0$, 1 et 2.
\Hzero
Sans objet. Les tests numériques correspondants sont les fiches~\ref{mw:m1a}
et~\ref{mw:m1d}.
\Stat
Sans objet (représentations graphiques).
\Loi
Sans objet.
\Pval
Sans objet.
\Usage
Les petits multiples sont la contrepartie visuelle directe du test de nullité de
l'ordonnée à l'origine~: \textbf{l'écart entre les deux droites \emph{est} le
test}. Ils permettent en outre trois lectures qu'aucun test ne fournit~:
identifier la colonne responsable d'un rejet, repérer le point qui l'entraîne, et
constater le rétrécissement du nuage à mesure que $n_j$ décroît --- rappel visuel
que les dernières colonnes ne reposent que sur deux ou trois points.

\medskip
Le graphique des ordonnées à l'origine résume les mêmes informations sous forme
compacte, les colonnes significatives étant mises en évidence. Celui de la
famille $\alpha$ chiffre l'enjeu du choix de pondération~: des courbes
confondues signifient que la question est sans portée pratique sur le triangle
considéré~; des courbes qui divergent sur une colonne désignent celle où
l'hypothèse de proportionnalité ou de variance est la plus fragile.

\medskip
Limite~: seules les neuf premières années de développement sont représentées en
petits multiples, au-delà la lisibilité s'effondre. Les colonnes suivantes
comptent de toute façon moins de trois observations et ne sont pas testées.
\end{fiche}

\begin{fiche}{Influence et leviers du triangle (Merz--Wüthrich)}
\refl{T.~Mack (1993), \emph{op.\ cit.}, pour la structure de pondération~;
R.\,D. Cook (1977) et D.\,C. Hoaglin, R.\,E. Welsch (1978), \emph{op.\ cit.},
pour les notions de levier et d'influence transposées ici.}
\Cadre
Trois représentations. Le facteur $\hat f_j$ étant une moyenne pondérée des
facteurs individuels de poids $C(i,j)$, le levier \emph{exact} de la cellule
$(i,j)$ dans son estimation vaut $h(i,j) = C(i,j)/S_j$, de somme~1 par colonne
(vérifié numériquement). L'influence exacte se mesure par le DFBETA
\[
  \hat f_j^{(-i)} = \frac{\sum_{i'} C(i',j+1) - C(i,j+1)}{\sum_{i'} C(i',j) - C(i,j)},
\]
obtenu sans réajustement itératif. \textbf{Diagnostics, non tests.}
\Hzero
Sans objet.
\Stat
Sans objet (représentations graphiques). Les quantités tracées sont $h(i,j)$, la
variation relative $\big(\hat f_j^{(-i)} - \hat f_j\big)/\hat f_j$, et la
décomposition de la réserve et de la variance de processus par année de
survenance (champ \code{\$terme\_variance} de \code{mw\_msep()}).
\Loi
Sans objet.
\Pval
Sans objet.
\Usage
Le premier graphique croise levier et résidu de Mack~: il sépare les cellules
atypiques par leur \emph{poids} de celles qui sont réellement influentes. Le
deuxième donne le DFBETA cellule par cellule, directement interprétable~: de
combien le facteur de développement bougerait si cette cellule était retirée. Le
troisième compare, année de survenance par année de survenance, la part de la
réserve et la part de la variance de processus de la MSEP.

\medskip
\textbf{Réserve de lecture importante.} Ce troisième graphique est une
\emph{décomposition} et non un leave-one-out. Le terme d'erreur d'estimation de
la MSEP (champ \code{\$terme\_covariance} de \code{mw\_msep()}) comporte des
termes croisés, partagés entre paires d'années de survenance, qui ne sont pas
attribuables à une année isolée~: seule la variance de processus, attribuable
année par année, y figure.
Il n'existe donc pas, pour la méthode Merz--Wüthrich, d'équivalent exact du
jackknife sur $\sigma_{\mathrm{USP}}$ disponible pour la méthode lognormale~;
c'est le diagnostic de concentration (fiche~\ref{mw:m6conc}) qui en tient lieu.
\end{fiche}

\clearpage
% =============================================================================
\section{Méthode Merz--Wüthrich --- cadre et spécification}
\label{sec:mw}
% =============================================================================

\subsection{Positionnement par rapport à la méthode lognormale}

L'annexe~XVII ouvre deux voies pour l'écart-type propre du risque de réserve. La
section~C (méthode n\textsuperscript{o}~1) repose sur un modèle lognormal appliqué
à deux vecteurs annuels~; la section~D (méthode n\textsuperscript{o}~2) repose sur
un modèle de chaîne d'échelle appliqué à un \emph{triangle} de paiements cumulés,
et sur l'erreur de prédiction à un an de Merz--Wüthrich. Les deux méthodes visent
le même paramètre mais reposent sur des hypothèses de nature différente, ce qui
commande des batteries de tests distinctes.

\begin{center}
\small
\rowcolors{2}{tablerow}{white}
\begin{tabular}{@{}p{3.1cm}p{5.7cm}p{6.0cm}@{}}
\rowcolor{tableheader}
& \textcolor{financeblue}{\textbf{Méthode lognormale (B et C)}} &
\textcolor{financeblue}{\textbf{Merz--Wüthrich (D)}} \\
Donnée d'entrée & deux vecteurs $x_t$, $y_t$ de longueur $T$ &
triangle $C(i,j)$, $(I+1)\times(J+1)$ \\
Modèle & $\mathbb{E}[Y_t] = \beta x_t$, variance quadratique, \textbf{loi lognormale} &
chaîne d'échelle~: $\mathbb{E}[C(i,j+1)\mid C(i,j)] = f_j\,C(i,j)$, variance
proportionnelle à $C(i,j)$ \\
Loi spécifiée & oui, entièrement & \textbf{non}~: seuls les deux premiers moments \\
Paramètres estimés & $\beta$, $\delta$, $\gamma$ (maximum de vraisemblance) &
$f_j$, $\sigma_j^2$ (moindres carrés pondérés) \\
Estimation de l'écart-type & fonction $\sigma(\delta,\gamma)$ &
$\sqrt{\mathrm{MSEP}}$ rapportée à la réserve \\
Correction de taille finie & $\sqrt{(T+1)/(T-1)}$ & aucune \\
Durée de crédibilité & années de survenance (B) ou exercices (C) &
années de survenance, section~G(3)(c) \\
Nature du bootstrap & paramétrique (loi entièrement spécifiée) &
semi-paramétrique, par rééchantillonnage des résidus \\
\end{tabular}
\end{center}

\subsection{Spécification réglementaire}

Les formules ci-dessous sont la transcription littérale de l'annexe~XVII,
section~D, paragraphes~3 à~5, sous leur \emph{forme imprimée} et non sous une
forme algébriquement équivalente plus compacte~: un relecteur doit pouvoir les
comparer au texte terme à terme. Les années de survenance sont indexées
$i = 0,\dots,I$, les années de développement $j = 0,\dots,J$, et $C(i,j)$ désigne
les sinistres cumulés.

\medskip
\noindent
\emph{Version du texte et pagination.} Les formules ont été relevées sur la
version consolidée en vigueur de l'annexe~XVII. Pour le paragraphe~5, le
libellé publié en 2015 diffère de celui en vigueur~: c'est ce dernier qui est
reproduit ci-dessous et qu'implémente \code{mw\_msep()}. L'acte qui a opéré
cette correction n'est pas identifié à ce jour~; l'hypothèse d'un rectificatif
refondu sans marqueur de modification, la seule que suggèrent les pièces
consultées, est signalée comme une hypothèse et non comme un fait établi. La
pagination « L~12/277-278 », qui circule dans les documents de travail du
projet, n'a pas été vérifiée sur le \emph{Journal officiel}~; elle n'est donc
pas reprise ici.

\paragraph{Facteurs de développement.}
\[
  \hat f_j = \frac{\sum_{i=0}^{I-j-1} C(i,j+1)}{\sum_{i=0}^{I-j-1} C(i,j)},
  \qquad
  \hat C(i,j) = C(i,I-i)\,\hat f_{I-i}\cdots\hat f_{j-2}\,\hat f_{j-1}.
\]

\paragraph{Variances de développement.} Pour $j = 0,\dots,J-2$,
\[
  \hat\sigma_j^2 = \frac{1}{I-j-1}\sum_{i=0}^{I-j-1} C(i,j)
  \left(\frac{C(i,j+1)}{C(i,j)} - \hat f_j\right)^2,
\]
et, pour la dernière année de développement, le règlement impose l'extrapolation
\[
  \hat\sigma_{J-1}^2 = \min\!\left(\hat\sigma_{J-2}^2,\ \hat\sigma_{J-3}^2,\
  \frac{\hat\sigma_{J-2}^4}{\hat\sigma_{J-3}^2}\right).
\]

\paragraph{Erreur quadratique moyenne de prédiction.} En posant
$\hat Q_j = \hat\sigma_j^2/\hat f_j^2$,
$S_j = \sum_{i=0}^{I-j-1} C(i,j)$ et $S'_j = \sum_{i=0}^{I-j} C(i,j)$~:
\[
\begin{aligned}
  \mathrm{MSEP} = {} & \sum_{i=1}^{I}\hat C(i,J)^2
  \left(\frac{\hat Q_{I-i}}{C(i,I-i)} + \frac{\hat Q_{I-i}}{S_{I-i}}
  + \sum_{j=I-i+1}^{J-1}\frac{C(I-j,j)}{S'_j}\cdot\frac{\hat Q_j}{S_j}\right) \\[2pt]
  & + 2\sum_{i=1}^{I}\sum_{k=i+1}^{I}\hat C(i,J)\,\hat C(k,J)
  \left(\frac{\hat Q_{I-i}}{S_{I-i}}
  + \sum_{j=I-i+1}^{J-1}\frac{C(I-j,j)}{S'_j}\cdot\frac{\hat Q_j}{S_j}\right).
\end{aligned}
\]

\noindent
Le crochet, noté $\Delta_i$ comme dans \code{mw\_msep()}, figure
\textbf{deux fois}~: une première fois à l'intérieur de la première somme, à
côté du terme de variance de processus $\hat Q_{I-i}/C(i,I-i)$, et une seconde
fois dans la double somme, affectée du \textbf{facteur~2}. La première somme
contient donc les termes diagonaux $\hat C(i,J)^2\,\Delta_i$. Les bornes
$k = i+1,\dots,I$, l'indice $I-i$ porté par le crochet, le facteur~2 et
l'exposant~1 sur $C(I-j,j)/S'_j$ sont ceux du texte. Cette notation
$\Delta_i$, propre à la méthode Merz--Wüthrich, est sans rapport avec la marge
d'équivalence $\Delta$ du test TOST de la méthode lognormale.

\begin{encadre}
\textbf{Forme imprimée et forme de Merz--Wüthrich.} Pour $k > i$, le crochet
indexé par $i$ est celui indexé par $\min(i,k)$, et $2\sum_{i<k}$ reproduit
exactement la double somme complète hors diagonale~; les termes diagonaux étant
portés par la première somme, la formule imprimée ci-dessus est identiquement
l'erreur de prédiction à un an de Merz--Wüthrich (2008). Le moteur transcrit la
\emph{forme imprimée}~; cette identité est \emph{vérifiée} numériquement
(section~\ref{sec:verif}) et non présumée.

\medskip
\textbf{Restitution du résultat.} \code{mw\_msep()} renvoie la MSEP et deux
composantes~: \code{\$terme\_variance} $=\sum_{i=1}^{I}\hat C(i,J)^2\,
\hat Q_{I-i}/C(i,I-i)$, la variance de processus seule, et
\code{\$terme\_covariance} $=\sum_{i=1}^{I}\hat C(i,J)^2\,\Delta_i
+ 2\sum_{i<k}\hat C(i,J)\,\hat C(k,J)\,\Delta_i$, l'erreur d'estimation, ses
termes diagonaux compris. Leur somme est exactement la MSEP du texte. Ce
découpage est un choix de présentation~: le règlement ne prescrit aucune
décomposition.
\end{encadre}

\paragraph{Paramètre propre.}
\[
  \sigma_{(\text{res},s,\mathrm{USP})}
  = c\cdot\frac{\sqrt{\mathrm{MSEP}}}{\sum_{i=0}^{I}\big(\hat C(i,J) - C(i,I-i)\big)}
  + (1-c)\,\sigma_{(\text{res},s)} .
\]

\begin{encadre}
\textbf{Restriction d'implémentation.} Le paragraphe~2(e) autorise $I \ge J$, mais
la formule du paragraphe~4 fait intervenir $C(i, I-i)$, qui n'est définie que si
$I - i \le J$. Le cas $I > J$ produirait un trapèze dont le texte ne précise pas
le traitement. Le moteur impose donc $I = J$ (triangle carré) et refuse
explicitement les autres géométries, plutôt que de les traiter par une convention
implicite non prévue par le règlement.
\end{encadre}

\subsection{Hypothèses testables et résidus}

Le paragraphe~2(h) énonce quatre hypothèses, qui structurent la batterie de tests
au même titre que H1 à H4 pour la méthode lognormale~:

\begin{center}
\small
\rowcolors{2}{tablerow}{white}
\begin{tabular}{@{}p{1.5cm}p{6.6cm}p{6.7cm}@{}}
\rowcolor{tableheader}
\textcolor{financeblue}{\textbf{Code}} &
\textcolor{financeblue}{\textbf{Énoncé (annexe XVII, D(2)(h))}} &
\textcolor{financeblue}{\textbf{Correspondance lognormale}} \\
M1 & (iii) $\mathbb{E}[C(i,j+1)\mid C(i,j)] = f_j\,C(i,j)$ &
H1 --- même forme de proportionnalité sans constante \\
M2 & (iv) $\operatorname{Var}[C(i,j+1)\mid C(i,j)] = \sigma_j^2\,C(i,j)$ &
H2 --- mais puissance 1 imposée, sans mélange $\delta$ à estimer \\
M3 & (i) et (ii) indépendance entre années de survenance et entre montants
incrémentaux & H4 --- indépendance \\
M5 & \emph{aucune} & H3 --- \textbf{sans équivalent}~: la lognormalité n'est
pas requise \\
\end{tabular}
\end{center}

Le support des tests est le résidu standardisé de Mack
\[
  r(i,j) = \frac{\sqrt{C(i,j)}}{\hat\sigma_j}
  \left(\frac{C(i,j+1)}{C(i,j)} - \hat f_j\right),
\]
centré, d'échelle approximativement unitaire et non corrélé sous M1 à M3.

\begin{encadre}
\textbf{La normalité n'est pas une hypothèse de la méthode.} Le paragraphe~2(h)
ne spécifie que les deux premiers moments conditionnels. Les tests de normalité
figurent donc dans la famille M5 comme \emph{diagnostics}, et non comme
vérifications d'une hypothèse réglementaire. Ils ne deviennent pertinents que si
l'on souhaite exploiter la MSEP pour construire un quantile, usage qui excède la
méthode elle-même.
\end{encadre}

\subsection{Réemploi et adaptation des tests}

\begin{center}
\small
\rowcolors{2}{tablerow}{white}
\begin{tabular}{@{}p{4.9cm}p{2.3cm}p{7.6cm}@{}}
\rowcolor{tableheader}
\textcolor{financeblue}{\textbf{Test}} &
\textcolor{financeblue}{\textbf{Statut}} &
\textcolor{financeblue}{\textbf{Commentaire}} \\
Durbin--Watson, test des suites, Grubbs, Rosner, Shapiro--Wilk, Lilliefors &
réemployés & Appliqués aux résidus de Mack au lieu des résidus normalisés~; la
statistique est inchangée, seule la loi de référence est réévaluée. \\
Breusch--Pagan studentisé & adapté & Régression des carrés des résidus sur
$C(i,j)$ et non sur le volume $x_t$~: teste la puissance~1 de M2. \\
Test de proportionnalité & adapté & Régression pondérée des facteurs $F(i,j)$
sur $C(i,j)$~: une pente non nulle contredit M1. \\
Test des effets d'année calendaire & \textbf{nouveau} & Propre aux triangles~:
les diagonales représentent les exercices comptables. Sans équivalent dans la
méthode lognormale. \\
Corrélation entre années de développement & \textbf{nouveau} & Teste M3(ii)
entre colonnes adjacentes. \\
Mann--Kendall, Cox--Stuart, Spearman, RESET, TOST, sup-F, CUSUM &
sans objet & Ces tests portent sur une série chronologique de ratios, objet qui
n'existe pas dans la géométrie du triangle. \\
\end{tabular}
\end{center}

\subsubsection*{Synthèse de couverture}

\begin{center}
\small
\rowcolors{2}{tablerow}{white}
\begin{tabular}{@{}p{1.2cm}p{6.0cm}p{3.2cm}p{3.9cm}@{}}
\rowcolor{tableheader}
\textcolor{financeblue}{\textbf{Famille}} &
\textcolor{financeblue}{\textbf{Entrées}} &
\textcolor{financeblue}{\textbf{Type}} &
\textcolor{financeblue}{\textbf{p-value retenue}} \\
M1 & Tendance avec le cumulé, à colonne donnée & test & Monte-Carlo \\
M1 & \textbf{Ordonnée à l'origine, colonne par colonne} & test & Monte-Carlo \\
M1 & \textbf{Homogénéité de $f_j$ entre années de survenance} & test & Monte-Carlo \\
M1 & \textbf{Absence de courbure} & test & Monte-Carlo \\
M1 & \textbf{Famille $\alpha$ (pondération)} & test & Monte-Carlo \\
M2 & Hétéroscédasticité résiduelle contre le cumulé & test & Monte-Carlo \\
M2 & \textbf{Exposant de variance, colonne par colonne} & test & Monte-Carlo \\
M2 & Variance unitaire des résidus & diagnostic & --- \\
M3 & \textbf{Homogénéité des résidus entre années de survenance} & test & Monte-Carlo \\
M3 & Effets d'année calendaire & test & Monte-Carlo \\
M3 & Corrélation entre années de développement & test & Monte-Carlo \\
M3 & Autocorrélation (Durbin--Watson) & test & Monte-Carlo \\
M3 & Test des suites & test & Monte-Carlo \\
M4 & Cellule aberrante (Grubbs) & test & Monte-Carlo \\
M4 & Cellules aberrantes multiples (Rosner) & procédure de décision & --- \\
M5 & Shapiro--Wilk & test (diagnostic) & exacte simulée \\
M5 & Lilliefors & test (diagnostic), secondaire & asymptotique \\
M6 & Extrapolation de $\hat\sigma^2_{J-1}$ & diagnostic & --- \\
M6 & Concentration de la réserve & diagnostic & --- \\
\end{tabular}
\end{center}

\noindent
Dix-neuf entrées, dont quatorze tests d'hypothèse, une procédure de décision et
quatre diagnostics. Treize des quatorze tests reposent sur le bootstrap de
résidus~; seul Shapiro--Wilk dispose d'une loi nulle propre, pour la raison
exposée dans sa fiche. Les six entrées en gras ci-dessus sont celles qui
vérifient les hypothèses de chaîne d'échelle \emph{colonne par colonne}, comme
l'exige la formulation « pour toutes les années d'accident » du texte~; la
régression agrégée initiale, conservée en tête, ne le permettait pas.

\clearpage
% =============================================================================
\section{Hypothèse M1 --- proportionnalité des cumulés}
\label{sec:mwm1}
% =============================================================================

\begin{encadre}
\textbf{Énoncé.} annexe~XVII, D(2)(h)(iii)~: $\mathbb{E}[C(i,j+1)\mid C(i,j)] = f_j\,C(i,j)$.
\end{encadre}

Les fiches ci-dessous suivent la structure en six rubriques employée pour les
hypothèses H1 à H4 de la méthode lognormale. Les notations sont celles de la
section~\ref{sec:notations-mw}.

\begin{fiche}{M1 --- Absence de tendance des facteurs avec le cumulé, à colonne donnée}
\label{mw:m1}
\refl{T.~Mack, \emph{Distribution-free calculation of the standard error of chain
ladder reserve estimates}, \emph{ASTIN Bulletin}, 23(2), 1993, section~3.}
\Cadre
Régression auxiliaire menée sur l'ensemble des cellules du triangle~:
\[
  F(i,j) = \mu_j + b\,C(i,j) + \varepsilon_{ij},
  \qquad \text{poids } w_{ij} = C(i,j)/\hat\sigma_j^2 ,
\]
où $\mu_j$ est un \textbf{effet fixe d'année de développement}. Deux points de
construction sont déterminants et ont été établis par simulation.

\emph{L'effet fixe de colonne est indispensable.} Le facteur $f_j$ décroît
fortement avec $j$ (de $3{,}49$ à $1{,}02$ sur le triangle de contrôle) alors que
le cumulé $C(i,j)$ croît~: une régression agrégée \emph{sans} terme de colonne
capte cette relation mécanique et non une dépendance au volume. Sur des triangles
simulés \emph{exactement} sous $H_0$, une telle régression rejette dans
\textbf{100~\% des cas}. Avec l'effet fixe, la fréquence de rejet revient à
$0{,}068$ au seuil nominal de $0{,}10$.

\emph{La pondération doit être celle des moindres carrés généralisés.} Puisque
$\operatorname{Var}\big(F(i,j)\big) = \sigma_j^2/C(i,j)$, le poids exact est
$C(i,j)/\sigma_j^2$ et non $C(i,j)$ seul~: $\sigma_j^2$ varie de plusieurs ordres
de grandeur entre la première et la dernière colonne, et pondérer par $C(i,j)$
seul rend l'échelle résiduelle dominée par les premières colonnes --- au point
que le test ne rejette \emph{jamais} sous l'alternative comme sous l'hypothèse
nulle.
\Hzero
$H_0 : b = 0$ --- à année de développement donnée, le facteur de développement ne
dépend pas du niveau du cumulé.
\Stat
$t = \hat b / \widehat{\mathrm{se}}(\hat b)$ dans la régression pondérée à effet
fixe de colonne.
\Loi
Loi de Student approchée~: les cellules d'une même colonne partagent $\hat f_j$
et $\hat\sigma_j$, ce qui invalide les degrés de liberté nominaux. Loi de
référence~: bootstrap de résidus.
\Pval
p-value de Monte-Carlo retenue~; la p-value de Student est fournie comme
complément non simulé.
\Usage
Ce test est le pendant \emph{agrégé} du test d'ordonnée à l'origine par colonne
qui suit~: il gagne en puissance ce que l'autre gagne en localisation. Puissance
mesurée sous une alternative où le facteur dépend linéairement du volume~:
$100~\%$ de rejets au seuil de $10~\%$. Il ne dit pas, en revanche, quelle
colonne est en cause~: c'est le rôle du test par colonne et du graphique en
petits multiples.
\end{fiche}

\begin{fiche}{M1 --- Nullité de l'ordonnée à l'origine, colonne par colonne}
\label{mw:m1a}
\refl{T.~Mack, \emph{Distribution-free calculation of the standard error of chain
ladder reserve estimates}, \emph{ASTIN Bulletin}, 23(2), 1993, section~3 ;
combinaison des p-values : R.\,A. Fisher, \emph{Statistical Methods for Research
Workers}, 4\textsuperscript{e} éd., Oliver \& Boyd, 1932, §21.1.}
\Cadre
Pour chaque année de développement $j$ comptant au moins trois observations, on
ajuste la régression pondérée
\[
  C(i,j+1) = a_j + b_j\,C(i,j) + \varepsilon_{ij},
  \qquad \text{poids } w_{ij} = 1/C(i,j),
\]
sur les $n_j = I-j$ couples disponibles. Le choix du poids n'est pas arbitraire~:
il découle de l'hypothèse de variance M2, et il possède une propriété
remarquable établie par Mack --- \textbf{la régression pondérée \emph{sans}
constante et de poids $1/C(i,j)$ a pour estimateur des moindres carrés
exactement le facteur chain-ladder}~:
\[
  \hat b_j^{\,(\text{sans constante})}
  = \frac{\sum_i w_{ij} C(i,j)\,C(i,j+1)}{\sum_i w_{ij} C(i,j)^2}
  = \frac{\sum_i C(i,j+1)}{\sum_i C(i,j)} = \hat f_j .
\]
Cette identité a été vérifiée numériquement colonne par colonne, à $10^{-16}$
près. Elle est ce qui fait de l'ajout d'une constante le test \emph{naturel} de
l'hypothèse réglementaire. Hypothèses de validité non testées~: indépendance des
lignes, variance correctement pondérée, erreurs approximativement normales.
\Hzero
$H_0 : a_j = 0$ pour toute année de développement $j$, c'est-à-dire la
proportionnalité stricte exigée par D(2)(h)(iii).
\Stat
Statistique de Student $t_j = \hat a_j / \widehat{\mathrm{se}}(\hat a_j)$ dans
chaque colonne, puis combinaison des $K$ p-values par la méthode de Fisher~:
$X = -2\sum_{k} \ln p_k$.
\Loi
$X \sim \chi^2(2K)$ \emph{si les colonnes étaient indépendantes}. Elles ne le
sont pas tout à fait~: les colonnes $j$ et $j+1$ partagent le vecteur
$C(\cdot,j+1)$. La p-value du khi-deux est donc rapportée comme non simulée et
\emph{indicative}~; la loi de référence retenue est celle du bootstrap de
résidus.
\Pval
p-value de Monte-Carlo retenue. Calibration vérifiée par simulation sous $H_0$
(1\,500 triangles $10\times10$)~: 10,7~\% de rejets au seuil nominal de 10~\% et
4,7~\% au seuil de 5~\%.
\Usage
C'est le test central de M1, et celui que la régression agrégée ne pouvait pas
fournir. Il répond à la première des deux exigences du texte réglementaire~: une
proportionnalité \emph{sans constante}, vérifiée \emph{pour chaque} année de
développement. Le graphique des petits multiples (section~\ref{sec:graph-mw})
en donne la contrepartie visuelle~: la droite pleine est la proportionnalité
imposée par le règlement, la droite pointillée l'ajustement libre~; l'écart
entre les deux \emph{est} le test.

\medskip
Limite~: le nombre d'observations décroît de $I$ à~3 le long du triangle, si
bien que les dernières colonnes ne contribuent presque rien. Le tableau détaillé
restitué par le moteur donne la colonne responsable, ce qui est la seule lecture
actionnable.
\end{fiche}

\begin{fiche}{M1 --- Homogénéité de $f_j$ entre années de survenance}
\label{mw:m1b}
\refl{Annexe XVII, D(2)(h)(iii), dont la formulation « pour toutes les années
d'accident » constitue la seconde exigence ; C.~Spearman (1904),
\emph{op.\ cit.} pour la corrélation de rang.}
\Cadre
Dans chaque colonne $j$ comptant au moins quatre observations, on mesure la
corrélation de rang entre les facteurs individuels $F(i,j)$ et l'indice $i$ de
l'année de survenance. Hypothèses non testées~: continuité, absence d'ex
\ae quo.
\Hzero
$H_0$ : le facteur $f_j$ est \emph{commun à toutes les années de survenance}~;
les facteurs individuels ne présentent aucune tendance en $i$.
\Stat
$\hat\rho_{s,j}$ par colonne, p-values combinées par la méthode de Fisher.
\Loi
$\chi^2(2K)$ approchée, avec la même réserve d'indépendance que la fiche
précédente. Loi de référence~: bootstrap de résidus.
\Pval
p-value de Monte-Carlo retenue.
\Usage
Ce test répond à la \emph{seconde} exigence du texte, distincte de la première~:
même si la proportionnalité vaut dans chaque colonne, un facteur qui dériverait
d'une année de survenance à l'autre invaliderait le modèle. En pratique, une
telle dérive traduit un changement de politique de règlement ou de mix de
portefeuille. Elle n'est pas détectable par le test d'ordonnée à l'origine, qui
porte sur le niveau et non sur la pente.
\end{fiche}

\begin{fiche}{M1 --- Absence de courbure de la régression}
\label{mw:m1c}
\refl{Test du terme quadratique, dans l'esprit de J.\,B. Ramsey (1969),
\emph{op.\ cit.} ; appliqué ici colonne par colonne.}
\Cadre
Régression pondérée augmentée d'un terme quadratique,
$C(i,j+1) = a_j + b_j C(i,j) + c_j C(i,j)^2 + \varepsilon_{ij}$, poids
$1/C(i,j)$, sur les colonnes comptant au moins quatre observations.
\Hzero
$H_0 : c_j = 0$ dans chaque colonne.
\Stat
Statistique de Student sur $\hat c_j$, p-values combinées par Fisher.
\Loi
$\chi^2(2K)$ approchée. Loi de référence~: bootstrap de résidus.
\Pval
p-value de Monte-Carlo retenue.
\Usage
La linéarité et la nullité de la constante sont deux propriétés distinctes~: une
relation courbe peut passer par l'origine et satisfaire le test précédent tout en
contredisant $\mathbb{E}[C(i,j+1)\mid C(i,j)] = f_j C(i,j)$. Ce test ferme cette
possibilité. Limite sévère à faible profondeur~: estimer trois paramètres sur
quatre à six observations laisse très peu de degrés de liberté, et la puissance
est faible.
\end{fiche}

\begin{fiche}{M1 --- Stabilité du facteur selon la pondération (famille $\alpha$)}
\label{mw:m1d}
\refl{T.~Mack, \emph{Which stochastic model is underlying the chain ladder
method~?}, \emph{Insurance: Mathematics and Economics}, 15(2--3), 1994.}
\Cadre
On considère la famille d'estimateurs
\[
  \hat f_j^{(\alpha)}
  = \frac{\sum_i C(i,j)^{\alpha} F(i,j)}{\sum_i C(i,j)^{\alpha}},
  \qquad \alpha \in \{0,\,1,\,2\},
\]
où $\alpha = 0$ donne la moyenne simple des facteurs individuels, $\alpha = 1$ le
facteur chain-ladder du règlement, et $\alpha = 2$ une pondération par le carré
du volume. \textbf{Sous l'hypothèse M1, les trois estiment la même quantité}~:
l'espérance conditionnelle étant proportionnelle à $C(i,j)$, tout jeu de poids
positifs fournit un estimateur sans biais de $f_j$.
\Hzero
$H_0$ : les trois estimateurs visent le même $f_j$ dans chaque colonne.
\Stat
Amplitude relative moyenne, pondérée par les effectifs~:
\[
  A = \frac{\sum_j n_j\,
      \big(\max_\alpha \hat f_j^{(\alpha)} - \min_\alpha \hat f_j^{(\alpha)}\big)
      \big/ \hat f_j^{(1)}}{\sum_j n_j}.
\]
\Loi
\textbf{Aucune loi analytique}~: la statistique agrège trois estimateurs
fortement corrélés sur des colonnes de tailles décroissantes.
\Pval
p-value de Monte-Carlo exclusivement.
\Usage
Diagnostic de robustesse plus que test d'hypothèse. Une divergence marquée
signale soit que l'espérance n'est pas proportionnelle à $C(i,j)$ (violation de
M1), soit que la pondération en $C(i,j)$ est inadaptée (violation de M2)~: le
test ne sépare pas les deux causes, et doit donc être lu conjointement avec la
fiche~\ref{mw:m2b}. Son intérêt propre est de chiffrer l'enjeu~: si $A$ est
faible, le choix de pondération imposé par le règlement est sans conséquence
matérielle sur le paramètre final.
\end{fiche}

\begin{fiche}{M2 --- Hétéroscédasticité résiduelle contre le cumulé}
\label{mw:m2bp}
\refl{T.\,S. Breusch, A.\,R. Pagan (1979), \emph{op.\ cit.}~; version
studentisée de R.~Koenker (1981), \emph{op.\ cit.}, appliquée ici aux résidus de
Mack.}
\Cadre
Régression auxiliaire $r(i,j)^2 = c_0 + c_1\,C(i,j) + \eta$, de coefficient de
détermination $R^2_{\text{aux}}$, sur les $\sum_j n_j$ cellules du triangle.
Hypothèses non testées~: indépendance des cellules, validité de l'approximation
asymptotique.
\Hzero
$H_0 : c_1 = 0$ --- les résidus standardisés ne dépendent plus de $C(i,j)$, donc
la pondération en $C(i,j)$ imposée par M2 capture correctement la variance.
\Stat
$\mathrm{LM} = n\,R^2_{\text{aux}}$, où $n$ est le nombre de résidus.
\Loi
$\chi^2(1)$ asymptotique. L'approximation est doublement fragile ici~: le nombre
de cellules est modeste et les résidus d'une même colonne sont liés par
$\hat f_j$ et $\hat\sigma_j$.
\Pval
p-value de Monte-Carlo retenue~; la p-value du khi-deux est fournie comme
complément non simulé. \textbf{Celle-ci sous-rejette massivement}~: $0{,}5~\%$ de
rejets mesurés au seuil nominal de $10~\%$ (section~\ref{sec:calib-asympt-mw}),
car $\hat\sigma_j^2$ est estimé sur la colonne même, ce qui réduit la dispersion
des résidus standardisés et donc le $R^2$ de la régression auxiliaire.
\Usage
Ce test porte spécifiquement sur \emph{l'exposant~1} de la structure de variance
$\operatorname{Var}[C(i,j+1)\mid C(i,j)] = \sigma_j^2\,C(i,j)$. Un rejet suggère
un exposant différent --- variance constante ($\alpha = 0$) ou proportionnelle au
carré ($\alpha = 2$) --- alternatives discutées par Barnett \& Zehnwirth mais
\emph{non prévues par l'annexe~XVII}, qui impose $\alpha = 1$. Le rejet doit donc
être documenté comme une limite du modèle réglementaire, non comme un motif de
changement de méthode.
\end{fiche}

% -----------------------------------------------------------------------------
\begin{fiche}{M2 --- Adéquation de l'exposant de variance, colonne par colonne}
\label{mw:m2b}
\refl{Annexe XVII, D(2)(h)(iv) ; complément par colonne du test de
Breusch--Pagan de la fiche précédente.}
\Cadre
Dans chaque colonne $j$ comptant au moins quatre résidus, on mesure la
corrélation de rang entre $|r(i,j)|$ et $C(i,j)$.
\Hzero
$H_0$ : l'amplitude des résidus standardisés ne dépend pas de $C(i,j)$ à
l'intérieur d'une colonne, ce qui est la traduction directe de
$\operatorname{Var}[C(i,j+1)\mid C(i,j)] = \sigma_j^2\,C(i,j)$.
\Stat
$\hat\rho_{s,j}$ par colonne, p-values combinées par Fisher.
\Loi
$\chi^2(2K)$ approchée. Loi de référence~: bootstrap de résidus.
\Pval
p-value de Monte-Carlo retenue.
\Usage
Le test de Breusch--Pagan de la fiche précédente agrège toutes les colonnes~; ce
test-ci localise. Un rejet signale que l'exposant~1 imposé par le règlement est
inadapté --- l'alternative étant un exposant~0 (variance constante) ou~2
(proportionnelle au carré). \textbf{Le règlement ne laissant pas ce choix
ouvert}, un rejet doit être documenté comme une limite du modèle réglementaire
et non comme un motif de changement de méthode.
\end{fiche}

\begin{fiche}{M2 --- Variance unitaire des résidus de Mack (diagnostic)}
\label{mw:m2var}
\refl{Diagnostic d'échelle~; Mack (1993), \emph{op.\ cit.}, définition de
$\hat\sigma_j^2$.}
\Cadre
Résidus standardisés $r(i,j)$. \textbf{Diagnostic, non test}~: la rubrique est
renseignée par cohérence de présentation.
\Hzero
Sans objet. Par construction, $\hat\sigma_j^2$ est calibré colonne par colonne,
de sorte que la variance des résidus vaut approximativement~1 \emph{à
l'intérieur} de chaque colonne, sans que la variance globale soit contrainte.
\Stat
$\operatorname{Var}\big(r(i,j)\big)$ sur l'ensemble du triangle.
\Loi
Sans objet.
\Pval
Aucune. Le moteur signale un écart supérieur à $0{,}5$ en valeur absolue.
\Usage
Une variance globale nettement différente de~1 signale une mauvaise
spécification de $\hat\sigma_j$, le plus souvent concentrée sur les dernières
colonnes, où $n_j$ tombe à deux ou trois observations. À lire conjointement avec
la fiche sur l'extrapolation de $\hat\sigma^2_{J-1}$.
\end{fiche}

% -----------------------------------------------------------------------------
\clearpage
% =============================================================================
\section{Hypothèse M3 --- indépendance}
\label{sec:mwm3}
% =============================================================================

\begin{encadre}
\textbf{Énoncé.} annexe~XVII, D(2)(h)(i) et (ii)~: indépendance des années de survenance et des montants incrémentaux.
\end{encadre}

Les fiches ci-dessous suivent la structure en six rubriques employée pour les
hypothèses H1 à H4 de la méthode lognormale. Les notations sont celles de la
section~\ref{sec:notations-mw}.

\begin{fiche}{M3 --- Effets d'année calendaire (test de Mack)}
\label{mw:m3cal}
\refl{T.~Mack, \emph{Which stochastic model is underlying the chain ladder
method~?}, \emph{Insurance: Mathematics and Economics}, 15(2--3), 1994,
p.~133--138~; T.~Mack (1993), \emph{op.\ cit.}, section~5.}
\Cadre
Dans chaque colonne $j$, les facteurs observés $F(i,j)$ sont classés en
« grands » ($L$) et « petits » ($S$) par rapport à leur médiane, les valeurs
égales à la médiane étant écartées. Les diagonales $k = i+j$ correspondent aux
exercices comptables. Hypothèses non testées~: continuité des facteurs,
indépendance des colonnes.
\Hzero
$H_0$ : absence d'effet d'année calendaire --- la répartition des $L$ et des $S$
le long de chaque diagonale est purement aléatoire, tous les arrangements étant
équiprobables.
\Stat
Sur la diagonale $k$ comportant $n_k$ étiquettes dont $L_k$ « grandes »~:
$Z_k = \min(L_k,\,n_k - L_k)$, puis $Z = \sum_k Z_k$ et
\[
  \frac{Z - \mathbb{E}[Z]}{\sqrt{\operatorname{Var}[Z]}} .
\]
\Loi
Sous $H_0$, $L_k \sim \mathcal{B}(n_k, 1/2)$, d'où les moments \textbf{exacts},
avec $m_k = \lfloor (n_k-1)/2 \rfloor$~:
\[
  \mathbb{E}[Z_k] = \frac{n_k}{2} - \binom{n_k-1}{m_k}\frac{n_k}{2^{n_k}},
  \quad
  \operatorname{Var}[Z_k] = \frac{n_k(n_k-1)}{4}
  - \binom{n_k-1}{m_k}\frac{n_k(n_k-1)}{2^{n_k}} + \mathbb{E}[Z_k] - \mathbb{E}[Z_k]^2 .
\]
Ces deux expressions ont été vérifiées par énumération exhaustive de la loi de
$\min(L, n-L)$ pour $n = 2,\dots,12$~: concordance à $10^{-10}$ près. En
revanche, la normalité de la \emph{somme} $Z$ n'est qu'approchée.
\Pval
$p = 2\big(1 - \Phi(|Z^\star|)\big)$ pour la version asymptotique~; la p-value de
Monte-Carlo est celle retenue.
\Usage
C'est le test le plus informatif du dispositif Merz--Wüthrich, car il détecte ce
que ni M1 ni M2 ne voient~: un choc affectant un exercice comptable entier
(inflation, changement de cadence de règlement, modification de la politique de
provisionnement). Un tel choc viole l'indépendance des années de survenance
posée par D(2)(h)(i). Limite à $T = 8$~: la somme ne porte que sur $I-1$
diagonales exploitables, dont plusieurs ne comptent que deux ou trois
étiquettes, d'où le recours au bootstrap plutôt qu'à l'approximation normale.
\end{fiche}

% -----------------------------------------------------------------------------
\begin{fiche}{M3 --- Homogénéité des résidus entre années de survenance}
\label{mw:m3acc}
\refl{W.\,H. Kruskal, W.\,A. Wallis, \emph{Use of ranks in one-criterion variance
analysis}, \emph{Journal of the American Statistical Association}, 47(260),
1952, p.~583--621.}
\Cadre
Les résidus de Mack sont groupés par année de survenance $i$. Le test compare
les rangs entre groupes, sans hypothèse de loi. Hypothèses non testées~:
indépendance des résidus, continuité.
\Hzero
$H_0$ : les résidus ont la même distribution dans toutes les années de
survenance, ce qui est la traduction testable de l'indépendance stochastique des
années de survenance posée par D(2)(h)(i).
\Stat
$H$ de Kruskal--Wallis, sur $k-1$ degrés de liberté où $k$ est le nombre
d'années de survenance comportant au moins un résidu.
\Loi
$\chi^2(k-1)$ asymptotique, approximation médiocre lorsque les groupes comptent
une ou deux observations --- ce qui est le cas des dernières lignes du triangle.
Loi de référence~: bootstrap de résidus.
\Pval
p-value de Monte-Carlo retenue~; le khi-deux est fourni comme complément non
simulé.
\Usage
Complète le test des effets d'année calendaire, qui porte sur les diagonales~:
celui-ci porte sur les \emph{lignes}. Un rejet signale qu'une année de
survenance se comporte différemment des autres, ce qui viole l'hypothèse
d'indépendance et d'identité de distribution. Limite structurelle~: les groupes
sont très déséquilibrés, la première ligne comptant $J$ résidus et la dernière
un seul.
\end{fiche}

\begin{fiche}{M3 --- Corrélation entre années de développement adjacentes}
\label{mw:m3cor}
\refl{T.~Mack (1993), \emph{op.\ cit.}, section~4~; C.~Spearman (1904),
\emph{op.\ cit.} pour la corrélation de rang.}
\Cadre
Pour chaque paire de colonnes adjacentes $(k-1 \to k)$ et $(k \to k+1)$, on
calcule la corrélation de rang de Spearman entre les facteurs individuels
observés sur les mêmes années de survenance. L'agrégation est pondérée par
$n_k - 1$. Hypothèses non testées~: continuité, absence d'ex \ae quo.
\Hzero
$H_0$ : les facteurs de développement successifs ne sont pas corrélés, ce qui
est la traduction testable de D(2)(h)(ii), l'indépendance des montants
incrémentaux à l'intérieur d'une année de survenance.
\Stat
$\displaystyle \hat\rho = \frac{\sum_k (n_k-1)\,\hat\rho_{s,k}}{\sum_k (n_k-1)}$,
moyenne pondérée des corrélations de rang par paire de colonnes.
\Loi
\textbf{Aucune loi analytique}~: la géométrie du triangle fait que les paires de
colonnes se recouvrent partiellement et reposent sur des effectifs décroissants.
Aucune approximation asymptotique n'est mobilisable.
\Pval
p-value de Monte-Carlo exclusivement. C'est, avec la corrélation, le seul test
du dispositif pour lequel aucune p-value non simulée n'existe.
\Usage
Un rejet signale que la cadence de liquidation d'une année de survenance est
persistante~: une année qui se développe vite sur une colonne continue de le
faire sur la suivante. Cela contredit l'indépendance des incréments et biaise la
MSEP, qui suppose que les erreurs de projection ne se cumulent pas. Limite~: sur
un triangle $8\times8$, seules cinq ou six paires de colonnes comptent au moins
trois observations, ce qui rend le test très peu puissant.
\end{fiche}

% -----------------------------------------------------------------------------
\begin{fiche}{M3 --- Autocorrélation des résidus de Mack (Durbin--Watson)}
\label{mw:m3dw}
\refl{J.~Durbin, G.\,S. Watson (1950, 1951), \emph{op.\ cit.}}
\Cadre
Résidus de Mack ordonnés selon leur parcours du triangle (colonne par colonne,
années de survenance croissantes). Hypothèses non testées~: résidus
homoscédastiques et de moyenne nulle sous le modèle.
\Hzero
$H_0$ : absence d'autocorrélation d'ordre~1 des résidus de Mack.
\Stat
$\mathrm{DW} = \sum_{t=2}^{n}(z_t - z_{t-1})^2 / \sum_{t=1}^{n} z_t^2$, calculée
sur les résidus centrés.
\Loi
\textbf{Contrairement à la méthode lognormale, la loi exacte d'Imhof n'est pas
applicable ici}~: les résidus de Mack ne forment pas un vecteur gaussien de
matrice de covariance connue, la structure de dépendance dépendant de la
géométrie du triangle. La loi de référence est celle du bootstrap de résidus.
\Pval
p-value de Monte-Carlo exclusivement.
\Usage
Complément aux deux tests précédents~: l'ordre de parcours étant conventionnel,
un rejet est plus difficile à interpréter que pour une série chronologique. Il
faut alors se reporter aux graphiques des résidus par année de survenance, de
développement et calendaire, qui localisent la structure.
\end{fiche}

% -----------------------------------------------------------------------------
\begin{fiche}{M3 --- Test des suites sur les résidus de Mack}
\label{mw:m3runs}
\refl{A.~Wald, J.~Wolfowitz (1940), \emph{op.\ cit.}}
\Cadre
Suite des signes des résidus de Mack par rapport à leur médiane, dans le même
ordre de parcours que ci-dessus. On note $n_1$ le nombre de signes positifs,
$n_2$ le nombre de négatifs et $R$ le nombre de suites.
\Hzero
$H_0$ : la suite des signes est un arrangement aléatoire.
\Stat
$Z = \big(R - \mathbb{E}[R]\big)/\sqrt{\operatorname{Var}[R]}$, avec
$\mathbb{E}[R] = 2n_1n_2/n + 1$ et
$\operatorname{Var}[R] = 2n_1n_2(2n_1n_2-n)/\big(n^2(n-1)\big)$.
\Loi
La loi combinatoire exacte du nombre de suites suppose l'échangeabilité des
signes, mise en défaut ici par la dépendance intra-colonne~: elle n'est donc pas
retenue, contrairement à ce qui est fait pour la méthode lognormale. Loi de
référence~: bootstrap de résidus.
\Pval
p-value de Monte-Carlo retenue~; l'approximation normale est fournie comme
complément non simulé. \textbf{Celle-ci sur-rejette nettement}~: $15{,}3~\%$ de
rejets mesurés au seuil nominal de $10~\%$ (section~\ref{sec:calib-asympt-mw}).
La cause est identifiée~: la contrainte
$\sum_i \sqrt{C(i,j)}\,r(i,j) = 0$ rend les résidus d'une même colonne
négativement corrélés, donc leurs signes plus alternés que sous l'échangeabilité
supposée.
\Usage
Test non paramétrique robuste, sensible aussi bien au regroupement qu'à
l'alternance excessive des signes. Sa lecture souffre, comme celle de
Durbin--Watson, du caractère conventionnel de l'ordre de parcours.
\end{fiche}

% -----------------------------------------------------------------------------
\clearpage
% =============================================================================
\section{Diagnostics M4 --- points aberrants du triangle}
\label{sec:mwm4}
% =============================================================================

\begin{encadre}
\textbf{Énoncé.} hors énoncé réglementaire~: conditionne la lecture de M1 à M3.
\end{encadre}

Les fiches ci-dessous suivent la structure en six rubriques employée pour les
hypothèses H1 à H4 de la méthode lognormale. Les notations sont celles de la
section~\ref{sec:notations-mw}.

\begin{fiche}{M4 --- Cellule aberrante du triangle (Grubbs)}
\label{mw:m4gr}
\refl{F.\,E. Grubbs (1950, 1969), \emph{op.\ cit.}}
\Cadre
Résidus de Mack, considérés comme un échantillon unique. Hypothèse non testée et
déterminante~: normalité des résidus sous $H_0$, laquelle n'est \emph{pas} une
hypothèse du modèle de Mack (voir M5).
\Hzero
$H_0$ : aucun résidu aberrant --- toutes les cellules proviennent du même modèle.
L'alternative est la présence d'exactement une cellule aberrante.
\Stat
$G = \max_{i,j}\big|r(i,j) - \bar r\big| / s_r$.
\Loi
Borne de Bonferroni fondée sur la loi de Student, conservatrice~; elle suppose la
normalité et l'indépendance, toutes deux approximatives ici. Loi de référence~:
bootstrap de résidus.
\Pval
p-value de Monte-Carlo retenue~; la borne de Bonferroni est fournie comme
majorant non simulé. \textbf{Ce majorant est extrêmement conservateur} sur des
résidus de Mack~: aucun rejet en 600 simulations sous $H_0$ au seuil de $10~\%$
(section~\ref{sec:calib-asympt-mw}). Il ne doit donc jamais servir de critère de
décision.
\Usage
Le moteur indique les coordonnées $(i,j)$ de la cellule la plus extrême, ce qui
rend le résultat directement actionnable~: il s'agit typiquement d'un sinistre
exceptionnel ou d'une erreur de saisie sur une année de développement donnée.
\textbf{Aucune exclusion de cellule ne doit être opérée sur ce seul fondement}~:
la justification doit être métier.
\end{fiche}

% -----------------------------------------------------------------------------
\begin{fiche}{M4 --- Cellules aberrantes multiples (ESD généralisé de Rosner)}
\label{mw:m4ros}
\refl{B.~Rosner (1983), \emph{op.\ cit.}}
\Cadre
Généralisation itérative de Grubbs, appliquée aux résidus de Mack, avec
$k = \lfloor n/4 \rfloor$ candidats au plus. Hypothèse non testée~: normalité.
\Hzero
$H_0$ : aucune cellule aberrante. On teste séquentiellement les alternatives
« il y a exactement $i$ cellules aberrantes », $i = 1,\dots,k$.
\Stat
$R_i = \max\big|r - \bar r_{\mathcal{S}_i}\big| / s_{\mathcal{S}_i}$ sur
l'échantillon privé des $i-1$ cellules déjà écartées.
\Loi
Valeurs critiques $\lambda_i$ tabulées, fondées sur la loi de Student.
\Pval
\textbf{Aucune}~: il s'agit d'une procédure de décision à niveau $\alpha$ fixé,
non d'un test produisant une p-value.
\Usage
Résout le problème de masquage de Grubbs lorsque plusieurs cellules sont
atypiques. La littérature recommande $n \ge 25$ pour la fiabilité des valeurs
critiques~; un triangle $8\times8$ fournit $28$ résidus, ce qui place la
procédure à la limite de son domaine. À traiter comme un signal d'attention.
\end{fiche}

% -----------------------------------------------------------------------------
\clearpage
% =============================================================================
\section{Diagnostics M5 --- normalité des résidus de Mack}
\label{sec:mwm5}
% =============================================================================

\begin{encadre}
\textbf{Énoncé.} \textbf{hors annexe~XVII}~: la normalité n'est pas une hypothèse du modèle, qui ne spécifie que les deux premiers moments conditionnels.
\end{encadre}

Les fiches ci-dessous suivent la structure en six rubriques employée pour les
hypothèses H1 à H4 de la méthode lognormale. Les notations sont celles de la
section~\ref{sec:notations-mw}.

\begin{fiche}{M5 --- Shapiro--Wilk sur les résidus de Mack}
\label{mw:m5sw}
\refl{S.\,S. Shapiro, M.\,B. Wilk (1965), \emph{op.\ cit.}~; loi nulle évaluée
par simulation directe.}
\Cadre
Résidus de Mack. \textbf{La normalité n'est pas une hypothèse du modèle}~: le
paragraphe~D(2)(h) ne spécifie que les deux premiers moments conditionnels. Ce
test est donc un \emph{diagnostic} et non la vérification d'une exigence
réglementaire.
\Hzero
$H_0$ : les résidus de Mack sont issus d'une loi normale.
\Stat
$W$, statistique de Shapiro--Wilk.
\Loi
$W$ étant invariante par translation et changement d'échelle, sa loi sous $H_0$
ne dépend que de l'effectif~: elle est évaluée numériquement par simulation
directe de 20\,000 échantillons normaux.
\Pval
p-value \emph{exacte à l'erreur de Monte-Carlo près}, en
$\mathcal{O}(B_{\text{null}}^{-1/2})$.
\Usage
\textbf{Point de méthode déterminant.} Ce test ne peut pas passer par le
bootstrap de résidus employé pour tous les autres tests de la méthode, car
celui-ci tire dans la loi \emph{empirique} des résidus observés~: son hypothèse
nulle serait « les résidus suivent leur propre loi empirique », ce qui rendrait
le test vide de sens (la simulation confirme une fréquence de rejet nulle au
lieu de 5~\%). D'où le recours à une loi nulle propre, indépendante du triangle
ajusté. Le résultat n'est pertinent que si l'on souhaite exploiter la MSEP pour
construire un quantile, usage qui excède la méthode réglementaire.
\end{fiche}

% -----------------------------------------------------------------------------
\begin{fiche}{M5 --- Lilliefors sur les résidus de Mack}
\label{mw:m5li}
\refl{H.\,W. Lilliefors (1967), \emph{op.\ cit.}~; p-value~: G.\,E. Dallal,
L.~Wilkinson (1986), \emph{op.\ cit.}}
\Cadre
Identique à la fiche précédente. Variante \emph{secondaire}~: masquée par défaut
dans l'application, conservée dans l'export.
\Hzero
$H_0$ : les résidus de Mack sont issus d'une loi normale, de moyenne et de
variance quelconques.
\Stat
$D_{\mathrm{L}} = \sup_u \big|\hat F(u) - \Phi((u - \bar r)/s_r)\big|$.
\Loi
Loi de Lilliefors, tenant compte de l'estimation de la moyenne et de
l'écart-type sur l'échantillon.
\Pval
Approximation analytique de Dallal \& Wilkinson, prolongée par les polynômes de
Stephens. p-value non simulée, disponible dès $n \ge 5$.
\Usage
Complément de recoupement à Shapiro--Wilk, moins puissant. Même réserve de fond~:
la normalité n'est pas exigée par la méthode.
\end{fiche}

% -----------------------------------------------------------------------------
\clearpage
% =============================================================================
\section{Diagnostics M6 --- robustesse de l'estimation}
\label{sec:mwm6}
% =============================================================================

\begin{encadre}
\textbf{Énoncé.} hors énoncé réglementaire~: mesure la concentration du résultat et la part de convention dans le calcul.
\end{encadre}

Les fiches ci-dessous suivent la structure en six rubriques employée pour les
hypothèses H1 à H4 de la méthode lognormale. Les notations sont celles de la
section~\ref{sec:notations-mw}.

\begin{fiche}{M6 --- Extrapolation de $\hat\sigma^2_{J-1}$ (diagnostic)}
\label{mw:m6sig}
\refl{Annexe XVII, section D, paragraphe~5(d)(ii), seconde ligne.}
\Cadre
La variance de la dernière année de développement ne peut être estimée, faute
d'observations. \textbf{Diagnostic, non test.}
\Hzero
Sans objet.
\Stat
La valeur retenue,
$\hat\sigma^2_{J-1} = \min\big(\hat\sigma^2_{J-2},\ \hat\sigma^2_{J-3},\
\hat\sigma^4_{J-2}/\hat\sigma^2_{J-3}\big)$.
\Loi
Sans objet~: la quantité n'est pas estimée mais \emph{imposée} par une règle
d'extrapolation.
\Pval
Aucune.
\Usage
Ce diagnostic existe pour rendre visible dans le dossier un point que le calcul
masque~: une des entrées de la MSEP ne provient pas des données mais d'une
convention réglementaire, reposant sur deux colonnes seulement. Comme la
dernière année de développement contribue à la MSEP de toutes les années de
survenance récentes --- celles qui portent l'essentiel de la réserve --- cette
convention a un effet non négligeable sur le paramètre final. Elle doit être
mentionnée explicitement.
\end{fiche}

% -----------------------------------------------------------------------------
\begin{fiche}{M6 --- Concentration de la réserve sur la dernière année}
\label{mw:m6conc}
\refl{Diagnostic de concentration~; décomposition du paragraphe~4.}
\Cadre
Part de la réserve totale portée par la dernière année de survenance,
$\big(\hat C(I,J) - C(I,0)\big) / \sum_i \big(\hat C(i,J) - C(i,I-i)\big)$.
\textbf{Diagnostic, non test.}
\Hzero
Sans objet.
\Stat
La part elle-même. Le moteur signale un dépassement de 40~\%.
\Loi
Sans objet.
\Pval
Aucune.
\Usage
La dernière année de survenance n'est observée que sur une seule cellule~: son
ultime résulte du produit de \emph{tous} les facteurs de développement, et son
erreur de prédiction domine la MSEP. Une part élevée signifie que le paramètre
propre repose principalement sur la ligne la moins informative du triangle. Ce
diagnostic joue, pour la méthode Merz--Wüthrich, le rôle que joue le jackknife
pour la méthode lognormale~: il mesure la concentration du résultat sur une
observation unique.
\end{fiche}

\newpage

\partie{Partie IV}{Validation, références et synthèse}

% =============================================================================
\section{Vérification empirique de l'implémentation}
\label{sec:verif}
% =============================================================================

L'implémentation a été validée par trois voies indépendantes.

\subsection*{Recalcul indépendant des statistiques}

Chaque statistique a été confrontée à un recalcul direct ou à une identité connue~: égalité
$F = t_b^2$ et égalité des p-values de Fisher et de Student en régression simple~; égalité de
$W'$ (Shapiro--Francia) et du carré du coefficient de corrélation du QQ-plot~; recalcul par
boucle explicite de $A^2$, $W^2$, $D$ et du $S$ de Mann--Kendall~; cohérence de $D$ avec
\code{ks.test}~; vérification de la borne $G \le (T-1)/\sqrt{T}$ pour Grubbs~; et vérification
numérique que la p-value bilatérale minimale de Cox--Stuart vaut exactement $0{,}0625$ pour
$T = 10$, ce qui confirme l'affirmation de la section~4.

\subsection*{Taille effective et puissance}

Sur 3\,000 échantillons normaux de taille 200, la fréquence de rejet au seuil nominal de 5~\%
ressort à 4,9~\% (D'Agostino), 4,5~\% (Anscombe--Glynn), 4,1~\% (Jarque--Bera) et 5,6~\%
(Wald--Wolfowitz)~; la borne de Bonferroni de Grubbs donne 4,6~\% pour $T=20$, conformément à son
caractère conservateur. Côté puissance, D'Agostino détecte l'asymétrie d'une loi exponentielle
dans 100~\% des cas et Anscombe--Glynn les queues d'une loi de Student à 4 degrés de liberté
dans 92,5~\% des cas.

\subsection*{Calibration des p-values de Monte-Carlo}

C'est le point le plus important, puisque ces p-values sont celles retenues pour la majorité des
tests --- mesuré sur le jeu de contrôle (méthode prime, $T = 8$, $B = 999$), 26 des 50 lignes de
la table auditable retiennent une p-value de Monte-Carlo, contre 9 une p-value exacte, 1 une
p-value quasi-exacte et 2 une p-value asymptotique, les 12 dernières étant des diagnostics ou
indicateurs sans p-value retenue. Sur 400 jeux de données simulés \emph{depuis le modèle} de
l'annexe~XVII, la fréquence de
rejet au seuil nominal de 10~\% a été mesurée pour chacune des statistiques bootstrapées.

\begin{center}
\small
\begin{tabular}{@{}lcc@{}}
\toprule
\textbf{Statistique} & \textbf{$T = 10$} & \textbf{$T = 20$} \\
\midrule
Anderson--Darling      & 0,092 & 0,113 \\
Cramér--von Mises      & 0,075 & 0,102 \\
Kolmogorov--Smirnov    & 0,092 & 0,115 \\
Shapiro--Wilk          & 0,090 & 0,110 \\
Shapiro--Francia       & 0,083 & 0,107 \\
Jarque--Bera           & 0,087 & 0,095 \\
Durbin--Watson         & 0,092 & 0,105 \\
Ljung--Box (retard 1)  & 0,095 & 0,090 \\
sup-F                  & 0,083 & 0,083 \\
OLS-CUSUM              & 0,075 & 0,087 \\
Grubbs                 & 0,090 & 0,085 \\
\bottomrule
\end{tabular}
\end{center}

Toutes les valeurs sont compatibles avec le niveau nominal de 10~\% à l'intérieur de
l'intervalle de confiance à 95~\% ($\pm 0{,}029$ pour 400 réplications). La procédure de
bootstrap paramétrique est donc correctement calibrée, y compris à $T = 10$ où aucune des lois
asymptotiques utilisées par les tests classiques ne l'est.

\medskip
\begin{encadre}
\textbf{Réserve.} Cette validation est conditionnelle au modèle~: elle établit que, \emph{si} les
données sont engendrées par le modèle de l'annexe~XVII, les p-values de Monte-Carlo ont le bon
niveau. Elle ne dit rien de la puissance du dispositif contre des écarts réalistes au modèle,
laquelle reste faible à ces tailles d'échantillon.
\end{encadre}

\newpage
\subsection{Vérification de l'implémentation Merz--Wüthrich}

Le triangle de Taylor \& Ashe (1983), référence classique de la littérature, sert
de jeu de contrôle.

\begin{center}
\small
\rowcolors{2}{tablerow}{white}
\begin{tabular}{@{}p{6.9cm}p{4.1cm}p{3.9cm}@{}}
\rowcolor{tableheader}
\textcolor{financeblue}{\textbf{Contrôle}} &
\textcolor{financeblue}{\textbf{Résultat}} &
\textcolor{financeblue}{\textbf{Statut}} \\
Facteurs de développement $\hat f_j$ & 3,4906 ; 1,7473 ; 1,4574 ; \dots &
conformes aux valeurs publiées \\
Réserve chain-ladder totale & 18\,680\,856 & conforme \\
Erreur type \emph{ultime} de Mack & 2\,447\,095 & conforme à la valeur publiée \\
MSEP à un an $=$ transcription littérale du paragraphe~D(5) & écart relatif nul
(précision machine) & conforme au texte \\
$\sqrt{\mathrm{MSEP}}$ à un an $=$ \code{ChainLadder::CDR} (Merz--Wüthrich
2008) & 1\,778\,967,66336\newline contre 1\,778\,967,66335758 & écart relatif
$-2{,}6\times10^{-15}$ \\
Erreur à un an $<$ erreur ultime & 1\,778\,968 $<$ 2\,447\,095 &
conforme à la théorie \\
Homogénéité~: $\lambda\!\cdot\!C \Rightarrow$ CV invariant & vérifié à
$10^{-10}$ & \\
Triangle déterministe $\Rightarrow \mathrm{MSEP} = 0$ &
$2{,}7\times10^{-18}$ & \\
Moments de $Z$ (test calendaire) par énumération exhaustive & $n = 2,\dots,12$ &
exacts \\
\end{tabular}
\end{center}

\noindent
La conformité de l'erreur type ultime de Mack à sa valeur publiée valide
simultanément $\hat f_j$, $\hat\sigma_j^2$ et la règle d'extrapolation de
$\hat\sigma_{J-1}^2$, qui sont les entrées de la MSEP à un an.

\medskip
\begin{encadre}
\textbf{Sur quoi repose la validation de la MSEP à un an.} Elle repose sur deux
voies distinctes~: d'une part la confrontation du moteur à une transcription
littérale du paragraphe~D(5), réécrite depuis le texte et implémentée
séparément dans la batterie de tests unitaires~; d'autre part une comparaison
\emph{externe} avec \code{ChainLadder::CDR} (paquet R \code{ChainLadder}
0.2.18), implémentation publique de Merz--Wüthrich (2008), reproduite à
$-2{,}6\times10^{-15}$ en écart relatif. Nous n'avons pas identifié de valeur
\emph{publiée} de la MSEP à un an sur ce triangle~: la comparaison externe porte
sur une implémentation de référence, non sur un nombre publié.

\medskip
\textbf{Ce que les versions antérieures de ce document affirmaient à tort.}
Jusqu'à la correction de la transcription du paragraphe~D(5)
(section~\ref{sec:mw}), la validation de cette grandeur était présentée comme
reposant sur une « double implémentation indépendante ». Elle ne l'était
pas~: les deux implémentations reprenaient la même transcription fautive du
paragraphe~D(5) --- termes diagonaux $\hat C(i,J)^2\,\Delta_i$ et facteur~2
omis --- et ne pouvaient donc rien détecter. La MSEP en était sous-estimée
(1\,519\,876 au lieu de 1\,778\,968 sur ce triangle), dans le sens favorable à
l'entreprise. C'est ce défaut que la transcription de la
section~\ref{sec:mw} et les deux voies ci-dessus corrigent.
\end{encadre}

\subsection{Calibration des p-values asymptotiques (Merz--Wüthrich)}
\label{sec:calib-asympt-mw}

Les p-values non simulées de la méthode Merz--Wüthrich sont rapportées à titre de
complément, la p-value retenue étant celle du bootstrap. Il importe néanmoins de
savoir \emph{de combien} elles s'en écartent, et pourquoi. Fréquences de rejet
mesurées sur 600 triangles $10\times10$ simulés exactement sous $H_0$, au seuil
nominal de $10~\%$ (intervalle de confiance à 95~\%~: $\pm 0{,}024$)~:

\begin{center}
\small
\rowcolors{2}{tablerow}{white}
\begin{tabular}{@{}p{5.1cm}p{1.9cm}p{8.0cm}@{}}
\rowcolor{tableheader}
\textcolor{financeblue}{\textbf{Test}} &
\textcolor{financeblue}{\textbf{Rejets à 10 \%}} &
\textcolor{financeblue}{\textbf{Interprétation}} \\
Ordonnée à l'origine (Fisher) & 0,092 & Conforme \\
Homogénéité de $f_j$ (Fisher) & 0,125 & Léger sur-rejet, dans l'incertitude \\
Courbure (Fisher) & 0,112 & Conforme \\
Exposant de variance (Fisher) & 0,097 & Conforme \\
Kruskal--Wallis & 0,135 & Sur-rejet~: groupes très déséquilibrés \\
Effets d'année calendaire & 0,123 & Sur-rejet~: approximation normale d'une somme
portant sur peu de diagonales \\
Test des suites & \textbf{0,153} & \textbf{Sur-rejet marqué.} La contrainte
$\sum_i \sqrt{C(i,j)}\,r(i,j) = 0$ rend les résidus d'une colonne négativement
corrélés, d'où une alternance des signes plus fréquente que sous
l'échangeabilité supposée par l'approximation normale \\
Grubbs (borne de Bonferroni) & \textbf{0,000} & \textbf{Extrêmement
conservatrice.} La borne est un majorant, et les résidus étant déjà standardisés
par colonne, l'écart au maximum est faible \\
Breusch--Pagan ($\chi^2$) & \textbf{0,005} & \textbf{Très conservatrice.}
$\hat\sigma_j^2$ est estimé sur la colonne même~; les résidus standardisés ont
une dispersion réduite, donc un $R^2$ auxiliaire trop faible \\
\end{tabular}
\end{center}

\begin{encadre}
\textbf{Conséquence pratique.} Les écarts observés entre p-value asymptotique et
p-value de Monte-Carlo ne traduisent pas une erreur de calcul mais la
\emph{non-validité} des lois de référence classiques sur des résidus de triangle,
qui ne sont ni indépendants ni échangeables. Trois tests s'en écartent fortement
et dans des sens opposés~: le test des suites sur-rejette, Grubbs et
Breusch--Pagan sous-rejettent massivement. Dans les trois cas, c'est la p-value
de Monte-Carlo qui est retenue par le moteur, et elle seule doit être utilisée en
dossier. La p-value non simulée n'est conservée que pour documenter l'écart.
\end{encadre}

\subsection{Calibration des p-values de Monte-Carlo (Merz--Wüthrich)}

Sur 160 triangles $10\times10$ simulés sous le modèle de Mack, fréquence de rejet
au seuil nominal de 10~\% (intervalle de confiance à 95~\%~: $\pm 0{,}046$)~:

\begin{center}
\small
\rowcolors{2}{tablerow}{white}
\begin{tabular}{@{}p{6.4cm}cc@{}}
\rowcolor{tableheader}
\textcolor{financeblue}{\textbf{Statistique}} &
\textcolor{financeblue}{\textbf{rejets à 10 \%}} &
\textcolor{financeblue}{\textbf{rejets à 5 \%}} \\
Effets d'année calendaire & 0,062 & 0,019 \\
Corrélation entre années de développement & 0,119 & 0,037 \\
Breusch--Pagan sur résidus de Mack & 0,100 & 0,062 \\
Grubbs & 0,113 & 0,044 \\
Durbin--Watson & 0,100 & 0,050 \\
Test des suites & 0,056 & 0,025 \\
Proportionnalité (pente) & 0,069 & 0,044 \\
\end{tabular}
\end{center}

\noindent
Toutes les valeurs sont compatibles avec les niveaux nominaux à l'intérieur de
l'intervalle de confiance.



\clearpage
% =============================================================================
\section{Compléments bibliographiques}
\label{sec:biblio}
% =============================================================================

Les fiches de la partie~III citent les articles \emph{fondateurs} de chaque
test~: ils établissent la statistique et sa loi. Les affirmations relatives à la
\emph{qualité de l'approximation à $T = 8$} et à la \emph{validité du bootstrap}
relèvent en revanche d'une littérature distincte, rassemblée ici. Chaque entrée
précise la portée exacte de la référence, c'est-à-dire l'affirmation qu'elle
soutient dans le présent document.

\subsection*{Vitesse de convergence et bornes d'erreur}

\begin{itemize}[itemsep=4pt]
  \item A.\,C. Berry, \emph{The accuracy of the Gaussian approximation to the
  sum of independent variates}, \emph{Transactions of the American Mathematical
  Society}, 49(1), 1941, p.~122--136 ; C.-G. Esseen, \emph{On the Liapounoff
  limit of error in the theory of probability}, \emph{Arkiv för Matematik,
  Astronomi och Fysik}, 28A(9), 1942, p.~1--19.
  \emph{Portée.} Les fiches citaient jusqu'ici uniquement les
  auteurs des tests. Le théorème de Berry--Esseen est la référence qui permet
  d'énoncer une vitesse $\mathcal{O}(T^{-1/2})$ pour les approximations
  normales, et donc de montrer qu'à $T = 8$ la borne est de l'ordre de
  $0{,}35$, c'est-à-dire sans valeur pratique. Aucune référence antérieure du
  document ne supportait cette affirmation.

  \item H.~Callaert, P.~Janssen, \emph{The Berry-Esseen theorem for
  $U$-statistics}, \emph{The Annals of Statistics}, 6(2), 1978, p.~417--421.
  \emph{Portée.} Le $S$ de Mann--Kendall est une $U$-statistique et
  non une somme de variables indépendantes~: le Berry--Esseen classique ne
  s'applique pas directement. Cette référence est celle qui couvre exactement le
  cas rencontré. La référence existante (Mann, 1945) établit la statistique,
  pas sa vitesse de convergence.
\end{itemize}

\subsection*{Lois exactes employées à la place des approximations normales}

\begin{itemize}[itemsep=4pt]
  \item M.\,G. Kendall, J.\,D. Gibbons, \emph{Rank Correlation Methods},
  5\textsuperscript{e} éd., Edward Arnold / Oxford University Press, 1990.
  \emph{Portée.} Le document citait Kendall (1975,
  4\textsuperscript{e} éd.). L'édition de 1990 est celle qui est aujourd'hui
  la référence courante pour la distribution exacte de $S$ et ses tables. Les
  deux éditions sont conservées~: la 4\textsuperscript{e} pour l'antériorité de
  la méthode, la 5\textsuperscript{e} pour le résultat exact effectivement
  implémenté.

  \item F.~Swed, C.~Eisenhart, \emph{Tables for testing randomness of grouping
  in a sequence of alternatives}, \emph{The Annals of Mathematical Statistics},
  14(1), 1943, p.~66--87.
  \emph{Portée.} Wald \& Wolfowitz (1940), déjà cité, établit le
  test~; Swed \& Eisenhart fournit la distribution combinatoire exacte du
  nombre de suites, qui est celle implémentée dans le moteur, l'approximation
  normale n'étant conservée qu'à titre indicatif.

  \item D.\,J. Best, D.\,E. Roberts, \emph{Algorithm AS 89: The upper tail
  probabilities of Spearman's $\rho$}, \emph{Journal of the Royal Statistical
  Society, Series C (Applied Statistics)}, 24(3), 1975, p.~377--379.
  \emph{Portée.} La référence Spearman (1904) établit le
  coefficient. Best \& Roberts documente le calcul des probabilités exactes en
  petit échantillon, qui est ce que \code{cor.test(..., exact = TRUE)} met en
  \oe uvre et que le moteur retient à $T = 8$.
\end{itemize}

\subsection*{Théorie du bootstrap et des tests de Monte-Carlo}

\begin{itemize}[itemsep=4pt]
  \item M.~Dwass, \emph{Modified randomization tests for nonparametric
  hypotheses}, \emph{The Annals of Mathematical Statistics}, 28(1), 1957,
  p.~181--187 ; A.\,C.\,A. Hope, \emph{A simplified Monte Carlo significance
  test procedure}, \emph{Journal of the Royal Statistical Society, Series B},
  30(3), 1968, p.~582--598.
  \emph{Portée.} Ce sont les références qui établissent que le test
  de Monte-Carlo fondé sur $\hat p = (1 + \#\{S^{(b)} \ge S_{\text{obs}}\})/(B+1)$
  est de niveau \emph{exact} pour $B$ fini lorsque $\alpha(B+1)$ est entier. Le
  document invoquait cette propriété en citant Efron (1979), qui ne la
  démontre pas. C'est une correction de référence, non un simple ajout.

  \item P.\,H. Jöckel, \emph{Finite sample properties and asymptotic efficiency
  of Monte Carlo tests}, \emph{The Annals of Statistics}, 14(1), 1986,
  p.~336--347 ; R.\,A. Marriott, \emph{Barnard's Monte Carlo tests: how many
  simulations?}, \emph{Journal of the Royal Statistical Society, Series C},
  28(1), 1979, p.~75--77.
  \emph{Portée.} Ces travaux quantifient la perte de puissance
  induite par un $B$ fini, et fondent le choix d'un $B$ de l'ordre du millier.
  Ils justifient la ligne (B) du tableau~\ref{tab:conv}.

  \item P.\,J. Bickel, D.\,A. Freedman, \emph{Some asymptotic theory for the
  bootstrap}, \emph{The Annals of Statistics}, 9(6), 1981, p.~1196--1217 ;
  K.~Singh, \emph{On the asymptotic accuracy of Efron's bootstrap},
  \emph{The Annals of Statistics}, 9(6), 1981, p.~1187--1195.
  \emph{Portée.} Efron (1979), seule référence bootstrap du
  document, introduit la méthode mais n'en établit pas la consistance. Ces deux
  articles sont ceux qui démontrent la validité asymptotique, et donc ce qui
  fonde --- et ce qui limite --- l'emploi du bootstrap ici~: la garantie est
  \emph{asymptotique en $T$}, ce qui est précisément le point à signaler à
  $T = 8$.

  \item P.~Hall, \emph{The Bootstrap and Edgeworth Expansion}, Springer, 1992 ;
  A.\,C. Davison, D.\,V. Hinkley, \emph{Bootstrap Methods and their
  Application}, Cambridge University Press, 1997.
  \emph{Portée.} Hall (1992) est l'ouvrage de référence sur les
  raffinements d'ordre supérieur obtenus pour les statistiques pivotales, qui
  explique pourquoi le bootstrap améliore l'approximation sans la rendre exacte.
  Davison \& Hinkley (1997) traite spécifiquement du bootstrap paramétrique et
  des tests de Monte-Carlo tels qu'employés ici.
\end{itemize}

\subsection*{Comportement en petit échantillon des tests employés}

\begin{itemize}[itemsep=4pt]
  \item N.~Davies, C.\,M. Triggs, P.~Newbold, \emph{Significance levels of the
  Box-Pierce portmanteau statistic in finite samples}, \emph{Biometrika}, 64(3),
  1977, p.~517--522.
  \emph{Portée.} Le document affirmait que l'approximation
  $\chi^2$ des tests portmanteau est mal calibrée en petit échantillon en
  citant Ljung \& Box (1978). Davies, Triggs \& Newbold est l'étude dédiée à
  cette question et constitue la référence appropriée. Ljung \& Box reste cité
  pour la statistique elle-même.

  \item K.\,O. Bowman, L.\,R. Shenton, \emph{Omnibus test contours for
  departures from normality based on $\sqrt{b_1}$ and $b_2$},
  \emph{Biometrika}, 62(2), 1975, p.~243--250 ; C.\,M. Jarque, A.\,K. Bera,
  \emph{A test for normality of observations and regression residuals},
  \emph{International Statistical Review}, 55(2), 1987, p.~163--172.
  \emph{Portée.} Bowman \& Shenton documente la lenteur de
  convergence des statistiques fondées sur $\sqrt{b_1}$ et $b_2$ vers leur loi
  limite, ce qui est le fondement de la mise en garde du
  tableau~\ref{tab:nature} sur Jarque--Bera. L'article de 1987 est la version
  complète du test, la note de 1980 dans \emph{Economics Letters} en étant la
  forme abrégée~: les deux sont conservées.

  \item S.\,S. Shapiro, M.\,B. Wilk, H.\,J. Chen, \emph{A comparative study of
  various tests for normality}, \emph{Journal of the American Statistical
  Association}, 63(324), 1968, p.~1343--1372.
  \emph{Portée.} Étude de puissance de référence pour les tests de
  normalité, qui documente leur faiblesse en très petit échantillon. Elle
  soutient l'affirmation, présente dans les fiches, selon laquelle une p-value
  élevée n'apporte quasiment aucune information à $T$ faible.

  \item G.\,S. Maddala, \emph{Introduction to Econometrics}, 3\textsuperscript{e}
  éd., Wiley, 2001, pour la discussion du caractère non exact du test RESET.
  \emph{Portée.} Ramsey (1969), seule référence citée jusqu'ici,
  développe le test à partir de résidus BLUS. Dans la forme usuellement
  implémentée --- et retenue ici --- les régresseurs auxiliaires sont des
  puissances des valeurs ajustées, qui dépendent de $y$~: la statistique n'est
  alors plus exactement de loi de Fisher. Cette nuance, absente de la version
  précédente de la documentation, motive le recours au bootstrap.

  \item J.\,P. Imhof, \emph{Computing the distribution of quadratic forms in
  normal variables}, \emph{Biometrika}, 48(3/4), 1961, p.~419--426.
  \emph{Portée.} Référence exacte de la méthode évoquée dans la
  fiche Durbin--Watson pour le calcul de la loi exacte de la statistique, dont
  le document mentionnait le nom sans la référencer.
\end{itemize}


\subsection{Traçabilité entre la documentation et le code}

\noindent
L'architecture complète du moteur est décrite en section~\ref{sec:archi}.
Le tableau ci-dessous relie chaque développement du présent document à son
implémentation.


Chaque affirmation quantitative de ce document correspond à une fonction
identifiable du moteur \code{R/engine.R}, seul fichier contenant de la logique
statistique ou actuarielle.

\begin{center}
\small
\begin{tabular}{@{}p{5.6cm}p{9.4cm}@{}}
\toprule
\textbf{Élément documenté} & \textbf{Implémentation dans \code{R/engine.R}} \\
\midrule
Noyau du modèle ($\pi_t$, $\hat\beta$, $z_t$) & \code{usp\_pi()}, \code{usp\_noyau()} \\
Optimisation sous contrainte $\delta\in[0,1]$ & \code{usp\_objectif()}, \code{usp\_ajuster()} \\
Lois exactes (Kendall, suites) & \code{mk\_p\_exacte()}, \code{runs\_p\_exacte()} \\
Statistiques de test & \code{stat\_*()}, \code{test\_*()} \\
Bootstrap paramétrique et p-values MC & \code{.stats\_bootstrapables()}, \code{usp\_bootstrap()} \\
Hiérarchie exacte $\succ$ MC $\succ$ asymptotique & fonction \code{add()} interne à \code{usp\_tests()} \\
Sortie de la hiérarchie (grandeur rivée par l'estimation) & branche finale de \code{add()}~:
\code{p\_retenue} et \code{nature\_p} à \code{NA}, verdict \code{INFO} \\
Réserve sur la base « ratios bruts » à $\hat\pi_t$ constant & champ \code{detail} des six lignes
de base \code{r} dans \code{usp\_tests()} \\
Erreur de Monte-Carlo & champ \code{err\_mc} de \code{usp\_bootstrap()} \\
Jackknife, profils, rapport de vraisemblance & \code{usp\_jackknife()}, \code{usp\_profil()} \\
Calibration et paramètre final & \code{usp\_parametre()}, \code{usp\_credibilite()} \\
Contrôles de validité & \code{engine\_valider\_donnees()}, \code{usp\_controle\_donnees()} \\
Quantités des graphiques & \code{engine\_plots\_data()} \\
Orchestration complète & \code{run\_engine()} \\
Table auditable des tests & \code{engine\_table\_tests()} \\
\bottomrule
\end{tabular}
\end{center}

\noindent
L'application Shiny (\code{app.R}, \code{R/display\_helpers.R}) ne contient
aucun de ces calculs~: elle collecte les entrées, appelle \code{run\_engine()} et
affiche l'objet retourné. Le moteur s'exécute indépendamment de toute interface,
ce qui permet une revue et une reproduction externes~:
\code{source("R/engine.R")} puis \code{run\_engine(xt = ..., yt = ...)}.

\subsection*{Méthode Merz--Wüthrich}

Les références ci-dessous fondent la section~\ref{sec:mw} et les fiches M1 à M6.
Elles n'ont pas d'équivalent dans le corpus lognormal, la méthode reposant sur un
modèle et un mode d'inférence différents.

\begin{itemize}[itemsep=4pt]
  \item T.~Mack, \emph{Distribution-free calculation of the standard error of
  chain ladder reserve estimates}, \emph{ASTIN Bulletin}, 23(2), 1993,
  p.~213--225.
  \emph{Portée.} Article fondateur du modèle de chaîne d'échelle stochastique~:
  il établit les estimateurs $\hat f_j$ et $\hat\sigma_j^2$, l'erreur type
  \emph{ultime}, et surtout l'énoncé des trois hypothèses reprises au
  paragraphe~D(2)(h) du règlement. C'est la valeur publiée de l'erreur ultime sur
  le triangle de Taylor \& Ashe qui a servi à valider notre implémentation de
  $\hat f_j$, $\hat\sigma_j^2$ et de la règle d'extrapolation.

  \item T.~Mack, \emph{Which stochastic model is underlying the chain ladder
  method~?}, \emph{Insurance: Mathematics and Economics}, 15(2--3), 1994,
  p.~133--138.
  \emph{Portée.} Source du test des effets d'année calendaire et de la
  discussion des hypothèses testables. C'est de cet article que proviennent les
  moments de $Z_k = \min(L_k, S_k)$, que nous avons vérifiés par énumération
  exhaustive.

  \item M.~Merz, M.\,V. Wüthrich, \emph{Modelling the claims development result
  for solvency purposes}, \emph{Casualty Actuarial Society E-Forum}, automne
  2008, p.~542--568.
  \emph{Portée.} Article dont dérive la formule d'erreur de prédiction \emph{à un
  an} du paragraphe~D(5). La distinction entre erreur à un an et erreur ultime
  est le fondement même de la méthode~n\textsuperscript{o}~2, et l'ordre
  $\sqrt{\mathrm{MSEP}}_{\text{1 an}} < \sqrt{\mathrm{MSEP}}_{\text{ultime}}$ que
  nous avons vérifié numériquement en découle. C'est également l'implémentation
  publique de cet article (\code{ChainLadder::CDR}) qui fournit la valeur de
  comparaison externe de notre MSEP à un an (section~\ref{sec:verif}).

  \item M.\,V. Wüthrich, M.~Merz, \emph{Stochastic Claims Reserving Methods in
  Insurance}, Wiley, 2008.
  \emph{Portée.} Ouvrage de référence rassemblant le cadre théorique complet
  (modèle de Mack, résultat de développement des sinistres, erreurs de
  prédiction). Il est cité pour le cadre général, l'article de 2008 l'étant pour
  la formule précise.

  \item P.\,D. England, R.\,J. Verrall, \emph{Stochastic claims reserving in
  general insurance}, \emph{British Actuarial Journal}, 8(3), 2002, p.~443--518.
  \emph{Portée.} Référence du bootstrap par rééchantillonnage des résidus, seul
  possible ici puisque le modèle de Mack ne spécifie aucune loi. Elle fonde la
  nature \emph{semi-paramétrique} des p-values de Monte-Carlo de cette méthode,
  et donc la réserve exposée en section~\ref{sec:boot-mw}.

  \item G.\,C. Taylor, F.\,R. Ashe, \emph{Second moments of estimates of
  outstanding claims}, \emph{Journal of Econometrics}, 23(1), 1983, p.~37--61.
  \emph{Portée.} Source du triangle de dix années de survenance utilisé comme jeu
  de contrôle. Son statut de référence partagée par la littérature est ce qui
  rend la validation reproductible par un tiers.

  \item G.~Barnett, B.~Zehnwirth, \emph{Best estimates for reserves},
  \emph{Proceedings of the Casualty Actuarial Society}, 87, 2000, p.~245--321.
  \emph{Portée.} Discussion des exposants alternatifs de la structure de
  variance ($\alpha = 0$, $1$ ou $2$). Citée dans la fiche M2 pour situer
  l'exposant~1 imposé par le règlement parmi les choix possibles, et non pour
  recommander une alternative --- l'annexe~XVII ne laissant pas ce choix ouvert.
\end{itemize}

\subsection*{Points sur lesquels la littérature ne permet pas de conclure}

Conformément à l'exigence de ne pas produire de conclusion artificiellement
précise, les points suivants sont signalés comme \emph{non tranchés} par la
littérature consultée.

\begin{enumerate}[itemsep=3pt]
  \item \textbf{Aucune borne d'erreur exploitable à $T = 8$} n'a pu être
  identifiée pour les statistiques de type multiplicateur de Lagrange
  (Breusch--Pagan, White) ni pour les statistiques portmanteau. Les résultats
  disponibles sont soit purement asymptotiques, soit des études de simulation
  portant sur des tailles supérieures. L'affirmation « la loi $\chi^2$ est
  inadéquate à $T = 8$ » est donc fondée sur l'absence de justification, non sur
  une borne quantifiée.
  \item \textbf{La vitesse de convergence du bootstrap paramétrique dans ce
  modèle} n'est pas établie. Les résultats de Bickel \& Freedman (1981), Singh
  (1981) et Hall (1992) portent sur des cadres où l'estimateur est
  asymptotiquement normal et le paramètre intérieur au domaine. Ici,
  $\hat\delta$ est fréquemment sur la frontière $\{0,1\}$, situation qui sort de
  ce cadre. Aucun taux ne peut donc être annoncé.
  \item \textbf{La loi exacte de la statistique agrégée $Z$} du test des effets
  d'année calendaire n'est pas établie. Les moments de chaque $Z_k$ sont exacts,
  mais la normalité de leur somme sur un nombre réduit de diagonales n'est
  justifiée par aucun résultat que nous ayons identifié~; d'où le recours au
  bootstrap.
  \item \textbf{Aucune valeur publiée de la MSEP à un an} n'a pu être identifiée
  pour le triangle de Taylor \& Ashe. La validation de cette grandeur repose donc
  sur la transcription littérale du paragraphe~D(5), sur des propriétés
  structurelles et sur la comparaison avec une implémentation externe
  (\code{ChainLadder::CDR}), non sur un nombre publié dans la littérature.
  \item \textbf{Le domaine de validité exact de la transformation de Royston}
  pour $4 \le T \le 11$ est documenté par son auteur, mais nous n'avons pas
  identifié d'étude indépendante en évaluant la précision spécifiquement à
  $T = 8$.
\end{enumerate}

\medskip
\begin{encadre}
\textbf{Articulation des deux corpus.} Les références des fiches établissent
\emph{ce que calcule} chaque test~; celles de la présente section établissent
\emph{ce que vaut} le résultat à la profondeur d'intérêt. Les secondes portent
sur les vitesses de convergence, les lois exactes effectivement implémentées et
les fondements théoriques du bootstrap et des tests de Monte-Carlo.
\end{encadre}


\newpage
\clearpage
% =============================================================================
\section{Synthèse et lecture des verdicts}
% =============================================================================

\begin{center}
\small
\begin{tabular}{@{}p{2.1cm}p{13.2cm}@{}}
\toprule
\textbf{Verdict} & \textbf{Signification} \\
\midrule
\code{OK} & p-value \emph{retenue} --- celle que désigne \code{nature\_p}, choisie dans l'ordre
exacte $\succ$ Monte-Carlo $\succ$ asymptotique --- au-dessus du seuil d'alerte, ou diagnostic
sans anomalie. Pour les tests de Student sur la pente et de Fisher, la logique est inversée~:
\code{OK} correspond à un $p$ \emph{faible}. \\
\code{ALERTE} & p-value retenue dans la zone intermédiaire $[\alpha/2\,;\,\alpha]$ (par défaut
$[5\,\%\,;\,10\,\%]$), ou diagnostic structurel à documenter~: solution au bord sur $\delta$,
point influent isolé, jackknife entre 10~\% et 20~\%, intervalle bootstrap large. \\
\code{ECHEC} & p-value retenue sous le seuil strict, ou anomalie structurelle marquée. N'invalide pas
mécaniquement l'USP, mais impose une justification explicite~: nature de l'écart, retraitement
envisagé, sensibilité du paramètre final. \\
\code{INFO} & Diagnostic sans seuil statistique associé~: test non applicable faute d'effectif
($T < 8$ pour D'Agostino, $T < 20$ pour Anscombe--Glynn), grandeur rivée par l'estimation
(fiches~\ref{fiche:32} et~\ref{fiche:33}), indicateur descriptif, ou lecture graphique. \\
\bottomrule
\end{tabular}
\end{center}

\subsection*{Hiérarchie de lecture recommandée}

\begin{enumerate}[itemsep=2pt]
  \item \textbf{Contrôles de données} (section~\ref{sec:controles}) -- bloquants. Aucun résultat
        n'a de sens si un contrôle est en échec.
  \item \textbf{Convergence numérique} -- préalable technique~: part des démarrages convergents
        (fiche~\ref{fiche:43}) et écart de $\sum_t z_t^2$ à $T$ (fiche~\ref{fiche:33}). La
        nullité de la moyenne \emph{pondérée} des résidus n'en fait pas partie~: c'est une
        identité algébrique, vraie même sans convergence (fiche~\ref{fiche:31}). La moyenne
        simple $\bar z$ est, elle, un diagnostic (fiche~\ref{fiche:32}).
  \item \textbf{Jackknife et intervalle bootstrap} (section~\ref{sec:robustesse}) -- la mesure la
        plus directe de la solidité du calibrage à $T$ faible.
  \item \textbf{Position de $\hat\delta$ et test du rapport de vraisemblance}
        (fiches~\ref{fiche:17} et~\ref{fiche:40}) -- identifiabilité de la structure de variance.
  \item \textbf{Ruptures et points influents} (section~\ref{sec:stabilite}) -- susceptibles
        d'expliquer tout rejet observé par ailleurs.
  \item \textbf{Tests des hypothèses H1 à H4} (sections~\ref{sec:h1} à~\ref{sec:h4}), en lisant
        pour chacun la p-value \emph{retenue} par le moteur, celle que désigne le champ
        \code{nature\_p}. La hiérarchie appliquée est \emph{exacte $\succ$ Monte-Carlo
        $\succ$ asymptotique} (section~\ref{sec:nature})~: la p-value de Monte-Carlo n'est
        pas la première servie, elle vient \emph{après} la p-value exacte lorsque celle-ci
        existe. Mesuré sur le jeu de contrôle (méthode prime, $T = 8$, $B = 999$)~: des 50
        lignes de la table auditable, 9 retiennent une p-value exacte, et 8 d'entre
        elles disposent pourtant d'une \code{p\_mc} qui n'est pas retenue~: lire
        systématiquement \code{p\_mc} reviendrait à écarter huit p-values exactes au profit
        d'estimations entachées d'une erreur de Monte-Carlo (section~\ref{sec:deuxerreurs}).
        Les trois colonnes restent affichées~: leur comparaison est un contrôle utile, mais
        elle ne modifie pas le verdict.
\end{enumerate}

\subsection*{Tests structurellement inopérants à faible $T$}

\begin{center}
\small
\begin{tabular}{@{}p{4.4cm}p{2.6cm}p{8.0cm}@{}}
\toprule
\textbf{Test} & \textbf{Seuil de validité} & \textbf{Nature du problème} \\
\midrule
Cox--Stuart & $T \ge 12$ & Rejet à 5~\% \emph{impossible} pour $m \le 5$ paires. \\
Anscombe--Glynn & $T \ge 20$ & Non défini en-deçà~; marqué \code{INFO}. \\
D'Agostino (asymétrie) & $T \ge 8$ & Non défini en-deçà~; puissance faible jusqu'à $T \approx 20$. \\
Jarque--Bera & $T \ge 50$ & $b_2$ borné par $\approx T$~; convergence $\chi^2$ très lente. \\
Ljung--Box / Box--Pierce & $T \ge 30$ & Approximation $\chi^2$ mal calibrée. \\
Wald--Wolfowitz & $n_1,n_2 \ge 10$ & Approximation normale~; granularité de $R$. \\
Breusch--Pagan / White & $T \ge 30$ & Convergence asymptotique lente. \\
Rosner ESD & $T \ge 25$ & Valeurs critiques peu fiables en-deçà. \\
Durbin--Watson & --- & Bornes invalides sur résidus de maximum de vraisemblance. \\
sup-F de Quandt & --- & Loi de Fisher inapplicable (point de rupture estimé). \\
\bottomrule
\end{tabular}
\end{center}

\vspace{1em}
\begin{encadre}
\textbf{Rappel final.} Le décompte des tests en échec ne se substitue pas à l'analyse qualitative
exigée par le règlement d'exécution (UE)~2015/498~: la note de justification du dossier de
candidature doit expliquer, segment par segment, en quoi la méthode retenue est adaptée au profil
de risque de l'organisme, y compris pour les segments non retenus. Aucun test isolé, ni aucun
comptage de tests, ne peut établir cette démonstration.
\end{encadre}


\partie{Annexe}{Flux détaillés entre fonctions}

\clearpage
% =============================================================================
\section{Annexe --- flux détaillés entre fonctions}
\label{sec:annexe-flux}
% =============================================================================

La figure~\ref{fig:archi} donne la vue d'ensemble de la chaîne de calcul. La
présente annexe descend au niveau de la fonction individuelle. Elle poursuit un
objectif précis~: rendre visible, pour chaque test, \emph{combien} de fonctions
concourent à son résultat et \emph{lesquelles} sont partagées avec d'autres
tests. Cette information conditionne la portée d'une revue de code~: une
anomalie dans une fonction mutualisée se propage à plusieurs tests, alors qu'une
anomalie dans une fonction dédiée reste circonscrite.

\subsection{Convention graphique}

\begin{center}
\small
\rowcolors{2}{tablerow}{white}
\begin{tabular}{@{}p{4.4cm}p{10.6cm}@{}}
\rowcolor{tableheader}
\textcolor{financeblue}{\textbf{Forme ou couleur}} &
\textcolor{financeblue}{\textbf{Signification}} \\
Bloc bleu clair & Entrée ou donnée intermédiaire \\
Bloc blanc, contour bleu & Fonction \textbf{dédiée} à un seul test \\
Bloc orangé, contour épais & Fonction \textbf{mutualisée} entre plusieurs tests
ou entre les deux méthodes \\
Bloc gris & Fonction \textbf{interne} de support, préfixée d'un point, jamais
appelée directement par \code{usp\_tests()} ou \code{mw\_tests()} \\
Bloc vert & Sortie : p-value ou grandeur restituée \\
Flèche pleine & Transmission d'un objet, dont le nom figure sur la flèche \\
Flèche pointillée & Dépendance de second rang (mise en cache, loi de référence) \\
\end{tabular}
\end{center}

\medskip
\noindent
Les fonctions mutualisées appellent une vigilance particulière. Le tableau
suivant les recense.

\begin{center}
\small
\rowcolors{2}{tablerow}{white}
\begin{tabular}{@{}p{4.4cm}p{1.9cm}p{8.5cm}@{}}
\rowcolor{tableheader}
\textcolor{financeblue}{\textbf{Fonction}} &
\textcolor{financeblue}{\textbf{Portée}} &
\textcolor{financeblue}{\textbf{Tests concernés}} \\
\code{usp\_credibilite()} & LN + MW & Calibration des deux méthodes~: seule
fonction de calcul commune aux deux branches \\
\code{test\_grubbs()} & LN + MW & Grubbs sur $z_t$, sur ratios bruts, et sur
résidus de Mack \\
\code{test\_rosner()} & LN + MW & ESD généralisé, dans les deux méthodes \\
\code{test\_runs()} & LN + MW & Test des suites, dans les deux méthodes \\
\code{stat\_dw()} & LN + MW & Durbin--Watson, dans les deux méthodes \\
\code{.shapiro\_sur()} & LN + MW & Shapiro--Wilk (Royston et loi nulle simulée),
dans les deux méthodes \\
\code{stat\_lilliefors()}, \code{lillie\_p()} & LN + MW & Lilliefors, dans les
deux méthodes \\
\code{test\_lm\_complet()} & LN & Student sur la constante, Student sur la
pente, Fisher global, $R^2$~: \textbf{quatre entrées issues d'un seul appel} \\
\code{usp\_noyau()}, \code{usp\_pi()} & LN & Tout le noyau lognormal~:
résidus $z_t$, donc l'intégralité de H2, H3, H4 \\
\code{.stats\_bootstrapables()} & LN & Les 34 statistiques bootstrapées de la
méthode lognormale, en un seul appel \\
\code{mw\_ajuster()} & MW & $\hat f_j$, $\hat\sigma_j^2$, $\hat C(i,j)$~: entrée
de la MSEP, des résidus, de tous les tests M1 à M6 \\
\code{mw\_residus()} & MW & Support de M1, M2, M3, M4 et M5 \\
\code{.mw\_stats()} & MW & Les 7 statistiques bootstrapées de la méthode \\
\end{tabular}
\end{center}

\newpage
\subsection{Graphe chapeau : les deux branches et leurs points de partage}

\begin{figure}[H]
\centering
\begin{tikzpicture}[
  font=\footnotesize,
  ent/.style  = {rectangle, rounded corners=2pt, draw=financeblue, fill=tableheader,
                 text=financeblue, align=center, inner sep=3pt, text width=34mm,
                 font=\footnotesize\bfseries},
  ded/.style  = {rectangle, rounded corners=2pt, draw=financeblue, fill=white,
                 align=center, inner sep=3pt, text width=34mm},
  mut/.style  = {rectangle, rounded corners=2pt, draw=orange!80!black, line width=1.1pt,
                 fill=orange!12, align=center, inner sep=3pt, text width=34mm},
  sor/.style  = {rectangle, rounded corners=2pt, draw=greenline, line width=0.9pt,
                 fill=greenline!8, align=center, inner sep=3pt, text width=36mm},
  fl/.style   = {-{Latex[length=1.8mm]}, draw=financeblue, line width=0.55pt},
  flo/.style  = {-{Latex[length=1.8mm]}, draw=orange!80!black, line width=0.7pt},
  flx/.style  = {font=\scriptsize\itshape, text=greenline, inner sep=1pt}
]
\node[ent] (xy) {\code{xt}, \code{yt}};
\node[ent, right=42mm of xy] (tri) {\code{triangle}};

\node[ded, below=6mm of xy] (fit) {\code{usp\_ajuster()}};
\node[ded, below=6mm of tri] (amw) {\code{mw\_ajuster()}};

\node[ded, below=6mm of fit] (zt) {résidus $z_t$\\ \scriptsize via \code{usp\_noyau()}};
\node[ded, below=6mm of amw] (rij) {résidus $r(i,j)$\\ \scriptsize via \code{mw\_residus()}};

\node[ded, below=6mm of zt] (bln) {\code{usp\_bootstrap()}};
\node[ded, below=6mm of rij] (bmw) {\code{mw\_bootstrap()}};

\node[mut] at ($(bln)!0.5!(bmw) + (0,-24mm)$) (part)
  {fonctions de statistique mutualisées\\[1pt]
   \scriptsize \code{stat\_dw()}, \code{test\_runs()}, \code{test\_grubbs()},\\
   \scriptsize \code{test\_rosner()}, \code{.shapiro\_sur()},\\
   \scriptsize \code{stat\_lilliefors()}, \code{lillie\_p()}};

\node[ded] at (bln |- part.south) [yshift=-16mm] (tln) {\code{usp\_tests()}\\ \scriptsize H1--H4};
\node[ded] at (bmw |- part.south) [yshift=-16mm] (tmw) {\code{mw\_tests()}\\ \scriptsize M1--M6};

% usp_credibilite() est placee SOUS la rangee des deux fonctions de tests,
% faute de quoi les fleches de convergence la longent au lieu d'y entrer.
\node[mut] at ($(tln)!0.5!(tmw) + (0,-18mm)$) (cred)
  {\code{usp\_credibilite()}\\ \scriptsize seule fonction de calibration commune};
\node[sor, below=6mm of cred] (out) {$\sigma_{\mathrm{USP}}$ et table des tests};

\draw[fl] (xy) -- (fit);   \draw[fl] (tri) -- (amw);
\draw[fl] (fit) -- node[flx,right] {\code{fit}} (zt);
\draw[fl] (amw) -- node[flx,right] {$\hat f_j$, $\hat\sigma_j$} (rij);
\draw[fl] (zt) -- (bln);   \draw[fl] (rij) -- (bmw);
% Toutes les liaisons de convergence et de divergence sont COUDEES : des
% segments diagonaux partant d'un meme point se chevauchent et deviennent
% illisibles des que les blocs sont proches.
% Entrees et sorties decalees verticalement : sans ce decalage, le segment
% descendant et le segment remontant partageraient la meme abscisse et se
% superposeraient.
\draw[flo] (bln.south) |- ([yshift=6mm]part.west);
\draw[flo] (bmw.south) |- ([yshift=6mm]part.east);
\draw[fl] ([yshift=-6mm]part.west) -| (tln.north);
\draw[fl] ([yshift=-6mm]part.east) -| (tmw.north);
\draw[fl] (tln.south) |- (cred.west);
\draw[fl] (tmw.south) |- (cred.east);
\draw[fl] (cred) -- (out);
\end{tikzpicture}
\caption{Vue chapeau. Les deux branches sont disjointes jusqu'aux fonctions de
statistique, dont sept sont mutualisées (en orangé), puis se rejoignent sur
\code{usp\_credibilite()}. Une anomalie dans un bloc orangé affecterait
\emph{les deux} méthodes.}
\label{fig:annexe-chapeau}
\end{figure}

\newpage
\subsection{Méthode lognormale : hypothèse H1}

\begin{figure}[H]
\centering
\begin{tikzpicture}[
  font=\scriptsize,
  ent/.style = {rectangle, rounded corners=2pt, draw=financeblue, fill=tableheader,
                text=financeblue, align=center, inner sep=2.5pt, text width=28mm,
                font=\scriptsize\bfseries},
  ded/.style = {rectangle, rounded corners=2pt, draw=financeblue, fill=white,
                align=center, inner sep=2.5pt, text width=42mm},
  mut/.style = {rectangle, rounded corners=2pt, draw=orange!80!black, line width=1.1pt,
                fill=orange!12, align=center, inner sep=2.5pt, text width=42mm},
  sor/.style = {rectangle, rounded corners=2pt, draw=greenline, fill=greenline!8,
                align=center, inner sep=2.5pt, text width=30mm},
  fl/.style  = {-{Latex[length=1.6mm]}, draw=financeblue, line width=0.5pt},
  li/.style  = {draw=financeblue, line width=0.5pt}
]
\node[ent] (r) {ratios $r_t = y_t/x_t$};
\node[mut, right=20mm of r] (lm) {\code{test\_lm\_complet()}};
\node[sor, right=20mm of lm, yshift=13.5mm]  (o1) {Student constante};
\node[sor, right=20mm of lm, yshift=4.5mm]   (o2) {Student pente};
\node[sor, right=20mm of lm, yshift=-4.5mm]  (o3) {Fisher global};
\node[sor, right=20mm of lm, yshift=-13.5mm] (o4) {$R^2$ (indicateur)};

\node[ded, below=17mm of lm] (tost) {\code{test\_tost\_intercept()}};
\node[sor, right=20mm of tost] (o5) {Équivalence (TOST)};
\node[ded, below=6mm of tost] (reset) {\code{test\_reset()}};
\node[sor, right=20mm of reset] (o6) {RESET};
\node[ded, below=6mm of reset] (mk) {\code{test\_mann\_kendall()}
  + \code{mk\_p\_exacte()}\\ \scriptsize via \code{.mk\_loi\_exacte()}, \code{.poly\_mult()}};
\node[sor, right=20mm of mk] (o7) {Mann--Kendall};
\node[ded, below=7mm of mk] (cs) {\code{test\_cox\_stuart()}};
\node[sor, right=20mm of cs] (o8) {Cox--Stuart};
\node[ded, below=6mm of cs] (sp) {\code{stats::cor.test()}};
\node[sor, right=20mm of sp] (o9) {Spearman ($\times2$)};

\draw[fl] (r) -- (lm);
% Les quatre sorties issues du meme appel sont desservies par un tronc commun
% puis des coudes, plutot que par quatre diagonales partant d'un meme point.
\coordinate (bus) at ([xshift=8mm]lm.east);
\draw[li] (lm.east) -- (bus);
\foreach \o in {o1,o2,o3,o4} \draw[fl] (bus) |- (\o.west);
\foreach \a/\b in {tost/o5, reset/o6, mk/o7, cs/o8, sp/o9} \draw[fl] (\a) -- (\b);
\foreach \b in {tost,reset,mk,cs,sp} \draw[fl] (r) |- (\b);
\end{tikzpicture}
\caption{Hypothèse H1. \code{test\_lm\_complet()} alimente \textbf{quatre}
entrées de la table des tests à partir d'un unique ajustement~: une anomalie s'y
propagerait aux quatre. Mann--Kendall illustre la configuration inverse~: un
seul test, mais quatre fonctions dont deux internes.}
\label{fig:annexe-h1}
\end{figure}

\subsection{Méthode lognormale : hypothèse H2}

\begin{figure}[H]
\centering
\begin{tikzpicture}[
  font=\scriptsize,
  ent/.style = {rectangle, rounded corners=2pt, draw=financeblue, fill=tableheader,
                text=financeblue, align=center, inner sep=2.5pt, text width=28mm,
                font=\scriptsize\bfseries},
  ded/.style = {rectangle, rounded corners=2pt, draw=financeblue, fill=white,
                align=center, inner sep=2.5pt, text width=42mm},
  mut/.style = {rectangle, rounded corners=2pt, draw=orange!80!black, line width=1.1pt,
                fill=orange!12, align=center, inner sep=2.5pt, text width=42mm},
  sor/.style = {rectangle, rounded corners=2pt, draw=greenline, fill=greenline!8,
                align=center, inner sep=2.5pt, text width=30mm},
  fl/.style  = {-{Latex[length=1.6mm]}, draw=financeblue, line width=0.5pt}
]
\node[ent] (z) {résidus $z_t$};
\node[mut, right=20mm of z] (bp) {\code{test\_breusch\_pagan()}};
\node[sor, right=20mm of bp] (p1) {BP studentisé (Koenker)};
\node[ded, below=6mm of bp] (bp79) {\code{test\_breusch\_pagan\_}\\\code{original()}};
\node[sor, right=20mm of bp79] (p2) {BP original 1979};
\node[ded, below=6mm of bp79] (wh) {\code{test\_white()}};
\node[sor, right=20mm of wh] (p3) {White};
\node[ded, below=6mm of wh] (gq) {\code{test\_goldfeld\_quandt()}};
\node[sor, right=20mm of gq] (p4) {Goldfeld--Quandt};
\node[ded, below=6mm of gq] (bf) {\code{test\_brown\_forsythe()}};
\node[sor, right=20mm of bf] (p5) {Brown--Forsythe};
\node[ded, below=6mm of bf] (ks2) {\code{stats::ks.test()}};
\node[sor, right=20mm of ks2] (p6) {Smirnov 2 échantillons};

\foreach \a/\b in {bp/p1, bp79/p2, wh/p3, gq/p4, bf/p5, ks2/p6} \draw[fl] (\a) -- (\b);
\draw[fl] (z) -- (bp);
\foreach \b in {bp79,wh,gq,bf,ks2} \draw[fl] (z) |- (\b);
\end{tikzpicture}
\caption{Hypothèse H2. Chacun des six tests dispose de sa propre fonction~: la
portée d'une anomalie y est circonscrite à un test, contrairement à H1.}
\label{fig:annexe-h2}
\end{figure}

\subsection{Méthode lognormale : hypothèses H3 et H4, stabilité}

\begin{landscape}
\begin{figure}[H]
\centering
\begin{tikzpicture}[
  font=\tiny,
  ent/.style = {rectangle, rounded corners=2pt, draw=financeblue, fill=tableheader,
                text=financeblue, align=center, inner sep=2.5pt, text width=25mm,
                font=\scriptsize\bfseries},
  ded/.style = {rectangle, rounded corners=2pt, draw=financeblue, fill=white,
                align=center, inner sep=2.5pt, text width=38mm},
  mut/.style = {rectangle, rounded corners=2pt, draw=orange!80!black, line width=1.1pt,
                fill=orange!12, align=center, inner sep=2.5pt, text width=38mm},
  int/.style = {rectangle, rounded corners=2pt, draw=gray!70, fill=gray!12,
                align=center, inner sep=2.5pt, text width=38mm},
  sor/.style = {rectangle, rounded corners=2pt, draw=greenline, fill=greenline!8,
                align=center, inner sep=2.5pt, text width=32mm},
  fl/.style  = {-{Latex[length=1.6mm]}, draw=financeblue, line width=0.5pt},
  fp/.style  = {-{Latex[length=1.6mm]}, draw=gray!70, line width=0.5pt, dashed}
]
\node[ent] (z) {résidus $z_t$};

% --- Colonne des fonctions ---------------------------------------------------
\node[mut, right=13mm of z, yshift=30mm] (sw) {\code{.shapiro\_sur()}};
\node[int, below=4mm of sw] (swn) {\code{sw\_loi\_nulle()}\\ \scriptsize mise en cache};
\node[ded, below=5mm of swn] (sf) {\code{test\_shapiro\_}\\ \code{francia()}};
\node[ded, below=4mm of sf] (ad) {\code{stat\_ad()}\\ + \code{ad\_p\_stephens()}};
\node[ded, below=4mm of ad] (cv) {\code{stat\_cvm()}\\ + \code{cvm\_p\_stephens()}};
\node[mut, below=4mm of cv] (li) {\code{stat\_lilliefors()}\\ + \code{lillie\_p()}};
\node[ded, below=4mm of li] (ks) {\code{stat\_ks()}};
\node[ded, below=4mm of ks] (jb) {\code{test\_jarque\_bera()}};
\node[ded, below=4mm of jb] (da) {\code{test\_dagostino\_}\\ \code{skew()}};
\node[ded, below=4mm of da] (ag) {\code{test\_anscombe\_kurt()}};
\node[mut, below=6mm of ag] (dw) {\code{stat\_dw()}\\ + \code{dw\_p\_exacte()}};
\node[ded, below=5mm of dw] (lb) {\code{stats::Box.test()}};
\node[mut, below=4mm of lb] (ru) {\code{test\_runs()}\\ + \code{runs\_p\_exacte()}};
\node[ded, below=6mm of ru] (sf2) {\code{stat\_supF()},\\ \code{stat\_cusum()}};
\node[mut, below=4mm of sf2] (gr) {\code{test\_grubbs()},\\ \code{test\_rosner()}};

% --- Colonne des fonctions internes (lois de reference) ----------------------
\node[int, right=14mm of dw] (imh) {\code{.imhof\_p\_sup0()}};
\node[int, right=14mm of ru] (rd)  {\code{.runs\_dens()}};

% --- Colonne des sorties : abscisse COMMUNE ---------------------------------
% Auparavant chaque sortie etait positionnee relativement a son bloc de gauche,
% de largeurs inegales (fonction 44 mm, interne 34 mm), d'ou un desalignement de
% la derniere colonne. Toutes partagent desormais une abscisse de reference.
\coordinate (cx) at ([xshift=10mm]imh.east);
\node[sor, anchor=west] at (cx |- sw)  (s1)  {Shapiro--Wilk (Royston)};
\node[sor, anchor=west] at (cx |- swn) (s2)  {SW loi nulle simulée};
\node[sor, anchor=west] at (cx |- sf)  (s3)  {Shapiro--Francia};
\node[sor, anchor=west] at (cx |- ad)  (s4)  {Anderson--Darling};
\node[sor, anchor=west] at (cx |- cv)  (s5)  {Cramér--von Mises};
\node[sor, anchor=west] at (cx |- li)  (s6)  {Lilliefors};
\node[sor, anchor=west] at (cx |- ks)  (s7)  {KS contre $\mathcal{N}(0,1)$};
\node[sor, anchor=west] at (cx |- jb)  (s8)  {Jarque--Bera};
\node[sor, anchor=west] at (cx |- da)  (s9)  {D'Agostino};
\node[sor, anchor=west] at (cx |- ag)  (s10) {Anscombe--Glynn};
\node[sor, anchor=west] at (cx |- dw)  (s11) {Durbin--Watson};
\node[sor, anchor=west] at (cx |- lb)  (s12) {Ljung--Box, Box--Pierce};
\node[sor, anchor=west] at (cx |- ru)  (s13) {Test des suites};
\node[sor, anchor=west] at (cx |- sf2) (s14) {sup-F, CUSUM};
\node[sor, anchor=west] at (cx |- gr)  (s15) {Grubbs, Rosner};

% --- Liaisons ----------------------------------------------------------------
\draw[fl] (sw) -- (s1);
\draw[fp] (sw) -- (swn); \draw[fl] (swn) -- (s2);
\foreach \a/\b in {sf/s3, ad/s4, cv/s5, li/s6, ks/s7, jb/s8, da/s9, ag/s10,
                   lb/s12, sf2/s14, gr/s15} \draw[fl] (\a) -- (\b);
\draw[fp] (dw) -- (imh); \draw[fl] (imh) -- (s11);
\draw[fp] (ru) -- (rd);  \draw[fl] (rd) -- (s13);
\foreach \b in {sw,sf,ad,cv,li,ks,jb,da,ag,dw,lb,ru,sf2,gr}
  \draw[fl] (z) -| ([xshift=-4mm]\b.west) -- (\b);
\end{tikzpicture}
\caption{Hypothèses H3 et H4 et diagnostics de stabilité. Tous les tests
partent du même objet, les résidus $z_t$~: une erreur dans \code{usp\_noyau()}
affecterait l'intégralité de ce graphe. Trois fonctions internes (en gris)
fournissent des lois de référence~: \code{sw\_loi\_nulle()},
\code{.imhof\_p\_sup0()} et \code{.runs\_dens()}. \code{.shapiro\_sur()} alimente
\emph{deux} entrées distinctes de la table des tests.}
\label{fig:annexe-h3h4}
\end{figure}
\end{landscape}

\newpage
\subsection{Méthode Merz--Wüthrich : de l'ajustement aux tests M1 à M6}

\begin{landscape}
\begin{figure}[H]
\centering
\begin{tikzpicture}[
  font=\tiny,
  ent/.style = {rectangle, rounded corners=2pt, draw=financeblue, fill=tableheader,
                text=financeblue, align=center, inner sep=2.5pt, text width=30mm,
                font=\tiny\bfseries},
  ded/.style = {rectangle, rounded corners=2pt, draw=financeblue, fill=white,
                align=center, inner sep=3pt, text width=44mm},
  mut/.style = {rectangle, rounded corners=2pt, draw=orange!80!black, line width=1.1pt,
                fill=orange!12, align=center, inner sep=3pt, text width=44mm},
  int/.style = {rectangle, rounded corners=2pt, draw=gray!70, fill=gray!12,
                align=center, inner sep=3pt, text width=32mm},
  sor/.style = {rectangle, rounded corners=2pt, draw=greenline, fill=greenline!8,
                align=center, inner sep=3pt, text width=26mm},
  fl/.style  = {-{Latex[length=1.6mm]}, draw=financeblue, line width=0.5pt},
  fp/.style  = {-{Latex[length=1.6mm]}, draw=gray!70, line width=0.5pt, dashed},
  flx/.style = {font=\tiny\itshape, text=greenline, inner sep=1pt},
  lig/.style = {draw=gray!70, line width=0.5pt, dashed}
]
\node[ent] (tri) {triangle $C(i,j)$};
\node[ded, below=5mm of tri] (val) {\code{mw\_valider\_triangle()}};
\node[mut, below=5mm of val] (aj) {\code{mw\_ajuster()}\\
  \scriptsize $\hat f_j$, $\hat\sigma^2_j$, $\hat C(i,j)$, $\hat Q_j$, $S_j$, $S'_j$};

\node[ded, left=14mm of aj, yshift=-16mm] (msep) {\code{mw\_msep()}};
\node[ded, below=5mm of msep] (par) {\code{mw\_parametre()}};
\node[sor, below=5mm of par] (sig) {$\sigma_{\mathrm{USP}}$};

\node[mut, below=16mm of aj] (res) {\code{mw\_residus()}\\ \scriptsize résidus $r(i,j)$};

% --- Colonne des fonctions de test ------------------------------------------
\node[ded, right=16mm of res, yshift=34mm] (m1)
  {\code{mw\_test\_ordonnee\_}\\ \code{origine()},
   \code{mw\_test\_}\\ \code{homogeneite\_f()},
   \code{mw\_test\_courbure()},\\ \code{mw\_famille\_alpha()}};
\node[ded, below=5mm of m1] (m2)
  {\code{stats::lm(r\textasciicircum2 \textasciitilde\ C)},\\
   \code{mw\_test\_}\\ \code{exposant\_variance()}};
\node[ded, below=5mm of m2] (cal) {\code{mw\_test\_annees\_}\\ \code{calendaires()}};
\node[ded, below=5mm of cal] (cor) {\code{mw\_stat\_}\\ \code{correlation\_dev()},
   \code{mw\_test\_}\\ \code{homogeneite\_accident()}};
\node[mut, below=5mm of cor] (dw) {\code{stat\_dw()}, \code{test\_runs()}};
\node[mut, below=5mm of dw] (gr) {\code{test\_grubbs()}, \code{test\_rosner()}};
\node[mut, below=5mm of gr] (sw) {\code{.shapiro\_sur()}, \code{sw\_p\_loi\_nulle()},\\
   \code{stat\_lilliefors()}, \code{lillie\_p()}};

% --- Fonctions internes : colonne intermediaire dediee -----------------------
% .fisher_combine() dessert quatre tests de M1 et un de M2 ; il est place en
% vis-a-vis de ces deux blocs plutot qu'a gauche, ou il chevauchait mw_residus().
\node[int] at ([xshift=48mm]$(m1)!0.5!(m2)$) (fi)
  {\code{.fisher\_combine()}\\ \scriptsize dessert 4 tests M1 + 1 test M2};
\node[int] at ([xshift=48mm]cal) (mo) {\code{.mack\_moments\_Z()}};

% --- Colonne des sorties : abscisse commune ---------------------------------
\coordinate (cx) at ([xshift=16mm]fi.east);
\node[sor, anchor=west] at (cx |- m1)  (q1) {M1 (5 entrées)};
\node[sor, anchor=west] at (cx |- m2)  (q2) {M2 (3 entrées)};
\node[sor, anchor=west] at (cx |- cal) (q3) {M3 année calendaire};
\node[sor, anchor=west] at (cx |- cor) (q4) {M3 (3 entrées)};
\node[sor, anchor=west] at (cx |- dw)  (q5) {M3 DW, suites};
\node[sor, anchor=west] at (cx |- gr)  (q6) {M4 aberrants};
\node[sor, anchor=west] at (cx |- sw)  (q7) {M5 normalité};

\node[ded, below=7mm of res] (boot) {\code{mw\_bootstrap()}};
\node[int, below=5mm of boot] (st)
  {\code{mw\_simuler\_triangle()},\\ \code{.mw\_stats()}};
\node[sor, below=5mm of st] (pmc) {p-values de Monte-Carlo};

% --- Liaisons ----------------------------------------------------------------
\draw[fl] (tri) -- (val); \draw[fl] (val) -- (aj);
\draw[fl] (aj) -| node[flx,above left] {$\hat f_j$, $\hat Q_j$} (msep);
\draw[fl] (msep) -- (par); \draw[fl] (par) -- (sig);
\draw[fl] (aj) -- node[flx,right] {$\hat f_j$, $\hat\sigma_j$} (res);
\foreach \b in {m1,m2,cal,cor,dw,gr,sw}
  \draw[fl] (res) -| ([xshift=-5mm]\b.west) -- (\b);
\foreach \a/\b in {cor/q4, dw/q5, gr/q6, sw/q7} \draw[fl] (\a) -- (\b);
% Convergence puis divergence autour de .fisher_combine() : troncs communs et
% coudes, pour eviter quatre diagonales issues d'un meme point.
\coordinate (bin)  at ([xshift=-6mm]fi.west);
\coordinate (bout) at ([xshift=6mm]fi.east);
\draw[lig] (m1.east) -| (bin); \draw[lig] (m2.east) -| (bin);
\draw[fp] (bin) -- (fi.west);
\draw[lig] (fi.east) -- (bout);
\draw[fl] (bout) |- (q1.west); \draw[fl] (bout) |- (q2.west);
\draw[fp] (cal) -- (mo); \draw[fl] (mo) -- (q3);
\draw[fl] (res) -- (boot);
\draw[fp] (boot) -- (st); \draw[fl] (st) -- (pmc);
\end{tikzpicture}
\caption{Méthode Merz--Wüthrich. \code{mw\_ajuster()} et \code{mw\_residus()}
sont les deux points de passage obligés~: la première alimente à la fois la MSEP
et les résidus, la seconde l'intégralité des tests M1 à M5. Quatre blocs de
tests reposent sur des fonctions déjà employées par la méthode lognormale
(en orangé), ce qui mutualise la revue de code mais étend la portée d'une
anomalie éventuelle.}
\label{fig:annexe-mw}
\end{figure}
\end{landscape}

\subsection{Lecture en termes de revue de code}

\begin{center}
\small
\rowcolors{2}{tablerow}{white}
\begin{tabular}{@{}p{4.4cm}p{3.0cm}p{7.6cm}@{}}
\rowcolor{tableheader}
\textcolor{financeblue}{\textbf{Configuration}} &
\textcolor{financeblue}{\textbf{Exemples}} &
\textcolor{financeblue}{\textbf{Conséquence pour la revue}} \\
Un test, une fonction dédiée & White, Goldfeld--Quandt, Brown--Forsythe,
Jarque--Bera, corrélation entre colonnes & Revue circonscrite~: une anomalie
n'affecte que ce test. \\
Un test, plusieurs fonctions & Mann--Kendall (\code{test\_mann\_kendall()} +
\code{mk\_p\_exacte()} + deux internes), Durbin--Watson (\code{stat\_dw()} +
\code{dw\_p\_exacte()} + \code{.imhof\_p\_sup0()}) & La statistique et la loi de
référence sont séparées~: il faut vérifier les deux, et leur cohérence
d'appariement. \\
Une fonction, plusieurs tests & \code{test\_lm\_complet()} (4 entrées),
\code{.shapiro\_sur()} (2 entrées) & Une anomalie se propage à toutes les
entrées issues de l'appel. \\
Fonction commune aux deux méthodes & \code{test\_grubbs()},
\code{test\_rosner()}, \code{test\_runs()}, \code{stat\_dw()},
\code{.shapiro\_sur()}, \code{stat\_lilliefors()}, \code{usp\_credibilite()} &
Portée maximale~: une anomalie affecterait les deux dispositifs, donc les deux
dossiers. Priorité de revue. \\
\end{tabular}
\end{center}

\noindent
Cette hiérarchie explique l'ordre de vérification retenu lors du développement~:
les fonctions communes et les noyaux (\code{usp\_noyau()}, \code{mw\_ajuster()})
ont fait l'objet de contrôles indépendants --- recalcul manuel, identités
analytiques, comparaison à des valeurs publiées --- avant les fonctions dédiées,
dont la portée d'erreur est limitée à un seul test.


\end{document}
